%%%%%%%%%%%%%%%%%%%%%%%%%%%%%%%%%%%%%%%%%%%%%%%%%%%%%%%%%%%%%%%%%%%%%%%%%%%%%%%%
%% APPENDIX H: SUPPORTING MATERIAL FOR CHAPTER 3 (METHODOLOGY AND DATA)
%% Merged from: Algorithms, Datasets, Preprocessing Pipelines, Model Architectures
%%%%%%%%%%%%%%%%%%%%%%%%%%%%%%%%%%%%%%%%%%%%%%%%%%%%%%%%%%%%%%%%%%%%%%%%%%%%%%%%

%% Supplementary Material (merged into Chapter 3)
% [SPACE-FIX] clearpage removed to eliminate empty page
\phantomsection\label{app:ch3_support_appx}

\section{Supplementary Material: Algorithms, Datasets, and Architectures}

This section provides detailed supplementary material for the methodology and data collection. It includes the mathematical foundations of all algorithms used, complete dataset specifications, preprocessing pipeline implementations, and layer-by-layer model architecture specifications. These details support reproducibility and enable independent verification of the experimental setup.

%% ====================================================================
%% ALGORITHMS AND MATHEMATICAL FOUNDATIONS (from old Appendix)
%% ====================================================================

\label{app:algorithms}

\section{EEG Signal Processing Algorithms}

Three signal processing algorithms form the mathematical foundation of the EEG preprocessing pipeline. Table~\ref{app_tab:signal_proc_algorithms} consolidates their transfer functions, key parameters, and roles in the RGAIG framework.
\needspace{8\baselineskip}
{\tablestripes\tablefontsize
\begin{xltabular}{\textwidth}{@{} p{2.5cm} L p{3.5cm} @{}}
\caption{EEG Signal Processing Algorithms: Mathematical Foundations}\label{app_tab:signal_proc_algorithms} \\
\toprule
\thead{Algorithm} & \thead{Transfer Function / Definition} & \thead{Parameters and Role} \\
\midrule
\endfirsthead
\multicolumn{3}{@{}l}{\textit{(\,Tab.~\ref{app_tab:signal_proc_algorithms} continued from previous page\,)}} \\
\toprule
\thead{Algorithm} & \thead{Transfer Function / Definition} & \thead{Parameters and Role} \\
\midrule
\endhead
\bottomrule
\endlastfoot
Band-pass filter (Butterworth) & $H(s) = \dfrac{1}{\sqrt{1 + \left(\frac{s}{\omega_c}\right)^{2n}}}$ & $\omega_c$ = cutoff frequency; $n$ = 4 (filter order); retains 0.5--45\,Hz \\
\addlinespace
Continuous Wavelet Transform & $W_f(a,b) = \dfrac{1}{\sqrt{|a|}} \displaystyle\int_{-\infty}^{\infty} f(t) \psi^*\!\left(\frac{t-b}{a}\right) dt$ & $\psi(t)$ = Morlet wavelet; $a$ = scale; $b$ = translation; 64 log-spaced bins \\
\addlinespace
Power Spectral Density (Welch) & $P_{xx}(f) = \dfrac{1}{KMU} \displaystyle\sum_{k=0}^{K-1} \left| \sum_{n=0}^{M-1} x_k(n) w(n) e^{-j2\pi fn} \right|^2$ & $K$ = segments; $M$ = segment length; $U$ = normalisation; $w(n)$ = Hann window \\
\end{xltabular}}
\vspace{0.5em}\noindent\textit{Note.} The Butterworth band-pass filter removes DC drift ($<$0.5\,Hz) and high-frequency noise ($>$45\,Hz). CWT provides multi-resolution time-frequency analysis for non-stationary EEG signals. PSD (Welch method) estimates frequency-domain power for band-power feature extraction.

\section{Classical and Deep Learning Algorithm Foundations}

Table~\ref{app_tab:algorithm_foundations} consolidates the mathematical foundations of the five core algorithms used across the RGAIG classification pipeline. Each formula is referenced at its point of use in the methodology chapter; full derivations are provided below the table.
\needspace{8\baselineskip}
\noindent Table~\ref{app_tab:algorithm_foundations} algorithm mathematical foundations: classical and deep learning.

{\tablestripes\tablefontsize
\begin{xltabular}{\textwidth}{@{} p{1.4cm} p{1.4cm} L p{1.8cm} @{}}
\caption{Algorithm Mathematical Foundations: Classical and Deep Learning}\label{app_tab:algorithm_foundations} \\
\toprule
\thead{Algorithm} & \thead{Paradigm} & \thead{Core Formula} & \thead{RGAIG Usage} \\
\midrule
\endfirsthead
\multicolumn{4}{@{}l}{\textit{(\,Tab.~\ref{app_tab:algorithm_foundations} continued from previous page\,)}} \\
\toprule
\thead{Algorithm} & \thead{Paradigm} & \thead{Core Formula} & \thead{RGAIG Usage} \\
\midrule
\endhead
\bottomrule
\endlastfoot
SVM & Margin-based & $\min_{w,b,\xi} \frac{1}{2}\|w\|^2 + C\sum_{i}\xi_i$\newline s.t.\ $y_i(w^T\phi(x_i) + b) \geq 1 - \xi_i$ & Baseline classifier (all EEG domains) \\
\addlinespace
Random Forest & Ensemble bagging & $\hat{y} = \frac{1}{B} \sum_{b=1}^{B} T_b(x)$\newline ($B$ trees, $T_b$ = tree $b$ prediction) & VotingClassifier component (deployed) \\
\addlinespace
1D-CNN & Convolutional & $y[n] = \sum_{k=0}^{K-1} x[n-k] \cdot w[k] + b$ & Temporal pattern extraction from raw EEG \\
\addlinespace
LSTM & Recurrent & $f_t = \sigma(W_f [h_{t-1}, x_t] + b_f)$;\newline $C_t = f_t C_{t-1} + i_t \tilde{C}_t$;\newline $h_t = o_t \tanh(C_t)$ & Temporal dynamics in CNN-LSTM hybrid \\
\addlinespace
Self-Attention & Transformer & $\text{Attention}(Q,K,V) = \text{softmax}\!\left(\frac{QK^T}{\sqrt{d_k}}\right)V$ & Transformer encoder for EEG sequences \\
\end{xltabular}}
\vspace{0.5em}\noindent\textit{Note.} SVM = support vector machine; CNN = convolutional neural network; LSTM = long short-term memory; $\sigma$ = sigmoid function; $f_t$, $i_t$, $o_t$ = forget, input, output gates; $Q$, $K$, $V$ = query, key, value matrices; $d_k$ = key dimensionality. Full LSTM gate equations: $i_t = \sigma(W_i [h_{t-1}, x_t] + b_i)$, $\tilde{C}_t = \tanh(W_C [h_{t-1}, x_t] + b_C)$, $o_t = \sigma(W_o [h_{t-1}, x_t] + b_o)$.

\section{Evaluation Metrics}

Five evaluation metrics are used throughout the experimental evaluation. Table~\ref{app_tab:eval_metrics_formulae} consolidates their mathematical definitions, interpretation guidelines, and the acceptance thresholds applied in this study.
\needspace{8\baselineskip}
{\tablestripes\tablefontsize
\begin{xltabular}{\textwidth}{@{} p{1.8cm} L p{2.5cm} p{1.8cm} @{}}
\caption{Evaluation Metrics: Mathematical Definitions and Acceptance Thresholds}\label{app_tab:eval_metrics_formulae} \\
\toprule
\thead{Metric} & \thead{Formula} & \thead{Interpretation} & \thead{Threshold} \\
\midrule
\endfirsthead
\multicolumn{4}{@{}l}{\textit{(\,Tab.~\ref{app_tab:eval_metrics_formulae} continued from previous page\,)}} \\
\toprule
\thead{Metric} & \thead{Formula} & \thead{Interpretation} & \thead{Threshold} \\
\midrule
\endhead
\bottomrule
\endlastfoot
Accuracy & $\dfrac{TP + TN}{TP + TN + FP + FN}$ & Overall correct classification rate & $\geq$85\% \\
\addlinespace
Precision & $\dfrac{TP}{TP + FP}$ & Proportion of positive predictions that are correct & $\geq$80\% \\
\addlinespace
Recall & $\dfrac{TP}{TP + FN}$ & Proportion of actual positives correctly identified & $\geq$80\% \\
\addlinespace
F1-Score & $2 \times \dfrac{\text{Precision} \times \text{Recall}}{\text{Precision} + \text{Recall}}$ & Harmonic mean of precision and recall & $\geq$0.80 \\
\addlinespace
AUC-ROC & $\displaystyle\int_0^1 TPR(FPR^{-1}(x))\, dx$ & Discrimination across all classification thresholds & $\geq$0.90 \\
\end{xltabular}}
\vspace{0.5em}\noindent\textit{Note.} TP = true positive; TN = true negative; FP = false positive; FN = false negative; TPR = true positive rate; FPR = false positive rate; AUC-ROC = area under the receiver operating characteristic curve. Thresholds are predefined acceptance criteria evaluated in Chapter~\ref{chap:analysis}.


%% ====================================================================
%% DATASET SPECIFICATIONS (from old Appendix)
%% ====================================================================

\label{app:datasets}

\section{KAU ASD-EEG Dataset (ASD)}

The King Abdulaziz University ASD-EEG Dataset \citep{djemal2017asd} provides resting-state EEG recordings for autism spectrum disorder classification. Table~\ref{tab:kau_asd_eeg} summarises the dataset parameters.
\needspace{8\baselineskip}
{\tablestripes\tablefontsize
\begin{xltabular}{\textwidth}{@{} p{4cm} L @{}}
\caption{KAU ASD-EEG Dataset Specifications}\label{tab:kau_asd_eeg} \\
\toprule
\thead{Parameter} & \thead{Value} \\
\midrule
\endfirsthead
\multicolumn{2}{@{}l}{\textit{(\,Tab.~\ref{tab:kau_asd_eeg} continued from previous page\,)}} \\
\toprule
\thead{Parameter} & \thead{Value} \\
\midrule
\endhead
\bottomrule
\endlastfoot
Total Subjects & 19 \\
ASD Subjects & 9 \\
Control Subjects & 10 \\
EEG Channels & 19 (10-20 international system) \\
Sampling Rate & 256 Hz \\
Recording Format & EDF (European Data Format) \\
Paradigm & Resting-state EEG \\
Site & KAU Hospital, Jeddah \\
Modality & Resting-state EEG (19-channel, 10-20 system) \\
\end{xltabular}}

\section{PhysioNet Parkinson's Dataset}

\noindent Table~\ref{tab:physionet} documents the PhysioNet Parkinson's disease dataset specifications.
\needspace{8\baselineskip}
{\tablestripes\tablefontsize
\begin{xltabular}{\textwidth}{@{} p{4cm} L @{}}
\caption{PhysioNet PD Dataset Specifications}\label{tab:physionet} \\
\toprule
\thead{Parameter} & \thead{Value} \\
\midrule
\endfirsthead
\multicolumn{2}{@{}l}{\textit{(\,Tab.~\ref{tab:physionet} continued from previous page\,)}} \\
\toprule
\thead{Parameter} & \thead{Value} \\
\midrule
\endhead
\bottomrule
\endlastfoot
Total Subjects & 52 \\
PD Patients & 26 \\
Healthy Controls & 26 \\
EEG Channels & 32 \\
Sampling Rate & 256 Hz \\
Recording Duration & 5 minutes \\
\end{xltabular}}

\section{DEAP Dataset (Stress/Emotion)}

\noindent Table~\ref{tab:deap} presents the DEAP dataset specifications for stress and emotion recognition.
\needspace{8\baselineskip}
{\tablestripes\tablefontsize
\begin{xltabular}{\textwidth}{@{} p{4cm} L @{}}
\caption{DEAP Dataset Specifications}\label{tab:deap} \\
\toprule
\thead{Parameter} & \thead{Value} \\
\midrule
\endfirsthead
\multicolumn{2}{@{}l}{\textit{(\,Tab.~\ref{tab:deap} continued from previous page\,)}} \\
\toprule
\thead{Parameter} & \thead{Value} \\
\midrule
\endhead
\bottomrule
\endlastfoot
Total Subjects & 32 \\
EEG Channels & 32 \\
Peripheral Channels & 8 \\
Sampling Rate & 512 Hz \\
Videos Watched & 40 per subject \\
Recording Length & 60s per trial \\
Labels & Valence, Arousal, Dominance, Liking \\
\end{xltabular}}

\section{SAM40 Dataset }

\noindent Table~\ref{tab:sam40} reports the SAM40 four-level stress dataset specifications.
\needspace{8\baselineskip}
{\tablestripes\tablefontsize
\begin{xltabular}{\textwidth}{@{} p{4cm} L @{}}
\caption{SAM40 Dataset Specifications}\label{tab:sam40} \\
\toprule
\thead{Parameter} & \thead{Value} \\
\midrule
\endfirsthead
\multicolumn{2}{@{}l}{\textit{(\,Tab.~\ref{tab:sam40} continued from previous page\,)}} \\
\toprule
\thead{Parameter} & \thead{Value} \\
\midrule
\endhead
\bottomrule
\endlastfoot
Total Subjects & 40 \\
EEG Channels & 32 \\
Sampling Rate & 128 Hz \\
Stress Levels & 4 (baseline, low, medium, high) \\
Protocol & Stroop task, arithmetic \\
\end{xltabular}}

\section{BCI Competition IV Dataset 2a}

\noindent Table~\ref{tab:bci} summarises the BCI Competition IV Dataset~2a specifications for motor imagery classification.
\needspace{8\baselineskip}
{\tablestripes\tablefontsize
\begin{xltabular}{\textwidth}{@{} p{4cm} L @{}}
\caption{BCI Competition IV 2a Specifications}\label{tab:bci} \\
\toprule
\thead{Parameter} & \thead{Value} \\
\midrule
\endfirsthead
\multicolumn{2}{@{}l}{\textit{(\,Tab.~\ref{tab:bci} continued from previous page\,)}} \\
\toprule
\thead{Parameter} & \thead{Value} \\
\midrule
\endhead
\bottomrule
\endlastfoot
Total Subjects & 9 \\
EEG Channels & 22 \\
EOG Channels & 3 \\
Sampling Rate & 250 Hz \\
Classes & 4 (left hand, right hand, feet, tongue) \\
Trials per Session & 288 \\
Sessions & 2 per subject \\
\end{xltabular}}

\section{TCIA Cancer Radiomics Dataset}

\noindent Table~\ref{tab:tcia} documents the TCIA cancer radiomics dataset parameters.
\needspace{8\baselineskip}
{\tablestripes\tablefontsize
\begin{xltabular}{\textwidth}{@{} p{4cm} L @{}}
\caption{TCIA Dataset Specifications}\label{tab:tcia} \\
\toprule
\thead{Parameter} & \thead{Value} \\
\midrule
\endfirsthead
\multicolumn{2}{@{}l}{\textit{(\,Tab.~\ref{tab:tcia} continued from previous page\,)}} \\
\toprule
\thead{Parameter} & \thead{Value} \\
\midrule
\endhead
\bottomrule
\endlastfoot
Total Patients & 500+ \\
Imaging Modality & CT, MRI, PET \\
Cancer Types & Lung, Breast, Brain, Liver \\
Image Format & DICOM \\
Annotations & Tumor segmentation masks \\
Features Extracted & 150+ radiomic features \\
\end{xltabular}}


%% ====================================================================
%% PREPROCESSING PIPELINES (from old Appendix)
%% ====================================================================

\label{app:preprocessing}

\section{EEG Preprocessing Pipeline}

The EEG preprocessing pipeline consists of five sequential stages, each designed to transform raw electrophysiological recordings into analysis-ready feature vectors. Table~\ref{tab:eeg_pipeline} summarizes the complete pipeline with input specifications, processing operations, output formats, and key parameters for each stage.
\needspace{8\baselineskip}
{\tablestripes\tablefontsize
\begin{xltabular}{\textwidth}{@{} p{0.4cm} p{1.8cm} p{1.8cm} L p{2.5cm} @{}}
\caption{EEG Preprocessing Pipeline Stages}\label{tab:eeg_pipeline} \\
\toprule
\thead{\#} & \thead{Stage} & \thead{Input} & \thead{Operation} & \thead{Output} \\
\midrule
\endfirsthead
\multicolumn{5}{@{}l}{\textit{(\,Tab.~\ref{tab:eeg_pipeline} continued from previous page\,)}} \\
\toprule
\thead{\#} & \thead{Stage} & \thead{Input} & \thead{Operation} & \thead{Output} \\
\midrule
\endhead
\bottomrule
\endlastfoot
1 & Data loading and quality check & Raw EDF files & Load multichannel EEG data; identify corrupted channels (standard deviation $< 10^{-10}$); mark and exclude bad channels & Validated raw EEG object with bad channel annotations \\
2 & Bandpass and notch filtering & Validated raw EEG & Apply 4th-order Butterworth bandpass filter (0.5--45 Hz); apply notch filter at 50 Hz to remove power line noise & Filtered EEG signal with physiological frequency range preserved \\
3 & ICA artifact removal & Filtered EEG & Fit Independent Component Analysis (20 components, seed = 42); automatically detect EOG artifacts via correlation; exclude artifact components & Clean EEG signal with eye and muscle artifacts removed \\
4 & Epoch segmentation & Clean continuous EEG & Segment into fixed-length epochs (4.0 s duration, 0.5 s overlap); apply baseline correction; reject epochs exceeding amplitude threshold & Set of artifact-free epochs \\
5 & Feature extraction & Epoch set & Compute power spectral density (Welch's method, 0.5--45 Hz); extract band powers: delta (0.5--4 Hz), theta (4--8 Hz), alpha (8--13 Hz), beta (13--30 Hz), gamma (30--45 Hz); compute spectral entropy & Numeric feature vector per epoch (6 features per channel) \\
\end{xltabular}}
\vspace{0.5em}\noindent\textit{Note.} EDF = European Data Format; ICA = Independent Component Analysis; EOG = electrooculography; PSD = power spectral density. The MNE-Python library was used for all EEG processing operations.

\section{Imaging Preprocessing Pipeline}

The medical imaging preprocessing pipeline transforms raw DICOM series into quantitative radiomic feature sets suitable for classification. Table~\ref{tab:imaging_pipeline} details each stage of this pipeline.
\needspace{8\baselineskip}
{\tablestripes\tablefontsize
\begin{xltabular}{\textwidth}{@{} p{0.4cm} p{1.8cm} p{1.8cm} L p{2.5cm} @{}}
\caption{Medical Imaging Preprocessing Pipeline Stages}\label{tab:imaging_pipeline} \\
\toprule
\thead{\#} & \thead{Stage} & \thead{Input} & \thead{Operation} & \thead{Output} \\
\midrule
\endfirsthead
\multicolumn{5}{@{}l}{\textit{(\,Tab.~\ref{tab:imaging_pipeline} continued from previous page\,)}} \\
\toprule
\thead{\#} & \thead{Stage} & \thead{Input} & \thead{Operation} & \thead{Output} \\
\midrule
\endhead
\bottomrule
\endlastfoot
1 & DICOM loading & Raw DICOM directory & Load DICOM image series; reconstruct 3D volumetric representation from sequential slices & 3D image volume with spatial metadata \\
2 & Intensity normalization & 3D image volume & Apply z-score normalization (subtract mean, divide by standard deviation) or min-max scaling to [0, 1] range & Intensity-standardized 3D volume \\
3 & Region of interest segmentation & Normalized volume + segmentation mask & Apply binary mask to isolate tumor or anatomical region; extract masked sub-volume & Isolated region of interest volume \\
4 & Radiomic feature extraction & ROI volume + mask & Extract 150+ quantitative features across five classes: first-order statistics, gray-level co-occurrence matrix (GLCM), gray-level run-length matrix (GLRLM), gray-level size-zone matrix (GLSZM), and 3D shape descriptors & Numeric feature vector per patient \\
\end{xltabular}}
\vspace{0.5em}\noindent\textit{Note.} DICOM = Digital Imaging and Communications in Medicine; ROI = region of interest; GLCM = gray-level co-occurrence matrix; GLRLM = gray-level run-length matrix; GLSZM = gray-level size-zone matrix. The SimpleITK and PyRadiomics libraries were used for imaging operations.

\section{Time-Frequency Transformation}

Two time-frequency transformation methods are employed to convert one-dimensional temporal signals into two-dimensional spectrograms suitable for CNN-based analysis. Table~\ref{tab:timefreq_methods} compares the two methods.
\needspace{8\baselineskip}
{\tablestripes\tablefontsize
\begin{xltabular}{\textwidth}{@{} p{2.8cm} L L @{}}
\caption{Time-Frequency Transformation Methods}\label{tab:timefreq_methods} \\
\toprule
\thead{Attribute} & \thead{Short-Time Fourier Transform (STFT)} & \thead{Continuous Wavelet Transform (CWT)} \\
\midrule
\endfirsthead
\multicolumn{3}{@{}l}{\textit{(\,Tab.~\ref{tab:timefreq_methods} continued from previous page\,)}} \\
\toprule
\thead{Attribute} & \thead{Short-Time Fourier Transform (STFT)} & \thead{Continuous Wavelet Transform (CWT)} \\
\midrule
\endhead
\bottomrule
\endlastfoot
Input & Single-channel EEG signal & Single-channel EEG signal \\
Window/wavelet & Hann window (256 samples) & Morlet wavelet \\
Scale/frequency range & Determined by sampling rate and window size & 1--128 scales \\
Time resolution & Fixed (determined by window length) & Adaptive (better at high frequencies) \\
Frequency resolution & Fixed (determined by window length) & Adaptive (better at low frequencies) \\
Output & 2D spectrogram (frequency $\times$ time) & 2D scalogram (scale $\times$ time) \\
Strength & Uniform resolution across all frequencies & Multi-resolution analysis of non-stationary signals \\
Application in this study & Baseline spectral representation & Primary representation for CNN-LSTM models \\
\end{xltabular}}
\vspace{0.5em}\noindent\textit{Note.} STFT = Short-Time Fourier Transform; CWT = Continuous Wavelet Transform; CNN = convolutional neural network; LSTM = long short-term memory. Both methods produce magnitude-only representations (absolute values of complex coefficients).


%% ====================================================================
%% MODEL ARCHITECTURES (from old Appendix)
%% ====================================================================

\label{app:architectures}

\section{1D-CNN Architecture for Raw EEG}

\noindent Table~\ref{tab:1dcnn} presents the layer-by-layer specification of the 1D-CNN architecture used for raw EEG classification.
\needspace{8\baselineskip}
{\tablestripes\tablefontsize
\begin{xltabular}{\textwidth}{@{} p{2.2cm} p{2.5cm} L R @{}}
\caption{1D-CNN Architecture}\label{tab:1dcnn} \\
\toprule
\thead{Layer} & \thead{Type} & \thead{Output Shape} & \thead{Parameters} \\
\midrule
\endfirsthead
\multicolumn{4}{@{}l}{\textit{(\,Tab.~\ref{tab:1dcnn} continued from previous page\,)}} \\
\toprule
\thead{Layer} & \thead{Type} & \thead{Output Shape} & \thead{Parameters} \\
\midrule
\endhead
\bottomrule
\endlastfoot
Input & Input & (None, 1000, 32) & 0 \\
Conv1D\_1 & Conv1D & (None, 996, 64) & 10,304 \\
BatchNorm1 & BatchNorm & (None, 996, 64) & 256 \\
MaxPool1 & MaxPool1D & (None, 498, 64) & 0 \\
Conv1D\_2 & Conv1D & (None, 494, 128) & 41,088 \\
BatchNorm2 & BatchNorm & (None, 494, 128) & 512 \\
MaxPool2 & MaxPool1D & (None, 247, 128) & 0 \\
Conv1D\_3 & Conv1D & (None, 243, 256) & 164,096 \\
Global\-Avg\-Pool & Global\-Avg\-Pool1D & (None, 256) & 0 \\
Dense1 & Dense & (None, 128) & 32{,}896 \\
Dropout & Dropout (0.5) & (None, 128) & 0 \\
Output & Dense (softmax) & (None, 2) & 258 \\
\midrule
\multicolumn{3}{l}{\textbf{Total Parameters}} & 249{,}410 \\
\end{xltabular}}

\section{CNN-LSTM Hybrid Architecture}

\noindent Table~\ref{tab:cnnlstm} documents the CNN-LSTM hybrid architecture with bidirectional LSTM and self-attention layers.
\needspace{8\baselineskip}
{\tablestripes\tablefontsize
\begin{xltabular}{\textwidth}{@{} p{2.2cm} p{2.5cm} L R @{}}
\caption{CNN-LSTM Architecture}\label{tab:cnnlstm} \\
\toprule
\thead{Layer} & \thead{Type} & \thead{Output Shape} & \thead{Parameters} \\
\midrule
\endfirsthead
\multicolumn{4}{@{}l}{\textit{(\,Tab.~\ref{tab:cnnlstm} continued from previous page\,)}} \\
\toprule
\thead{Layer} & \thead{Type} & \thead{Output Shape} & \thead{Parameters} \\
\midrule
\endhead
\bottomrule
\endlastfoot
Input & Input & (None, 128, 128, 1) & 0 \\
Conv2D\_1 & Conv2D & (None, 126, 126, 32) & 320 \\
MaxPool2D\_1 & MaxPool2D & (None, 63, 63, 32) & 0 \\
Conv2D\_2 & Conv2D & (None, 61, 61, 64) & 18,496 \\
MaxPool2D\_2 & MaxPool2D & (None, 30, 30, 64) & 0 \\
Conv2D\_3 & Conv2D & (None, 28, 28, 128) & 73,856 \\
MaxPool2D\_3 & MaxPool2D & (None, 14, 14, 128) & 0 \\
Reshape & Reshape & (None, 14, 1792) & 0 \\
Bi-LSTM & Bi\-directional LSTM & (None, 14, 256) & 1{,}970{,}176 \\
Attention & Self-Attention & (None, 256) & 65{,}792 \\
Dense1 & Dense & (None, 128) & 32{,}896 \\
Dropout & Dropout (0.5) & (None, 128) & 0 \\
Output & Dense (softmax) & (None, 2) & 258 \\
\midrule
\multicolumn{3}{l}{\textbf{Total Parameters}} & 2,161,794 \\
\end{xltabular}}

\section{DeepConvNet Architecture for BCI}

\noindent Table~\ref{tab:deepconvnet} reports the DeepConvNet architecture designed for BCI motor imagery classification.
\needspace{8\baselineskip}
{\tablestripes\tablefontsize
\begin{xltabular}{\textwidth}{@{} p{2.2cm} p{2.5cm} L R @{}}
\caption{DeepConvNet Architecture}\label{tab:deepconvnet} \\
\toprule
\thead{Layer} & \thead{Type} & \thead{Output Shape} & \thead{Parameters} \\
\midrule
\endfirsthead
\multicolumn{4}{@{}l}{\textit{(\,Tab.~\ref{tab:deepconvnet} continued from previous page\,)}} \\
\toprule
\thead{Layer} & \thead{Type} & \thead{Output Shape} & \thead{Parameters} \\
\midrule
\endhead
\bottomrule
\endlastfoot
Input & Input & (None, 1000, 22, 1) & 0 \\
Conv2D\_temp & Conv2D (25, 1) & (None, 976, 22, 25) & 650 \\
Conv2D\_spat & Conv2D (1, 22) & (None, 976, 1, 25) & 13,775 \\
BatchNorm1 & BatchNorm & (None, 976, 1, 25) & 100 \\
ELU & Activation & (None, 976, 1, 25) & 0 \\
MaxPool2D\_1 & MaxPool2D (3,1) & (None, 325, 1, 25) & 0 \\
Dropout\_1 & Dropout (0.5) & (None, 325, 1, 25) & 0 \\
Conv2D\_2 & Conv2D (10, 1) & (None, 316, 1, 50) & 12,550 \\
BatchNorm2 & BatchNorm & (None, 316, 1, 50) & 200 \\
ELU & Activation & (None, 316, 1, 50) & 0 \\
MaxPool2D\_2 & MaxPool2D (3,1) & (None, 105, 1, 50) & 0 \\
Flatten & Flatten & (None, 5250) & 0 \\
Dense & Dense & (None, 4) & 21,004 \\
\midrule
\multicolumn{3}{l}{\textbf{Total Parameters}} & 48,279 \\
\end{xltabular}}

\section{3D-CNN Architecture for Radiomics}

\noindent Table~\ref{tab:3dcnn} presents the 3D-CNN architecture for volumetric cancer radiomics classification.
\needspace{8\baselineskip}
{\tablestripes\tablefontsize
\begin{xltabular}{\textwidth}{@{} p{2.2cm} p{2.5cm} L R @{}}
\caption{3D-CNN Architecture for Cancer Radiomics}\label{tab:3dcnn} \\
\toprule
\thead{Layer} & \thead{Type} & \thead{Output Shape} & \thead{Parameters} \\
\midrule
\endfirsthead
\multicolumn{4}{@{}l}{\textit{(\,Tab.~\ref{tab:3dcnn} continued from previous page\,)}} \\
\toprule
\thead{Layer} & \thead{Type} & \thead{Output Shape} & \thead{Parameters} \\
\midrule
\endhead
\bottomrule
\endlastfoot
Input & Input & (None, 64, 64, 64, 1) & 0 \\
Conv3D\_1 & Conv3D & (None, 62, 62, 62, 32) & 896 \\
MaxPool3D\_1 & MaxPool3D & (None, 31, 31, 31, 32) & 0 \\
Conv3D\_2 & Conv3D & (None, 29, 29, 29, 64) & 55{,}360 \\
MaxPool3D\_2 & MaxPool3D & (None, 14, 14, 14, 64) & 0 \\
Conv3D\_3 & Conv3D & (None, 12, 12, 12, 128) & 221{,}312 \\
MaxPool3D\_3 & MaxPool3D & (None, 6, 6, 6, 128) & 0 \\
Global\-Avg\-Pool3D & GlobalAvgPool & (None, 128) & 0 \\
Dense1 & Dense & (None, 256) & 33{,}024 \\
Dropout & Dropout (0.5) & (None, 256) & 0 \\
Output & Dense (softmax) & (None, 2) & 514 \\
\midrule
\multicolumn{3}{l}{\textbf{Total Parameters}} & 311{,}106 \\
\end{xltabular}}




%% ====================================================================
%% ENGINEERING CONTENT MOVED FROM CHAPTER 3 (2026-03-31)
%% Each subsection below preserves the original content with app_ prefixed labels
%% ====================================================================

% [SPACE-FIX] clearpage removed to eliminate empty page
\subsection{Signal Standardisation Protocol}
\label{app:sec1_signal_standardisation}

%% SUPPRESSED FOR WORD COUNT
Table~\ref{tab:eeg_bands} documents the five canonical EEG frequency bands, their clinical significance, and their relevance to the five clinical domains in this study.
\needspace{3\baselineskip}
\noindent Table~\ref{app_tab:eeg_bands} eeg frequency bands.

{\tablestripes\tablefontsize
\begin{xltabular}{\textwidth}{@{} p{2cm} p{1.8cm} p{2.5cm} L L @{}}
\caption[EEG Frequency Bands]{EEG Frequency Bands: Neurophysiological Basis and Clinical Domain Relevance}\label{app_tab:eeg_bands} \\
\toprule
\thead{Band} & \thead{Range} & \thead{Brain State} & \thead{Clinical Significance} & \thead{Domain Relevance} \\
\midrule
\endfirsthead
\multicolumn{5}{@{}l}{\textit{(\,Tab.~\ref{app_tab:eeg_bands} continued from previous page\,)}} \\
\toprule
\thead{Band} & \thead{Range} & \thead{Brain State} & \thead{Clinical Significance} & \thead{Domain Relevance} \\
\midrule
\endhead
\bottomrule
\endlastfoot
Delta ($\delta$) & 0.5--4 Hz & Deep sleep & Brain damage; encephalopathy & \\
\addlinespace
Theta ($\theta$) & 4--8 Hz & Drowsiness; memory & Cognitive decline; ADHD & ASD ($\theta$/$\beta$ ratio); Stress \\
\addlinespace
Alpha ($\alpha$) & 8--13 Hz & Relaxed wakefulness & Eyes-closed baseline; asymmetry & Stress (frontal asymmetry); BCI ($\mu$ rhythm) \\
\addlinespace
Beta ($\beta$) & 13--30 Hz & Active thinking; motor & Motor planning; concentration & ; BCI \\
\addlinespace
Gamma ($\gamma$) & 30--45 Hz & Cognitive processing & Perception; consciousness & ASD-focused (high-level cognition) \\
\end{xltabular}}
\vspace{0.5em}\noindent\textit{Note.} Hz = hertz; ADHD = attention deficit hyperactivity disorder; $\mu$ rhythm = sensorimotor rhythm in the alpha range (8--12\,Hz) specific to motor cortex. The 0.5--45\,Hz band-pass filter (Figure~\ref{fig:eeg_preproc}) retains all five bands while excluding high-frequency muscle artifacts ($>$45\,Hz) and DC drift ($<$0.5\,Hz). Gamma range is truncated at 45\,Hz to avoid contamination from power-line interference (50/60\,Hz). Clinical significance informed by~\cite{niedermeyer2005eeg}.

%% MOVED TO SUPPLEMENTARY VOLUME — EEG data attributes table, 10-20 electrode system
%% table, and Fourier transform family table (Tables tab:eeg_data_attributes,
%% tab:1020_system, tab:fourier_family)
\subsubsection{EEG Signal Characteristics and Acquisition Standards}
\label{app_sec:eeg_data_attributes}

%% SUPPRESSED FOR WORD COUNT

The signal processing pipeline rests on four mathematical foundations: the Continuous Wavelet Transform (Equation~\ref{eq:cwt}), the Short-Time Fourier Transform (Equation~\ref{eq:stft}), Sample Entropy (Equation~\ref{eq:sample_entropy}), and Phase-Locking Value (Equation~\ref{eq:plv}). Each formula is introduced at its point of use in the subsections below; Appendix~\ref{app:ch3_support} provides full derivations and parameter sensitivity analyses.

\paragraph{Continuous Wavelet Transform (CWT).}
The CWT decomposes the EEG signal $x(t)$ using a Morlet wavelet:
\begin{equation}
W(a,b) = \frac{1}{\sqrt{a}} \int_{-\infty}^{\infty} x(t) \, \psi^{*}\!\left(\frac{t-b}{a}\right) dt
\label{app_eq:cwt}
\end{equation}
where $\psi(t) = \pi^{-1/4} e^{i\omega_0 t} e^{-t^2/2}$ is the Morlet wavelet ($\omega_0 = 6$), $a$ is the scale parameter, and $b$ is the translation parameter. The framework uses 64 logarithmically spaced frequency bins spanning 0.5--45~Hz with 128 time samples per epoch, producing 64$\times$128 scalogram images per channel.

\paragraph{Short-Time Fourier Transform (STFT).}
The STFT applies a sliding Hann window ($N = 256$ samples) with 75\% overlap and 512-point FFT:
\begin{equation}
X(m, \omega) = \sum_{n=0}^{N-1} x(n + mH) \, w(n) \, e^{-j\omega n}
\label{app_eq:stft}
\end{equation}
where $w(n)$ is the Hann window, $m$ is the frame index, and $H = N/4$ is the hop size. This produces 129$\times T$ spectrogram images, where $T$ depends on epoch length.

Table~\ref{tab:cwt_stft_specs} compares the two representations and their domain assignments.
\needspace{3\baselineskip}
\noindent Table~\ref{app_tab:cwt_stft_specs} time-frequency transformation specifications.

{\tablestripes\tablefontsize
\begin{xltabular}{\textwidth}{@{} p{2.5cm} L L @{}}
\caption{Time-Frequency Transformation Specifications}\label{app_tab:cwt_stft_specs} \\
\toprule
\thead{Parameter} & \thead{CWT (Morlet)} & \thead{STFT (Hann)} \\
\midrule
\endfirsthead
\multicolumn{3}{@{}l}{\textit{(\,Tab.~\ref{app_tab:cwt_stft_specs} continued from previous page\,)}} \\
\toprule
\thead{Parameter} & \thead{CWT (Morlet)} & \thead{STFT (Hann)} \\
\midrule
\endhead
\bottomrule
\endlastfoot
Wavelet/window & Morlet ($\omega_0 = 6$) & Hann ($N = 256$) \\
\addlinespace
Frequency range & 0.5--45 Hz (64 log bins) & 0--128 Hz (129 linear bins) \\
\addlinespace
Time resolution & 128 samples & Depends on epoch length \\
\addlinespace
Output shape & 64 $\times$ 128 per channel & 129 $\times$ $T$ per channel \\
\addlinespace
Overlap & N/A (continuous) & 75\% (hop = 64 samples) \\
\addlinespace
FFT size & N/A & 512 points \\
\addlinespace
Assigned model & 2D-CNN (scalograms; compared model) & Transformer (spectrograms; compared model) \\
\addlinespace
Clinical use & Multi-scale oscillatory patterns; alpha/beta/gamma band analysis & Drowsiness-specific alpha/theta transitions; steady-state evoked potentials \\
\end{xltabular}}
\vspace{0.5em}\noindent\textit{Note.} CWT = continuous wavelet transform; STFT = short-time Fourier transform; FFT = fast Fourier transform; Hz = hertz; CNN = convolutional neural network. CWT provides superior frequency resolution at low frequencies; STFT provides uniform time-frequency resolution.

Figure~\ref{fig:eeg_preproc} details the eight-step pipeline. Procedural descriptions are in Appendix~\ref{app:ch3_support}.
\needspace{18\baselineskip}
\begin{figure}[!htbp]
\centering
\resizebox{0.97\textwidth}{!}{%
\begin{tikzpicture}[
  pstep/.style={draw=ggunavy, fill=ggunavy!8, rounded corners=3pt, text width=13.5cm, minimum height=0.9cm, font=\small, align=left, inner sep=6pt},
  arr/.style={-{Stealth[length=4mm, width=3mm]}, very thick, ggunavy},
  node distance=0.8cm
]
\node[pstep, fill=lightblue] (s1) {\textbf{Step 1: Quality Check} --- Identify corrupted/missing recordings {\small (MNE-Python visual inspection)}};
\node[pstep, fill=lightorange, below=of s1] (s2) {\textbf{Step 2: Band-pass Filtering} --- Retain 0.5--45\,Hz; remove drift and noise {\small (Butterworth 4th-order)}};
\node[pstep, fill=lightteal, below=of s2] (s3) {\textbf{Step 3: Notch Filter} --- Remove power-line interference (50/60\,Hz) {\small (FIR notch filter)}};
\node[pstep, fill=lightpurple, below=of s3] (s4) {\textbf{Step 4: ICA Artifact Removal} --- Separate and remove ocular/muscular artifacts {\small (FastICA algorithm; ICLabel classification, $>$80\% non-brain rejection)}};
\node[pstep, fill=lightgreen, below=of s4] (s5) {\textbf{Step 5: Baseline Correction} --- Remove residual DC offset {\small (Epoch-mean subtraction)}};
\node[pstep, fill=lightyellow, below=of s5] (s6) {\textbf{Step 6: Normalisation} --- Standardise amplitude distributions {\small (Z-score normalisation)}};
\node[pstep, fill=lightpink, below=of s6] (s7) {\textbf{Step 7: Segmentation} --- Create fixed-length epochs (1--4\,s) {\small (Sliding window, 50\% overlap)}};
\node[pstep, fill=lightsky, below=of s7] (s8) {\textbf{Step 8: Transformation} --- Generate spectrograms, scalograms, CSP features {\small (STFT, CWT, CSP)}};
\draw[arr] (s1) -- (s2);
\draw[arr] (s2) -- (s3);
\draw[arr] (s3) -- (s4);
\draw[arr] (s4) -- (s5);
\draw[arr] (s5) -- (s6);
\draw[arr] (s6) -- (s7);
\draw[arr] (s7) -- (s8);
\end{tikzpicture}%
}% end resizebox
\caption{Eight-Step EEG Preprocessing Pipeline}
\label{app_fig:eeg_preproc}
\vspace{0.5em}\noindent\textit{Note.} ICA = independent component analysis; FIR = finite impulse response; CSP = common spatial pattern; STFT = short-time Fourier transform; CWT = continuous wavelet transform; Hz = hertz; s = seconds.
\end{figure}

\paragraph{ICA Component Rejection Criteria.} Independent components were classified using ICLabel~\cite{pion2019iclabel}; those with $>$80\% probability of non-brain origin were rejected across five artefact categories (Table~\ref{app_tab:ica_rejection}). This automated classification reduces subjective bias in manual artefact rejection and ensures reproducibility across datasets.
\needspace{8\baselineskip}
{\tablestripes\tablefontsize
\begin{xltabular}{\textwidth}{@{} p{2.5cm} L L @{}}
\caption{ICA Component Rejection Criteria: Five Artefact Categories}\label{app_tab:ica_rejection} \\
\toprule
\thead{Artefact Category} & \thead{Identification Signature} & \thead{Physiological Source} \\
\midrule
\endfirsthead
\multicolumn{3}{@{}l}{\textit{(\,Tab.~\ref{app_tab:ica_rejection} continued from previous page\,)}} \\
\toprule
\thead{Artefact Category} & \thead{Identification Signature} & \thead{Physiological Source} \\
\midrule
\endhead
\bottomrule
\endlastfoot
Ocular & Frontal topography; characteristic low-frequency time course & Eye blinks and saccades \\
\addlinespace
Muscular & Broadband high-frequency spectral profile & Electromyographic contamination \\
\addlinespace
Cardiac & Periodic QRS-complex morphology & Heart-related electrical activity \\
\addlinespace
Line noise & Narrowband spectral peaks at 50/60~Hz & Power-line interference \\
\addlinespace
Channel noise & Spatially focal topography & Single-channel electrode malfunction \\
\end{xltabular}}
\vspace{0.5em}\noindent\textit{Note.} ICA = independent component analysis; ICLabel threshold: $>$80\% probability of non-brain origin triggers automatic rejection~\cite{pion2019iclabel}. Hz = hertz.

\noindent Domain-specific ICA implementations varied by dataset: FastICA for ASD; 20-component ICA for PD; wavelet denoising complementing ICA for Stress; and amplitude thresholding as the primary rejection method for BCI (due to the motor-imagery paradigm's distinct artifact profile).

%% ============================================================================
%% EEG SIGNAL PROCESSING PIPELINE (from EEG-RAG paper)
%% ============================================================================
%% MOVED TO APPENDIX: Redundant EEG Signal Processing Pipeline (fig:eeg_signal_pipeline, tab:eeg_filter_params) — overlaps with fig:eeg_preproc

%% SUPPRESSED — redundant with fig:eeg_preproc (TikZ version preferred over PNG)
\iffalse
\noindent To provide an integrated view of how raw EEG recordings are transformed into model-ready feature vectors, Figure~\ref{fig:eeg_pipeline_visual} presents the end-to-end preprocessing and feature extraction pipeline. This pipeline operationalises the artefact removal methods described above and connects them to the feature engineering stages that follow.
\noindent The EEG signal preprocessing and feature extraction pipeline visual is presented in the main text as Figure~\ref{fig:eeg_pipeline_visual} (Chapter~\\ref{chap:methods}); not reproduced here to avoid duplication.
\fi



% [SPACE-FIX] clearpage removed to eliminate empty page
\subsection{PBS Image Preprocessing Pipeline}
\label{app:sec2_pbs_preproc}

%% SUPPRESSED FOR WORD COUNT
Unlike EEG preprocessing, which addresses temporal artefacts, the cancer domain requires spatial normalisation of microscopy images. Table~\ref{tab:img_preproc} summarises the five-step PBS image preprocessing pipeline, designed to standardise staining variation and class imbalance before classification.
\needspace{8\baselineskip}
\noindent Table~\ref{app_tab:img_preproc} five-step pbs image preprocessing pipeline.

{\tablestripes\tablefontsize
\begin{xltabular}{\textwidth}{@{} p{0.8cm} L L @{}}
\caption{Five-Step PBS Image Preprocessing Pipeline}\label{app_tab:img_preproc} \\
\toprule
\thead{Step} & \thead{Process} & \thead{Purpose} \\
\midrule
\endfirsthead
\multicolumn{3}{@{}l}{\textit{(\,Tab.~\ref{app_tab:img_preproc} continued from previous page\,)}} \\
\toprule
\thead{Step} & \thead{Process} & \thead{Purpose} \\
\midrule
\endhead
\bottomrule
\endlastfoot
1 & Resize to 512$\times$512 & Standardise spatial dimensions across images \\
2 & HSV colour thresholding & Segment cell regions from background; normalise stain intensity \\
3 & Normalisation to [0,1] & Scale pixel intensities for neural network input \\
4 & GAN augmentation & Generate synthetic training samples to balance classes and improve generalisation \\
5 & Data splitting & Stratified 80/20 train/test with 10-fold cross-validation \\
\end{xltabular}}
\vspace{0.5em}\noindent\textit{Note.} HSV = hue--saturation--value colour space; GAN = generative adversarial network; PBS = peripheral blood smear.

\noindent The five-step sequence is ordered deliberately: spatial standardisation (Steps~1--2) precedes intensity normalisation (Step~3), which precedes augmentation (Step~4), ensuring that GAN-generated images inherit the same normalised properties as real samples. This ordering prevents data leakage by applying augmentation only after the train/test split (Step~5), a safeguard that directly supports the internal validity claims tested in Chapter~\ref{chap:analysis}.

%% MOVED TO SUPPLEMENTARY VOLUME — Computer vision pipeline table (tab:cv_pipeline)
%% and signal vs image comparison table (tab:signal_vs_image)
\subsubsection{PBS Classification Pipeline}
\label{app_sec:cv_pipeline}

PBS microscopy images undergo parallel preprocessing: resizing to 512$\times$512, HSV colour thresholding for stain normalisation, pixel intensity scaling to [0,1], and GAN-based augmentation. The GAN generator (ConvTranspose2d layers) synthesises realistic cell images, while the discriminator provides auxiliary classification, and a ResNet-18 classifier performs the final 4-class ALL classification. Full pipeline specifications are in the Supplementary Volume.

%% ============================================================================
%% GAN ARCHITECTURE FOR CANCER DETECTION (from Cancer paper)
%% ============================================================================

Figure~\ref{fig:gan_pipeline} illustrates the end-to-end GAN-based PBS cancer classification pipeline, from input image preprocessing through GAN-based augmentation to the final four-class ALL output. The GAN addresses class imbalance by synthesising realistic cell images (generator) and provides auxiliary classification signals through the discriminator's learned features.
\needspace{20\baselineskip}
\begin{figure}[!htbp]
\centering
\resizebox{0.97\textwidth}{!}{%
\begin{tikzpicture}[
  pbox/.style={rectangle, draw=ggunavy, fill=ggunavy!8, rounded corners=3pt,
    text width=2.8cm, minimum height=1.6cm, align=center, font=\small, inner sep=5pt},
  gbox/.style={rectangle, draw=ggunavy, fill=ggunavy!18, rounded corners=3pt,
    text width=3.2cm, minimum height=1.6cm, align=center, font=\small, inner sep=5pt},
  obox/.style={rectangle, draw=ggunavy, fill=ggunavy!30, rounded corners=3pt,
    text width=2.6cm, minimum height=1.6cm, align=center, font=\small\bfseries, inner sep=5pt},
  arr/.style={-{Stealth[length=4mm, width=3mm]}, very thick, ggunavy},
  node distance=0.8cm
]
% Row 1: Input processing
\node[pbox, fill=lightblue] (input) {\textbf{\color{ggunavy}Input PBS}\\[2pt]20,000 images\\512$\times$512 RGB};
\node[pbox, fill=lightorange, right=of input] (resize) {\textbf{\color{ggunavy}Resize}\\[2pt]512$\times$512\\standardise};
\node[pbox, fill=lightteal, right=of resize] (norm) {\textbf{\color{ggunavy}Normalise}\\[2pt]$[0,1]$ range\\HSV threshold};

% Row 1: GAN
\node[gbox, fill=lightpurple, right=of norm] (gan) {\textbf{\color{ggunavy}GAN}\\[2pt]Generator:\\151,808 params\\Discriminator:\\262,297 params};

% Row 2: Classification
\node[pbox, fill=lightgreen, right=of gan] (resnet) {\textbf{\color{ggunavy}ResNet-18}\\[2pt]Classifier\\6,149,220 params};
\node[obox, fill=lightyellow, right=of resnet] (output) {\textbf{\color{ggunavy}4-Class Output}\\[2pt]Benign\\Early Pre-B\\Pre-B, Pro-B};

% Arrows
\draw[arr] (input) -- (resize);
\draw[arr] (resize) -- (norm);
\draw[arr] (norm) -- (gan);
\draw[arr] (gan) -- (resnet);
\draw[arr] (resnet) -- (output);
\end{tikzpicture}%
}
\caption{GAN-Based PBS Cancer Classification Pipeline}
\label{app_fig:gan_pipeline}
\vspace{0.5em}\noindent\textit{Note.} PBS = peripheral blood smear; GAN = generative adversarial network; RGB = red--green--blue; HSV = hue--saturation--value. The generator uses ConvTranspose2d layers to synthesise realistic cell images; the discriminator provides auxiliary classification features. Total pipeline: 6,563,325 trainable parameters. \textit{Source:}~\cite{asthana2025cancer}.
\end{figure}

Table~\ref{tab:gan_model_params} documents the parameter counts for each pipeline component. The ResNet-18 classifier accounts for 93.7\% of the total parameters, reflecting the transfer learning strategy where ImageNet-pretrained convolutional features are fine-tuned on PBS microscopy images. The GAN generator and discriminator together comprise only 414,105 parameters (6.3\%), making the augmentation component computationally lightweight.
\needspace{3\baselineskip}
\noindent Table~\ref{app_tab:gan_model_params} gan cancer classification pipeline: model component parameters.

{\tablestripes\tablefontsize
\begin{xltabular}{\textwidth}{@{} p{3.5cm} r p{2cm} L @{}}
\caption{GAN Cancer Classification Pipeline: Model Component Parameters}\label{app_tab:gan_model_params} \\
\toprule
\thead{Component} & \thead{Parameters} & \thead{Share (\%)} & \thead{Function} \\
\midrule
\endfirsthead
\multicolumn{4}{@{}l}{\textit{(\,Tab.~\ref{app_tab:gan_model_params} continued from previous page\,)}} \\
\toprule
\thead{Component} & \thead{Parameters} & \thead{Share (\%)} & \thead{Function} \\
\midrule
\endhead
\bottomrule
\endlastfoot
Generator & 151,808 & 2.3 & Synthesise realistic PBS cell images via ConvTranspose2d upsampling layers to augment training data \\
\addlinespace
Discriminator & 262,297 & 4.0 & Distinguish real from generated images; provide auxiliary classification features for the downstream classifier \\
\addlinespace
ResNet-18 classifier & 6,149,220 & 93.7 & Four-class ALL classification using ImageNet-pretrained features fine-tuned on PBS microscopy \\
\midrule
\textbf{Total} & \textbf{6,563,325} & \textbf{100.0} & \\
\end{xltabular}}
\vspace{0.5em}\noindent\textit{Note.} ALL = acute lymphoblastic leukaemia; PBS = peripheral blood smear; GAN = generative adversarial network. Parameter counts from~\cite{asthana2025cancer}. The generator uses five ConvTranspose2d layers (64$\to$512 channels); the discriminator mirrors this with five Conv2d layers. ResNet-18 is initialised with ImageNet weights and fine-tuned with a learning rate of $10^{-4}$.

%% ============================================================================
%% GAN INTERNAL ARCHITECTURE (from published cancer paper: Asthana et al.)
%% ============================================================================
Figure~\ref{fig:gan_architecture} presents the internal architecture of the three GAN components---Generator, Discriminator, and Classifier---showing the layer-by-layer progression from latent noise vector to four-class ALL classification output.
\needspace{18\baselineskip}
\begin{figure}[!htbp]
\centering
\resizebox{0.97\textwidth}{!}{%
\begin{tikzpicture}[
  genbox/.style={rectangle, draw=ggunavy, fill=lightpurple, rounded corners=3pt,
    thick, minimum width=4.2cm, minimum height=1.0cm,
    align=center, font=\small, inner sep=3pt},
  disbox/.style={rectangle, draw=ggunavy, fill=lightorange, rounded corners=3pt,
    thick, minimum width=4.2cm, minimum height=1.0cm,
    align=center, font=\small, inner sep=3pt},
  clsbox/.style={rectangle, draw=ggunavy, fill=lightgreen, rounded corners=3pt,
    thick, minimum width=4.2cm, minimum height=1.0cm,
    align=center, font=\small, inner sep=3pt},
  hdr/.style={rectangle, draw=ggunavy, fill=ggunavy!15, rounded corners=3pt,
    thick, minimum width=4.4cm, minimum height=0.55cm,
    align=center, font=\small\bfseries, inner sep=3pt},
  arr/.style={-{Stealth[length=4mm, width=3mm]}, very thick, ggunavy},
  node distance=0.8cm,
]

% ===== GENERATOR (left column) =====
\node[hdr] (gh) at (0,0) {Generator (G) --- 151,808 params};
\node[genbox, below=0.6cm of gh] (g1) {Input: Latent $z \sim \mathcal{N}(0,1)$, dim=100};
\node[genbox, below=of g1] (g2) {ConvTranspose2d (latent$\to$512), 4$\times$4};
\node[genbox, below=of g2] (g3) {BatchNorm2d(512) + ReLU};
\node[genbox, below=of g3] (g4) {ConvTranspose2d (512$\to$256), stride 2};
\node[genbox, below=of g4] (g5) {ConvTranspose2d (256$\to$128), stride 2};
\node[genbox, below=of g5] (g6) {ConvTranspose2d (128$\to$64), stride 2};
\node[genbox, below=of g6] (g7) {ConvTranspose2d (64$\to$3), Tanh};
\node[genbox, fill=lightblue, below=0.15cm of g7] (gout) {Output: 64$\times$64$\times$3 synthetic PBS image};
\draw[arr] (g1) -- (g2);
\draw[arr] (g3) -- (g4);
\draw[arr] (g5) -- (g6);
\draw[arr] (g6) -- (g7);
\draw[arr] (g7) -- (gout);

% ===== DISCRIMINATOR (middle column) =====
\node[hdr] (dh) at (5.5,0) {Discriminator (D) --- 262,297 params};
\node[disbox, below=0.6cm of dh] (d1) {Input: 64$\times$64$\times$3 image};
\node[disbox, below=of d1] (d2) {Conv2d (3$\to$64), LeakyReLU(0.2)};
\node[disbox, below=of d2] (d3) {Conv2d (64$\to$128), BN, LeakyReLU};
\node[disbox, below=of d3] (d4) {Conv2d (128$\to$256), BN, LeakyReLU};
\node[disbox, below=of d4] (d5) {Conv2d (256$\to$512), BN, LeakyReLU};
\node[disbox, below=of d5] (d6) {Adversarial Head: Conv2d (512$\to$1)};
\node[disbox, fill=lightyellow, below=0.15cm of d6] (dout) {Output: Real/Fake (Sigmoid)};
\draw[arr] (d1) -- (d2);
\draw[arr] (d3) -- (d4);
\draw[arr] (d4) -- (d5);
\draw[arr] (d5) -- (d6);
\draw[arr] (d6) -- (dout);

% ===== CLASSIFIER (right column) =====
\node[hdr] (ch) at (11,0) {ResNet-18 Classifier --- 6,149,220 params};
\node[clsbox, below=0.6cm of ch] (c1) {Input: 64$\times$64$\times$3 image};
\node[clsbox, below=of c1] (c2) {Initial Conv2d (3$\to$64), BN, ReLU};
\node[clsbox, below=of c2] (c3) {Layer1 (BasicBlock $\times$ 2)};
\node[clsbox, below=of c3] (c4) {Layer2 (BasicBlock $\times$ 2)};
\node[clsbox, below=of c4] (c5) {Layer3 (BasicBlock $\times$ 2)};
\node[clsbox, below=of c5] (c6) {Layer4 (BasicBlock $\times$ 2)};
\node[clsbox, below=of c6] (c7) {AdaptiveAvgPool2d + Flatten};
\node[clsbox, below=of c7] (c8) {Fully Connected (512$\to$4)};
\node[clsbox, fill=lightblue, below=0.15cm of c8] (cout) {Output: 4-class (Benign/Early/Pre/Pro)};
\draw[arr] (c1) -- (c2);
\draw[arr] (c3) -- (c4);
\draw[arr] (c5) -- (c6);
\draw[arr] (c7) -- (c8);
\draw[arr] (c8) -- (cout);

% ===== Inter-component arrows =====
\draw[arr, dashed, ggunavy!60] (gout.east) -- ++(0.4,0) |- (d1.west)
  node[pos=0.25, right, font=\small\itshape\color{ggunavy!60}] {generated};
\draw[arr, dashed, ggunavy!60] (dout.east) -- ++(0.4,0) |- ([yshift=-0.15cm]c1.west)
  node[pos=0.25, right, font=\small\itshape\color{ggunavy!60}] {features};

\end{tikzpicture}%
}
\caption[GAN Internal Architecture]{GAN Internal Architecture: Generator, Discriminator, and ResNet-18 Classifier for PBS Cancer Classification}
\label{app_fig:gan_architecture}
\vspace{0.5em}\noindent\textit{Note.} The Generator upsamples a 100-dimensional latent vector $z$ through five transposed convolutional layers to produce 64$\times$64 synthetic PBS images. The Discriminator distinguishes real from generated images using five convolutional layers with LeakyReLU activation. The ResNet-18 Classifier performs 4-class ALL subtype classification using ImageNet-pretrained weights fine-tuned with cross-entropy loss. Total parameters: 6,563,325. BN = batch normalisation. \textit{Source:} Asthana, P., Lalawat, R.\,S., Gond, S.\,S., \& Anand, S.\,K. (2025). A GAN-Based Framework for the Classification of Hematological Malignancies Using Peripheral Blood Smear Images.
\end{figure}

%% ============================================================================
%% 15-STEP PROCESSING FLOWCHART (from published cancer paper: Asthana et al.)
%% ============================================================================
Figure~\ref{fig:gan_15step} presents the complete 15-step processing sequence for GAN-based PBS image classification, showing the numbered flow from raw image loading through adversarial training to the final four-class output.
\needspace{20\baselineskip}
\begin{figure}[!htbp]
\centering
\resizebox{0.97\textwidth}{!}{%
\begin{tikzpicture}[
  sbox/.style={rectangle, draw=ggunavy, fill=ggunavy!8, rounded corners=3pt,
    thick, minimum width=2.4cm, minimum height=1.1cm,
    align=center, font=\small, inner sep=4pt},
  decision/.style={diamond, draw=ggunavy, fill=lightyellow, aspect=2,
    thick, align=center, font=\small, inner sep=2pt},
  terminal/.style={rectangle, draw=ggunavy, fill=ggunavy!20, rounded corners=8pt,
    thick, minimum width=1.6cm, minimum height=1.0cm,
    align=center, font=\small\bfseries, inner sep=3pt},
  num/.style={circle, draw=ggunavy, fill=ggunavy, text=white,
    font=\small\bfseries, minimum size=0.4cm, inner sep=0pt},
  arr/.style={-{Stealth[length=4mm, width=3mm]}, very thick, ggunavy},
  node distance=0.8cm
]

% Row 1: Steps 1-5 (preprocessing)
\node[terminal] (start) {Start};
\node[sbox, right=of start] (s1) {Load PBS\\Images};
\node[num, above left=0.01cm and 0.01cm of s1.north east] {1};
\node[sbox, right=of s1] (s2) {Resize to\\512$\times$512};
\node[num, above left=0.01cm and 0.01cm of s2.north east] {2};
\node[sbox, right=of s2] (s3) {Normalise\\$[0,1]$};
\node[num, above left=0.01cm and 0.01cm of s3.north east] {3};
\node[sbox, right=of s3] (s4) {Data\\Augment.};
\node[num, above left=0.01cm and 0.01cm of s4.north east] {4};
\node[sbox, right=of s4] (s5) {Train/Val\\Split};
\node[num, above left=0.01cm and 0.01cm of s5.north east] {5};

\draw[arr] (start) -- (s1);
\draw[arr] (s1) -- (s2);
\draw[arr] (s2) -- (s3);
\draw[arr] (s3) -- (s4);
\draw[arr] (s4) -- (s5);

% Row 2: Steps 6-10 (GAN training)
\node[sbox, below=0.8cm of s1] (s6) {Sample\\$z \sim \mathcal{N}(0,1)$};
\node[num, above left=0.01cm and 0.01cm of s6.north east] {6};
\node[sbox, fill=lightpurple, right=of s6] (s7) {Generator\\$G(z)$};
\node[num, above left=0.01cm and 0.01cm of s7.north east] {7};
\node[sbox, fill=lightorange, right=of s7] (s8) {Discriminator\\$D(x)$};
\node[num, above left=0.01cm and 0.01cm of s8.north east] {8};
\node[decision, right=0.6cm of s8] (s9) {Adv.\\Loss?};
\node[num, above=0.01cm of s9.north] {9};
\node[sbox, right=0.6cm of s9] (s10) {Update\\$G$ \& $D$};
\node[num, above left=0.01cm and 0.01cm of s10.north east] {10};

\draw[arr] (s5.south) -- ++(0,-0.3) -| (s6.north);
\draw[arr] (s6) -- (s7);
\draw[arr] (s7) -- (s8);
\draw[arr] (s8) -- (s9);
\draw[arr] (s9) -- node[above, font=\small] {Yes} (s10);
\draw[arr, dashed] (s10.south) -- ++(0,-0.3) -| (s7.south);

% Row 3: Steps 11-15 (classification)
\node[sbox, fill=lightgreen, below=0.8cm of s6] (s11) {Classifier\\$C(x)$};
\node[num, above left=0.01cm and 0.01cm of s11.north east] {11};
\node[sbox, right=of s11] (s12) {Softmax\\Output};
\node[num, above left=0.01cm and 0.01cm of s12.north east] {12};
\node[sbox, right=of s12] (s13) {Predict\\Class};
\node[num, above left=0.01cm and 0.01cm of s13.north east] {13};
\node[sbox, fill=lightblue, right=of s13] (s14) {Output:\\Benign/Early/\\Pre/Pro};
\node[num, above left=0.01cm and 0.01cm of s14.north east] {14};
\node[terminal, right=of s14] (s15) {End};
\node[num, above=0.01cm of s15.north] {15};

\draw[arr] (s9.south) -- node[right, font=\small] {No} ++(0,-0.5) -| (s11.north);
\draw[arr] (s11) -- (s12);
\draw[arr] (s12) -- (s13);
\draw[arr] (s13) -- (s14);
\draw[arr] (s14) -- (s15);

\end{tikzpicture}%
}
\caption[GAN Processing Flowchart]{Complete 15-Step Processing Flowchart for GAN-Based PBS Cancer Classification}
\label{app_fig:gan_15step}
\vspace{0.5em}\noindent\textit{Note.} Steps 1--5 handle preprocessing (resize, normalise, augment, split). Steps 6--10 perform GAN adversarial training (sample latent vector, generate synthetic images, discriminate, compute adversarial loss, update weights). Steps 11--15 perform classification (ResNet-18 classifier, softmax, predict 4-class ALL output). The adversarial loop (steps 7--10) iterates until convergence (50 epochs). PBS = peripheral blood smear. \textit{Source:}~\cite{asthana2025cancer}.
\end{figure}



% [SPACE-FIX] clearpage removed to eliminate empty page
\subsection{Feature Engineering and Disease-Specific Biomarkers}
\label{app:sec3_feature_eng}

%% SUPPRESSED FOR WORD COUNT
The feature engineering stage transforms preprocessed signals into discriminative representations. Table~\ref{tab:feature_eng} summarizes seven feature categories, their constituent features, applicable domains, and clinical meaning.
\needspace{8\baselineskip}
\noindent Table~\ref{app_tab:feature_eng} feature engineering summary across seven categories.

{\tablestripes\tablefontsize
\begin{xltabular}{\textwidth}{@{} p{1.8cm} L p{1.6cm} L @{}}
\caption{Feature Engineering Summary Across Seven Categories}\label{app_tab:feature_eng} \\
\toprule
\thead{Category} & \thead{Features} & \thead{Domains} & \thead{Clinical Meaning} \\
\midrule
\endfirsthead
\multicolumn{4}{@{}l}{\textit{(\,Tab.~\ref{app_tab:feature_eng} continued from previous page\,)}} \\
\toprule
\thead{Category} & \thead{Features} & \thead{Domains} & \thead{Clinical Meaning} \\
\midrule
\endhead
\bottomrule
\endlastfoot
Spectral & PSD, band powers (delta through gamma) & All EEG & Dominant brain oscillation states; abnormal patterns indicate dysfunction \\
Band ratios & Alpha/beta, theta/beta & Stress, ASD & Cognitive load and attentional regulation \\
Time-frequency & Wavelet coefficients~\cite{mallat1989wavelet}, STFT~\cite{welch1967use} & All EEG & Transient, non-stationary neural events \\
Entropy & Sample entropy~\cite{richman2000sampleentropy}, spectral entropy & ASD, PD & Signal complexity; reduced entropy suggests pathological rigidity \\
Spatial & CSP components~\cite{ramoser2000csp} & BCI & Maximizes variance between motor imagery classes \\
Connectivity & Coherence, PLV~\cite{stam2007coherence} & ASD, PD & Inter-regional neural synchrony; disrupted in ASD and PD \\
GAN-learned & Discriminator features, ResNet-18 deep features & Cancer & Cell morphology, staining patterns, and nuclear characteristics \\
\end{xltabular}}
\vspace{0.5em}\noindent\textit{Note.} PSD = power spectral density; CSP = common spatial pattern; PLV = phase-locking value; GAN = generative adversarial network; STFT = short-time Fourier transform; ASD = autism spectrum disorder;

%% SUPPRESSED FOR WORD COUNT
Table~\ref{tab:feature_127} provides the complete breakdown of 127 EEG features extracted per channel, organised into five categories.
\needspace{8\baselineskip}
\noindent Table~\ref{app_tab:feature_127} eeg feature extraction (127 features).

{\tablestripes\tablefontsize
\begin{xltabular}{\textwidth}{@{} p{1.8cm} c L p{2.5cm} @{}}
\caption[EEG Feature Extraction (127 Features)]{Complete EEG Feature Extraction: 127 Features Per Channel in Five Categories}\label{app_tab:feature_127} \\
\toprule
\thead{Category} & \thead{Count} & \thead{Features} & \thead{Extraction Method} \\
\midrule
\endfirsthead
\multicolumn{4}{@{}l}{\textit{(\,Tab.~\ref{app_tab:feature_127} continued from previous page\,)}} \\
\toprule
\thead{Category} & \thead{Count} & \thead{Features} & \thead{Extraction Method} \\
\midrule
\endhead
\bottomrule
\endlastfoot
Time-domain & 25 & Mean, variance, std, RMS, skewness, kurtosis, peak-to-peak, zero crossings, line length, energy, Hjorth (activity, mobility, complexity), percentiles (5th, 25th, 50th, 75th, 95th) & Signal statistics; Hjorth parameters via Eq.~\ref{eq:hjorth_activity} \\
\addlinespace
Entropy & 15 & Shannon, sample ($m$=2, $r$=0.2$\sigma$), spectral, permutation (order 3--5), approximate, fuzzy, multiscale & Non-linear dynamics \\
\addlinespace
Frequency & 45 & Band powers ($\delta$, $\theta$, $\alpha$, $\beta$, $\gamma$), relative powers, peak frequencies, bandwidths, spectral edge (95\%), band ratios ($\theta$/$\beta$, $\alpha$/$\theta$, $\delta$/$\alpha$), spectral entropy, spectral centroid, spectral flatness, spectral rolloff & PSD via Welch method (256-point FFT, 50\% overlap Hamming window) \\
\addlinespace
Connectivity & 39 & Phase Locking Value (PLV), weighted PLI (wPLI), coherence, imaginary coherence, correlation, mutual information & Cross-channel dependencies; Eq.~\ref{eq:plv} \\
\addlinespace
Nonlinear & 3 & Hurst exponent, Lyapunov exponent, correlation dimension & Dynamical systems analysis \\
\midrule
\textbf{Total} & \textbf{127} & & \\
\end{xltabular}}
\vspace{0.5em}\noindent\textit{Note.} PSD = power spectral density; PLV = phase locking value; wPLI = weighted phase lag index; FFT = fast Fourier transform. Feature count of 127 applies to EEG domains (ASD, PD, Stress, BCI); cancer PBS classification uses GAN + ResNet-18 deep features.

%% SUPPRESSED FOR WORD COUNT

\textbf{Hjorth Parameters.} Activity, mobility, and complexity characterise signal dynamics:
\begin{equation}
\text{Activity} = \sigma^2(x), \quad
\text{Mobility} = \sqrt{\frac{\sigma^2(x')}{\sigma^2(x)}}, \quad
\text{Complexity} = \frac{\text{Mobility}(x')}{\text{Mobility}(x)}
\label{app_eq:hjorth_activity}
\end{equation}

\textbf{Sample Entropy.} Measures signal regularity ($m = 2$, $r = 0.2\sigma$):
\begin{equation}
H_{\text{SampEn}} = -\ln\frac{A^m(r)}{B^m(r)}
\label{app_eq:sample_entropy}
\end{equation}

\textbf{Phase Locking Value (PLV).} Quantifies inter-channel phase synchrony:
\begin{equation}
\text{PLV}_{ij} = \left| \frac{1}{T} \sum_{t=1}^{T} e^{i(\phi_i(t) - \phi_j(t))} \right|
\label{app_eq:plv}
\end{equation}
where $\phi_i(t)$ and $\phi_j(t)$ are the instantaneous phases of channels $i$ and $j$ obtained via the Hilbert transform.

Figure~\ref{fig:feature_pipeline} illustrates the complete eight-stage feature engineering pipeline, from raw signal acquisition through filtering, artefact removal, segmentation, feature extraction, selection, evaluation, and normalisation to model-ready feature matrices.
\needspace{20\baselineskip}
\begin{figure}[!htbp]
\centering
\resizebox{0.97\textwidth}{!}{%
\begin{tikzpicture}[
  stage/.style={rectangle, rounded corners=4pt, draw=ggunavy, fill=ggunavy!8, text width=3cm, minimum height=1.4cm, align=center, font=\small, inner sep=5pt},
  arr/.style={-{Stealth[length=4mm, width=3mm]}, very thick, ggunavy},
  note/.style={font=\small\itshape, text=ggunavy!60, text width=3cm, align=center}
]
%% Row 1
\node[stage, fill=lightblue] (s1) at (0,3) {\textbf{1. Raw Signal}\\{\small EEG / Image}\\{\small acquisition}};
\node[stage, fill=lightorange] (s2) at (4,3) {\textbf{2. Bandpass Filter}\\{\small 0.5--45~Hz}\\{\small Butterworth 4th}};
\node[stage, fill=lightorange] (s3) at (8,3) {\textbf{3. Artefact Removal}\\{\small ICA decomposition}\\{\small EOG/EMG reject}};
\node[stage, fill=lightteal] (s4) at (12,3) {\textbf{4. Segmentation}\\{\small Epoch windowing}\\{\small 50\% overlap}};
%% Row 2
\node[stage, fill=lightteal] (s5) at (12,0) {\textbf{5. Feature Extract.}\\{\small Time, freq.,}\\{\small connectivity}};
\node[stage, fill=lightpurple] (s6) at (8,0) {\textbf{6. Feature Select.}\\{\small MI, ANOVA, SHAP}\\{\small Top-$k$ ranking}};
\node[stage, fill=lightpurple] (s7) at (4,0) {\textbf{7. Feature Eval.}\\{\small Importance scores}\\{\small Stability check}};
\node[stage, fill=lightgreen] (s8) at (0,0) {\textbf{8. Normalisation}\\{\small Z-score / MinMax}\\{\small Domain-specific}};
%% Arrows
\draw[arr] (s1) -- (s2);
\draw[arr] (s2) -- (s3);
\draw[arr] (s3) -- (s4);
\draw[arr] (s4) -- ++(0,-1) -| (s5);
\draw[arr] (s5) -- (s6);
\draw[arr] (s6) -- (s7);
\draw[arr] (s7) -- (s8);
%% Output
\node[rectangle, rounded corners=3pt, draw=ggunavy, fill=ggunavy, text=white, text width=3cm, minimum height=1cm, align=center, font=\small\bfseries] (out) at (0,-2) {Model Input\\{\normalfont\scriptsize Feature matrix}};
\draw[arr] (s8) -- (out);
\end{tikzpicture}%
}%
\caption[Eight-stage feature engineering pipeline]{Eight-Stage Feature Engineering Pipeline: From Raw EEG/Image Acquisition Through Filtering, Artefact Removal, Segmentation, Feature Extraction, Selection, Evaluation, and Normalisation to Model-Ready Feature Matrices}
\label{app_fig:feature_pipeline}
\end{figure}

\noindent The pipeline's bidirectional flow---from raw signal acquisition (Stage~1) through feature extraction and selection (Stages~5--7) to normalised model input (Stage~8)---ensures that every domain passes through identical structural stages, even when the internal parameters differ. This architectural consistency is what enables ASD-focused performance comparisons in Chapter~\ref{chap:analysis}.

Table~\ref{tab:feature_selection} documents the feature selection methods applied per domain and the resulting dimensionality reduction. SHAP is used as the final feature importance arbiter for ASD, with average dimensionality reduction of 62\%.
\needspace{3\baselineskip}
\noindent Table~\ref{app_tab:feature_selection} feature selection methods and results.

{\tablestripes\tablefontsize
\begin{xltabular}{\textwidth}{@{} >{\raggedright\arraybackslash}p{0.14\textwidth} >{\raggedright\arraybackslash}X c c c @{}}
\caption[Feature selection methods and results]{Feature Selection Methods and Dimensionality Reduction by Domain}\label{app_tab:feature_selection} \\
\toprule
\thead{Domain} & \thead{Methods Applied} & \thead{Initial} & \thead{Selected} & \thead{Reduction} \\
\midrule
\endfirsthead
\multicolumn{5}{@{}l}{\textit{(\,Tab.~\ref{app_tab:feature_selection} continued from previous page\,)}} \\
\toprule
\thead{Domain} & \thead{Methods Applied} & \thead{Initial} & \thead{Selected} & \thead{Reduction} \\
\midrule
\endhead
\bottomrule
\endlastfoot
Schizophrenia & MI + ANOVA $F$-test + SHAP & 45 & 18 & 60\% \\
\addlinespace
Epilepsy & MI + LASSO ($\lambda$=0.01) + SHAP & 52 & 22 & 58\% \\
\addlinespace
Stress & MI + ANOVA + mRMR + SHAP & 48 & 20 & 58\% \\
\addlinespace
Autism (ASD) & ANOVA + recursive feat.\ elim.\ + SHAP & 38 & 15 & 61\% \\
\addlinespace
Parkinson  & MI + SHAP (top-$k$, $k$=12) & 42 & 12 & 71\% \\
\addlinespace
Depression & MI + LASSO + SHAP & 46 & 18 & 61\% \\
\addlinespace
Cancer & MI + PCA (95\% var.) + SHAP & 107 & 35 & 67\% \\
\end{xltabular}}
\vspace{0.5em}\noindent\textit{Note.} MI = mutual information; mRMR = minimum redundancy maximum relevance; LASSO = least absolute shrinkage and selection operator; PCA = principal component analysis; SHAP = SHapley Additive exPlanations. SHAP is used as the final feature importance arbiter for ASD. Average dimensionality reduction is 62\%.

\subsection{Disease-Specific EEG Biomarkers}
\label{app_sec:disease_biomarkers}

%% SUPPRESSED FOR WORD COUNT
Table~\ref{tab:disease_biomarkers} maps seven conditions to their primary EEG biomarkers, detection methods, and clinical rationale.
\needspace{8\baselineskip}
\noindent Table~\ref{app_tab:disease_biomarkers} disease-specific eeg biomarkers.

{\tablestripes\tablefontsize
\begin{xltabular}{\textwidth}{@{} p{1.4cm} p{1.2cm} L L L @{}}
\caption[Disease-Specific EEG Biomarkers]{Disease-Specific EEG Biomarkers: Detection Methods and Clinical Rationale}\label{app_tab:disease_biomarkers} \\
\toprule
\thead{Disease} & \thead{Key Bands} & \thead{Biomarker} & \thead{Detection} & \thead{Clinical Rationale} \\
\midrule
\endfirsthead
\multicolumn{5}{@{}l}{\textit{(\,Tab.~\ref{app_tab:disease_biomarkers} continued from previous page\,)}} \\
\toprule
\thead{Disease} & \thead{Key Bands} & \thead{Biomarker} & \thead{Detection} & \thead{Clinical Rationale} \\
\midrule
\endhead
\bottomrule
\endlastfoot
Epilepsy & $\delta$, $\theta$, $\gamma$ & Spike amplitude, high-freq.\ oscillations & Spike detection (75th pctl.) & Epileptiform discharges indicate seizure onset zones \\
\addlinespace
Parkinson & $\beta$, $\theta$ & Excessive $\beta$ synchrony, tremor freq.\ (4--6 Hz) & Spectral analysis (60th pctl.) & $\beta$-band hypersynchrony is the hallmark of PD \\
\addlinespace
Alzheimer & $\theta$, $\delta$ & Cortical slowing ($\theta$/$\delta$ ratio elevation) & Spectral ratio (70th pctl.) & Progressive cortical slowing reflects neurodegeneration \\
\addlinespace
Schizo\-phrenia & $\gamma$, $\theta$ & Impaired $\gamma$ coherence, $\theta$ elevation & Coherence analysis (65th pctl.) & Disrupted $\gamma$ binding reflects impaired neural integration \\
\addlinespace
Depression & $\alpha$, $\theta$ & Frontal $\alpha$ asymmetry (right $>$ left) & Asymmetry analysis (55th pctl.) & Right-hemisphere dominance indicates withdrawal motivation \\
\addlinespace
ASD & $\gamma$, $\alpha$ & Altered $\gamma$ connectivity, reduced $\alpha$ power & Connectivity analysis (70th pctl.) & Disrupted neural synchrony reflects social cognition deficits \\
\addlinespace
Stress & $\beta$, $\alpha$ & $\beta$/$\alpha$ ratio elevation, frontal $\alpha$ suppression & Band ratio (60th pctl.) & Elevated $\beta$ and suppressed $\alpha$ indicate heightened cortical arousal \\
\end{xltabular}}
\vspace{0.5em}\noindent\textit{Note.} $\delta$ = delta (0.5--4 Hz); $\theta$ = theta (4--8 Hz); $\alpha$ = alpha (8--13 Hz); $\beta$ = beta (13--30 Hz); $\gamma$ = gamma ($>$30 Hz).  ASD = autism spectrum disorder. Biomarker thresholds are expressed as percentiles of the training distribution. These biomarkers anchor both classification features and RAG-generated clinical explanations.

\noindent To illustrate how time-frequency features are extracted from raw EEG, Figure~\ref{fig:cwt_scalogram_example} presents a representative CWT scalogram showing the joint time--frequency energy distribution. This visualisation demonstrates why wavelet-based features capture transient neural events that standard spectral analysis may miss.

\noindent The CWT scalogram visualisation is presented in the main text as Figure~\ref{fig:cwt_scalogram_example} (Chapter~\ref{chap:methods}); not reproduced here to avoid duplication.

%% ============================================================================
%% EEG FEATURE EXTRACTION METHODS (from ASD + EEG-RAG papers)
%% ============================================================================
%% MOVED TO APPENDIX: Redundant EEG Feature Taxonomy table (tab:eeg_feature_categories) — overlaps with tab:feature_127 and tab:feature_eng

\subsection{ASD-Focused Feature Extraction}

%% SUPPRESSED FOR WORD COUNT
A key challenge for ASD-focused diagnosis is that different diseases produce distinct biomarker signatures across different feature domains. Table~\ref{tab:feature_comparison} systematically compares the feature extraction methods applied across all the ASD domain, revealing both shared analytical strategies and domain-specific adaptations that inform the RGAIG pipeline design.
\needspace{8\baselineskip}
\noindent Table~\ref{app_tab:feature_comparison} asd-focused feature extraction comparison.

{\tablestripes\tablefontsize
\begin{xltabular}{\textwidth}{@{} p{1.2cm} L L L L L @{}}
\caption{ASD-Focused Feature Extraction Comparison}\label{app_tab:feature_comparison} \\
\toprule
\thead{Feature\newline Domain} & \thead{ASD} & \thead{PD} & \thead{Cancer} & \thead{Stress} & \thead{BCI} \\
\midrule
\endfirsthead
\multicolumn{6}{@{}l}{\textit{(\,Tab.~\ref{app_tab:feature_comparison} continued from previous page\,)}} \\
\toprule
\thead{Feature\newline Domain} & \thead{ASD} & \thead{PD} & \thead{Cancer} & \thead{Stress} & \thead{BCI} \\
\midrule
\endhead
\bottomrule
\endlastfoot
Time domain & Hjorth parameters, zero-crossing, statistical moments & Hjorth parameters, waveform shape analysis & N/A (pixel-level) & Statistical moments, Hjorth parameters & Statistical moments, zero-crossing rate \\
\addlinespace
Frequency & Band power ($\delta$, $\theta$, $\alpha$, $\beta$, $\gamma$), PSD, spectral entropy & Beta power (13--30\,Hz), alpha power, $\theta$/$\beta$ ratio & N/A & Frontal $\alpha$ asym\-metry, frontal $\theta$, high~$\beta$ & Band power, PSD, spectral entropy \\
\addlinespace
Time-frequency & CWT scalogram & Multi-scale CWT (3 scales) & N/A & VMD + CWT & CWT spectrogram \\
\addlinespace
Spatial & Topographic maps & Channel attention maps & CNN feature maps (Conv1--5) & Channel-level features & Topographic maps \\
\addlinespace
Connec\-tivity & PLV, coherence, Granger causality & Beta coherence, patho\-logical synchrony & N/A & Limited (cross-channel) & Limited (cross-channel) \\
\addlinespace
Key bio\-marker & $\theta$/$\beta$ ratio elevation & Beta-band elevation and synchrony & Cell mor\-phology (blast ratio) & Frontal $\alpha$ asym\-metry & $\alpha$/$\theta$ ratio reduction \\
\end{xltabular}}
\vspace{0.5em}\noindent\textit{Note.} CWT = continuous wavelet transform; VMD = variational mode decomposition; PLV = phase-locking value; PSD = power spectral density; CNN = convolutional neural network; $\delta$ = delta (0.5--4 Hz); $\theta$ = theta (4--8 Hz); $\alpha$ = alpha (8--12 Hz); $\beta$ = beta (13--30 Hz); $\gamma$ = gamma ($>$30 Hz).

\noindent The comparison reveals that time-domain and frequency-domain features are shared across all four EEG domains, while connectivity features (PLV, coherence) are unique to ASD and PD---reflecting the distinct pathophysiology of these disorders. Cancer relies entirely on CNN-learned spatial features rather than engineered signal features, confirming the need for the dual-track architecture. Building on this ASD-focused feature landscape, Table~\ref{tab:signal_to_image} documents the signal-to-image conversion methods that enable CNN-based classification of originally non-image data.
\needspace{8\baselineskip}
\noindent Table~\ref{app_tab:signal_to_image} signal-to-image conversion methods across domains.

{\tablestripes\tablefontsize
\begin{xltabular}{\textwidth}{@{} p{1.4cm} L p{2.2cm} L @{}}
\caption{Signal-to-Image Conversion Methods Across Domains}\label{app_tab:signal_to_image} \\
\toprule
\thead{Domain} & \thead{Conversion Method} & \thead{Output Dims.} & \thead{Rationale} \\
\midrule
\endfirsthead
\multicolumn{4}{@{}l}{\textit{(\,Tab.~\ref{app_tab:signal_to_image} continued from previous page\,)}} \\
\toprule
\thead{Domain} & \thead{Conversion Method} & \thead{Output Dims.} & \thead{Rationale} \\
\midrule
\endhead
\bottomrule
\endlastfoot
ASD (EEG) & CWT scalogram (Morlet wavelet) & 224 $\times$ 224 $\times$ 3 (RGB) & Captures multi-scale oscillatory patterns; compatible with ImageNet-pretrained CNNs \\
\addlinespace
& Multi-scale CWT (3 frequency scales) & 3 $\times$ 64 $\times$ 64 & Separates $\alpha$, $\beta$, $\gamma$ band contributions for channel attention mechanism \\
\addlinespace
Cancer (Image) & Direct input (already image) & 224 $\times$ 224 $\times$ 3 (RGB) & Microscopy images require no signal conversion; stain normalisation applied instead \\
\addlinespace
Stress (EEG) & CWT + topographic map overlay & 128 $\times$ 128 $\times$ 3 (RGB) & Combines time-frequency and spatial activation patterns \\
\addlinespace
BCI (EEG) & Spectrogram (STFT) & 224 $\times$ 224 $\times$ 3 (RGB) & Standard short-time Fourier transform captures drowsiness-specific $\alpha$/$\theta$ transitions \\
\end{xltabular}}
\vspace{0.5em}\noindent\textit{Note.} CWT = continuous wavelet transform; STFT = short-time Fourier transform; CNN = convolutional neural network; RGB = red--green--blue colour channels.

\noindent All four EEG domains converge on 224$\times$224$\times$3 output dimensions despite using different conversion methods (CWT, spectrogram, STFT). This dimensional standardisation enables the same CNN architecture to process inputs from different domains---a design choice that directly supports the ASD-focused generalisation tested in Hypothesis~H2.

%% ============================================================================
%% WEARABLE SENSOR INTEGRATION (from Wearable paper)
%% ============================================================================
%% MOVED TO APPENDIX: Wearable Sensor Integration (tab:wearable_sensors, fig:edge_deployment) — no wearable data used in experiments



% [SPACE-FIX] clearpage removed to eliminate empty page
\subsection{Classical Baseline Classifiers and Ensemble Selection}
\label{app:sec4_classifiers}

%% SUPPRESSED FOR WORD COUNT

Table~\ref{tab:ml_classifiers} summarises the classical baseline classifier selection rationale, key parameters, and domain suitability.
\needspace{8\baselineskip}
\noindent Table~\ref{app_tab:ml_classifiers} classical baseline classifier selection.

{\tablestripes\tablefontsize
\begin{xltabular}{\textwidth}{@{} p{1.4cm} p{1.4cm} L L p{1.6cm} @{}}
\caption[Classical Baseline Classifier Selection]{Classical Baseline Classifier Selection: Rationale, Parameters, and Domain Suitability}\label{app_tab:ml_classifiers} \\
\toprule
\thead{Model} & \thead{Paradigm} & \thead{Why Selected} & \thead{Key Parameters} & \thead{Primary Domain} \\
\midrule
\endfirsthead
\multicolumn{5}{@{}l}{\textit{(\,Tab.~\ref{app_tab:ml_classifiers} continued from previous page\,)}} \\
\toprule
\thead{Model} & \thead{Paradigm} & \thead{Why Selected} & \thead{Key Parameters} & \thead{Primary Domain} \\
\midrule
\endhead
\bottomrule
\endlastfoot
SVM-RBF & Margin-based & Non-linear boundaries in high-dimensional feature spaces; structural risk minimisation guarantee~\cite{cortes1995svm} & $C \in \{0.1, 1, 10, 100\}$; $\gamma$: auto or grid $\{10^{-4}\text{--}10^{-1}\}$ & All EEG \\
\addlinespace
Random Forest & Ensemble bagging & Variance reduction via 300 bootstrapped trees; robust to overfitting on small $n$; built-in Gini importance~\cite{breiman2001randomforests} & 300 trees; balanced class weights; $\sqrt{p}$ features per split & All EEG (in Voting\-Classifier) \\
\addlinespace
ExtraTrees & Ensemble bagging & Extremely randomised splits reduce variance further than RF; faster training via random thresholds~\cite{geurts2006extremely} & 300 trees; balanced class weights; random split thresholds & All EEG (in Voting\-Classifier) \\
\addlinespace
XGBoost & Gradient boosting & L1+L2 regularisation; native missing-value handling; state-of-the-art for tabular data~\cite{chen2016xgboost} & 300 rounds; LR 0.1; early stop 10 rounds; L2 = 0.01 & Cancer PBS \\
\addlinespace
LDA & Linear projection & Closed-form solution; minimal parameters ($\sim$100); overfitting-resistant on small datasets~\cite{ramoser2000csp} & Automatic (Fisher criterion) & BCI (with CSP) \\
\end{xltabular}}
\vspace{0.5em}\noindent\textit{Note.} SVM = support vector machine; RBF = radial basis function; LDA = linear discriminant analysis; CSP = common spatial pattern; LR = learning rate; RF = Random Forest. Feature dimensionality ranges from 95 to 380 depending on domain. CSP + LDA has been the BCI competition benchmark since 2005. The VotingClassifier (ExtraTrees + RF, soft voting) is the best-performing and final deployed model for all EEG-based domains.

\subsubsection{VotingClassifier Ensemble}
\label{app_sec:voting_ensemble}

For all EEG-based classification tasks (ASD, PD, Stress), the best-performing model is a \textbf{VotingClassifier} ensemble combining ExtraTreesClassifier and RandomForestClassifier with soft voting. This ensemble was selected after systematic comparison with both classical baselines and deep learning architectures (Section~\ref{sec:dl_architectures}). Table~\ref{app_tab:voting_config} summarises the four configuration components and their parameters.
\needspace{8\baselineskip}
{\tablestripes\tablefontsize
\begin{xltabular}{\textwidth}{@{} p{2.5cm} L L @{}}
\caption{VotingClassifier Ensemble Configuration}\label{app_tab:voting_config} \\
\toprule
\thead{Component} & \thead{Key Parameters} & \thead{Rationale} \\
\midrule
\endfirsthead
\multicolumn{3}{@{}l}{\textit{(\,Tab.~\ref{app_tab:voting_config} continued from previous page\,)}} \\
\toprule
\thead{Component} & \thead{Key Parameters} & \thead{Rationale} \\
\midrule
\endhead
\bottomrule
\endlastfoot
ExtraTrees Classifier & 300 trees; balanced class weights; random split thresholds & Reduces variance beyond standard RF by randomising both feature selection and split points~\cite{geurts2006extremely} \\
\addlinespace
Random Forest Classifier & 300 trees; balanced class weights; $\sqrt{p}$ features per split & Bootstrap aggregation with feature subsampling provides diversity~\cite{breiman2001randomforests} \\
\addlinespace
Soft voting & Probability averaging across both estimators & Outperforms hard voting by 1.5 pp on average (consistent with 1--3\% margins in ensemble literature~\cite{kittler1998combining, zhou2012ensemble}); retains more information than discrete majority voting \\
\addlinespace
Balanced class weights & Inverse-frequency weighting (\textsf{\small class\_weight=\allowbreak{}'balanced'}) & Compensates for class imbalance in clinical datasets (e.g., ASD control:case ratio of 1.4:1) \\
\end{xltabular}}
\vspace{0.5em}\noindent\textit{Note.} RF = Random Forest; pp = percentage points. Soft voting margin measured as mean difference between soft-voting and hard-voting accuracy across 5-fold CV (Chapter~\ref{chap:analysis}).

The VotingClassifier construction follows five sequential steps from data splitting through final evaluation, as summarised in Table~\ref{app_tab:voting_steps}.
\needspace{8\baselineskip}
{\tablestripes\tablefontsize
\begin{xltabular}{\textwidth}{@{} p{0.4cm} p{1.8cm} L p{2.5cm} @{}}
\caption{VotingClassifier Construction: Five Sequential Steps}\label{app_tab:voting_steps} \\
\toprule
\thead{\#} & \thead{Step} & \thead{Procedure} & \thead{Key Parameters} \\
\midrule
\endfirsthead
\multicolumn{4}{@{}l}{\textit{(\,Tab.~\ref{app_tab:voting_steps} continued from previous page\,)}} \\
\toprule
\thead{\#} & \thead{Step} & \thead{Procedure} & \thead{Key Parameters} \\
\midrule
\endhead
\bottomrule
\endlastfoot
1 & Subject-level splitting & Stratified 5-fold CV partitions subjects (not epochs) into training and test folds & Seed = 42; no data leakage \\
\addlinespace
2 & Data augmentation & Gaussian noise injection within training folds increases effective training set size & 0.1--0.5\% of feature SD; ASD: $n$=19 subjects, 300 augmented epochs (3$\times$) \\
\addlinespace
3 & Feature engineering & Extract features across band powers, statistical features, and Hjorth parameters & 140+ features per sample; 5 band powers, 7 statistical, 3 Hjorth \\
\addlinespace
4 & Ensemble fitting & ExtraTrees and RF trained independently; soft voting averages class probability estimates & 300 trees each; balanced weights \\
\addlinespace
5 & Cross-validated evaluation & Performance aggregated across all five folds & Accuracy, AUC-ROC, F1; bootstrap CI (1,000 resamples) \\
\end{xltabular}}
\vspace{0.5em}\noindent\textit{Note.} CV = cross-validation; SD = standard deviation; RF = Random Forest; CI = confidence interval. Augmentation applied exclusively within training folds to prevent data leakage.

Classification results for the VotingClassifier across all EEG domains are reported in Chapter~\ref{chap:analysis}. The ensemble's design on EEG tabular features is consistent with the broader deep learning literature, which demonstrates that tree-based ensembles outperform deep learning on structured/tabular data when the feature set is well-engineered~\cite{chen2016xgboost}.


% [SPACE-FIX] clearpage removed to eliminate empty page
\subsection{Deep Learning Architectures}
\label{app:sec5_dl_arch}


The framework evaluates ten deep learning architectures as \textit{compared models} for EEG domains and as the \textit{best-performing model} for imaging domains . For EEG-based classification, DL architectures were systematically compared against the VotingClassifier ensemble (Section~\ref{sec:voting_ensemble}) but did not outperform it on the engineered feature sets used in RGAIG. DL remains the deployed model for cancer imaging, where spatial feature hierarchies in histopathological images favour convolutional architectures. Table~\ref{tab:dl_architecture_comparison} summarises the ten architectures, their input modalities, primary use cases, and deployment status.
\needspace{8\baselineskip}
\noindent Table~\ref{app_tab:dl_architecture_comparison} deep learning architecture comparison.

{\tablestripes\tablefontsize
\begin{xltabular}{\textwidth}{@{} p{0.3cm} L L L p{1.5cm} @{}}
\caption[Deep Learning Architecture Comparison]{Deep Learning Architecture Comparison: Input Type, Primary Use, and Deployment Status}\label{app_tab:dl_architecture_comparison} \\
\toprule
\thead{\#} & \thead{Architecture} & \thead{Input Type} & \thead{Primary Use} & \thead{Deploy.\ Status} \\
\midrule
\endfirsthead
\multicolumn{5}{@{}l}{\textit{(\,Tab.~\ref{app_tab:dl_architecture_comparison} continued from previous page\,)}} \\
\toprule
\thead{\#} & \thead{Architecture} & \thead{Input Type} & \thead{Primary Use} & \thead{Deploy.\ Status} \\
\midrule
\endhead
\bottomrule
\endlastfoot
1 & 1D-CNN & Raw EEG time-series & Temporal pattern extraction from single-channel signals & Compared \\
\addlinespace
2 & 2D-CNN & CWT scalograms & Spectro-temporal image classification for EEG & Compared \\
\addlinespace
3 & CNN-LSTM (bi-directional + self-attention) & STFT spectrograms & Captures both spatial and temporal EEG dynamics & Compared \\
\addlinespace
4 & DeepConvNet~\cite{schirrmeister2017deep} & Raw multi-channel EEG & End-to-end deep convolutional decoding of EEG & Compared \\
\addlinespace
5 & EEGNet~\cite{lawhern2018eegnet} & Raw multi-channel EEG & Compact, generalisable EEG classification across paradigms & Compared \\
\addlinespace
6 & GAN generator (ConvTranspose2d) & PBS microscopy images & Synthetic image generation for class balancing & Deployed  \\
\addlinespace
7 & ResNet-18 classifier & PBS microscopy images & 4-class ALL leukaemia classification from cell images & Deployed  \\
\addlinespace
8 & Vision Transformer (ViT) & Patch-based image tokens & Global self-attention over image patches for classification & Compared \\
\addlinespace
9 & Transformer encoder & EEG feature sequences & Self-attention over temporal EEG segments & Compared \\
\addlinespace
10 & MLP meta-learner (Ultra Stacking) & Stacked base-model predictions & Multi-disease screening via meta-learned ensemble & Deployed (multi-disease) \\
\end{xltabular}}
\vspace{0.5em}\noindent\textit{Note.} CNN = convolutional neural network; LSTM = long short-term memory; CWT = continuous wavelet transform; STFT = short-time Fourier transform; GAN = generative adversarial network; PBS = peripheral blood smear; ALL = acute lymphoblastic leukaemia; ViT = vision transformer; MLP = multilayer perceptron. ``Compared'' = evaluated but not deployed (VotingClassifier outperformed on EEG features); ``Deployed'' = selected as best-performing model for the indicated domain. Full architecture specifications are in Appendix~C.

%% ENGINEERING DETAIL SUPPRESSED — DBA summary retained; full specifications in Appendix~C
%% (Architecture parameters, equations, layer configurations)
\iffalse
%% BEGIN SUPPRESSED ENGINEERING: DL Ensemble + Ultra Stacking Architecture Details
\subsubsection{Three-Model DL Ensemble}
\label{app_sec:dl_ensemble_detail}
Three DL architectures evaluated: EEGNet (2,548 params), 2D-CNN (148,226 params), Transformer (502,018 params).
Ensemble prediction: $P_{\text{ensemble}} = \sum_{i=1}^{3} w_i \cdot p_i$. Total: 652,792 params.
\subsubsection{Ultra Stacking Ensemble (Multi-Disease Classification)}
\label{app_sec:ultra_stacking_detail}
Three-layer stacking: Layer 1 (15+ base classifiers), Layer 2 (mutual information feature selection, top 300), Layer 3 (MLP meta-learner 128-64-32).
%% END SUPPRESSED ENGINEERING
\fi

\label{app_sec:dl_ensemble}
\label{app_sec:ultra_stacking}

The framework's model selection strategy follows an adaptive decision principle: the best-performing architecture is selected per clinical domain based on empirical evidence, not architectural preference. Three ensemble strategies were evaluated (Table~\ref{tab:ensemble_comparison}): (1)~VotingClassifier for single-condition EEG detection, (2)~a three-model DL ensemble as a compared alternative, and (3)~Ultra Stacking for multi-disease screening. Classification results are reported in Chapter~\ref{chap:analysis}. This adaptive selection represents a governance-aligned decision: organisations deploy the architecture that maximises diagnostic accuracy for their specific clinical use case rather than defaulting to the most complex model.

Table~\ref{tab:ensemble_comparison} compares the three ensemble strategies and their deployment contexts.
\needspace{3\baselineskip}
\noindent Table~\ref{app_tab:ensemble_comparison} ensemble architecture comparison.

{\tablestripes\tablefontsize
\begin{xltabular}{\textwidth}{@{} p{2.3cm} L p{3cm} L @{}}
\caption[Ensemble Architecture Comparison]{Ensemble Architecture Comparison: VotingClassifier vs Three-Model DL vs Ultra Stacking}\label{app_tab:ensemble_comparison} \\
\toprule
\thead{Aspect} & \thead{VotingClassifier\newline(Deployed)} & \thead{Three-Model DL\newline(Compared)} & \thead{Ultra Stacking\newline(Multi-Disease)} \\
\midrule
\endfirsthead
\multicolumn{4}{@{}l}{\textit{(\,Tab.~\ref{app_tab:ensemble_comparison} continued from previous page\,)}} \\
\toprule
\thead{Aspect} & \thead{VotingClassifier\newline(Deployed)} & \thead{Three-Model DL\newline(Compared)} & \thead{Ultra Stacking\newline(Multi-Disease)} \\
\midrule
\endhead
\bottomrule
\endlastfoot
Base models & ExtraTrees (300) + RF (300) & EEGNet, 2D-CNN, Transformer & 15+ classifiers (tree, neural, kernel) \\
\addlinespace
Input & 140+ engineered features & Raw EEG, CWT scalograms, STFT spectrograms & 47 engineered features per channel \\
\addlinespace
Fusion & Soft voting (probability averaging) & Weighted probability averaging & Stacking with MLP meta-learner \\
\addlinespace
Validation & 5-fold stratified CV & LOSO-CV & 5-fold stratified CV \\
\addlinespace
Use case & Single-condition EEG detection & Alternative EEG architecture & Multi-disease screening \\
\addlinespace
Inference & $<$5 ms (CPU) & 32 ms (GPU) & 15 ms (CPU) \\
\addlinespace
Status & \textbf{Best-performing; deployed} & Compared; not deployed & Deployed for multi-disease \\
\end{xltabular}}
\vspace{0.5em}\noindent\textit{Note.} RF = Random Forest; LOSO = leave-one-subject-out; CV = cross-validation; DL = deep learning; CWT = continuous wavelet transform; STFT = short-time Fourier transform; MLP = multilayer perceptron. The VotingClassifier is the final deployed model for all EEG-based domains (ASD, PD, Stress). DL architectures were systematically compared but did not outperform the ensemble on engineered features.

The key managerial decision is model selection: the VotingClassifier ensemble is deployed for all EEG domains (ASD, PD, Stress) because it outperforms DL alternatives on engineered features, while GAN+ResNet-18 is deployed for cancer imaging where spatial feature hierarchies favour convolutional architectures. This evidence-based selection tests H2: whether governance-driven model selection---choosing the best-performing architecture per domain rather than defaulting to complexity---yields higher classification accuracy.

%% ENGINEERING FIGURE SUPPRESSED — CNN-LSTM architecture diagram moved to Appendix C
\iffalse
\noindent The CNN-LSTM hybrid architecture diagram is presented in the main text as Figure~\ref{fig:cnn_lstm_arch} (Chapter~\\ref{chap:methods}); not reproduced here to avoid duplication.
\fi
%% Phantom label for cross-references
\phantomsection\label{app_fig:cnn_lstm_arch}

%% ============================================================================
%% CNN-ASD MODEL ARCHITECTURE (from published ICSIT 2025 paper: Asthana et al.)
%% ============================================================================
\noindent The CNN-ASD model architecture diagram is presented in the main text as Figure~\ref{fig:cnn_asd_arch} (Chapter~\\ref{chap:methods}); not reproduced here to avoid duplication.

%% ============================================================================
%% ASD CLASSIFICATION PIPELINE (from published ICSIT 2025 paper)
%% ============================================================================
\noindent The EEG-based ASD classification pipeline diagram is presented in the main text as Figure~\ref{fig:asd_classification_pipeline} (Chapter~\\ref{chap:methods}); not reproduced here to avoid duplication.


% [SPACE-FIX] clearpage removed to eliminate empty page
\subsection{Training Configuration}
\label{app:sec6_training}


Figure~\ref{fig:eeg_disease_finder_arch} presents the end-to-end system architecture of the Agentic AI Disease Finder~\cite{asthana2025eegmaster}, which integrates four main components summarised in Table~\ref{app_tab:disease_finder_components}. The architecture was validated across 10 publicly available EEG datasets encompassing 6,087 subjects and 5 neurological conditions (stress, anxiety, affect, epilepsy, and motor imagery), achieving 87.3$\pm$4.8\% mean accuracy under leave-one-subject-out (LOSO) cross-validation---a stringent protocol that prevents data leakage across subjects.
\needspace{8\baselineskip}
\noindent Table~\ref{app_tab:disease_finder_components} agentic ai disease finder: four architectural components.

{\tablestripes\tablefontsize
\begin{xltabular}{\textwidth}{@{} p{0.5cm} p{2.5cm} L p{3cm} @{}}
\caption{Agentic AI Disease Finder: Four Architectural Components}\label{app_tab:disease_finder_components} \\
\toprule
\thead{\#} & \thead{Component} & \thead{Function} & \thead{Key Specifications} \\
\midrule
\endfirsthead
\multicolumn{4}{@{}l}{\textit{(\,Tab.~\ref{app_tab:disease_finder_components} continued from previous page\,)}} \\
\toprule
\thead{\#} & \thead{Component} & \thead{Function} & \thead{Key Specifications} \\
\midrule
\endhead
\bottomrule
\endlastfoot
1 & Signal preprocessing & Five sequential artefact-removal stages transform raw EEG into clean epochs & Bandpass 0.5--45\,Hz; ICA; notch 50/60\,Hz; CAR; segmentation \\
\addlinespace
2 & Feature extraction & Dual-path extraction combining engineered features with time-frequency representations & 127 features per channel; CWT 64$\times$128; STFT 129$\times T$ \\
\addlinespace
3 & Ensemble classification & Three complementary deep learning architectures combined via soft voting & EEGNet (2.5K params); 2D-CNN (148K); Transformer (502K) \\
\addlinespace
4 & RAG explanation & Evidence-grounded explanation generation linking predictions to peer-reviewed literature & 50K+ PubMed abstracts; FAISS retrieval; llama3.2:3b \\
\end{xltabular}}
\vspace{0.5em}\noindent\textit{Note.} ICA = independent component analysis; CAR = common average reference; CWT = continuous wavelet transform; STFT = short-time Fourier transform; RAG = retrieval-augmented generation; FAISS = Facebook AI Similarity Search. \textit{Source:}~\cite{asthana2025eegmaster}.
\needspace{20\baselineskip}
\begin{figure}[!htbp]
\centering
\resizebox{0.97\textwidth}{!}{%
\begin{tikzpicture}[
  node distance=0.8cm and 0.15cm,
  procbox/.style={rectangle, draw=ggunavy, fill=ggunavy!8, rounded corners=2pt,
    minimum height=0.9cm, minimum width=1.8cm, align=center, font=\small},
  featbox/.style={rectangle, draw=orange!70!black, fill=orange!10, rounded corners=2pt,
    minimum height=0.9cm, minimum width=1.8cm, align=center, font=\small},
  modbox/.style={rectangle, draw=teal!80!black, fill=teal!10, rounded corners=2pt,
    minimum height=0.9cm, minimum width=1.8cm, align=center, font=\small},
  ragbox/.style={rectangle, draw=purple!70!black, fill=purple!10, rounded corners=2pt,
    minimum height=0.9cm, minimum width=1.8cm, align=center, font=\small},
  outbox/.style={rectangle, draw=red!70!black, fill=red!10, rounded corners=2pt,
    minimum height=0.9cm, minimum width=2.0cm, align=center, font=\small},
  grplabel/.style={font=\small\bfseries, text=gray!70!black},
  arrow/.style={-{Stealth[length=4pt]}, thick, ggunavy}
]
%% Row 1: Preprocessing pipeline
\node[grplabel] at (-1.2,0) {Preprocessing};
\node[procbox, fill=blue!15] (raw) at (1.2,0) {\textbf{Raw EEG}\\$C \times T$};
\node[procbox, right=of raw] (resamp) {\textbf{Resample}\\256\,Hz};
\node[procbox, right=of resamp] (bandp) {\textbf{Bandpass}\\0.5--45\,Hz};
\node[procbox, right=of bandp] (notch) {\textbf{Notch}\\50/60\,Hz};
\node[procbox, right=of notch] (ica) {\textbf{ICA}\\Artefact};
\node[procbox, right=of ica] (car) {\textbf{CAR}\\Ref};
\node[procbox, right=of car] (seg) {\textbf{Segment}\\to Epochs};

\draw[arrow] (raw) -- (resamp);
\draw[arrow] (resamp) -- (bandp);
\draw[arrow] (bandp) -- (notch);
\draw[arrow] (notch) -- (ica);
\draw[arrow] (ica) -- (car);
\draw[arrow] (car) -- (seg);

%% Row 2: Features
\node[grplabel] at (-1.2,-1.8) {Features};
\node[featbox, fill=orange!20] (feat127) at (2.2,-1.8) {\textbf{127 Features}\\Time, Frequency,\\Entropy, Conn.};
\node[featbox, fill=yellow!20, right=2.5cm of feat127] (tfr) {\textbf{TF Repr.}\\CWT $64 \times 128$};
\node[featbox, fill=yellow!15, right=of tfr] (stft) {\textbf{STFT}\\$129 \times T$};

\draw[arrow] ([xshift=-1.0cm]seg.south) -- ++(0,-0.3) -| (feat127.north);
\draw[arrow] (seg.south) -- ++(0,-0.3) -| (tfr.north);
\draw[arrow] ([xshift=1.0cm]seg.south) -- ++(0,-0.3) -| (stft.north);

%% Row 3: Models
\node[grplabel] at (-1.2,-3.6) {Models};
\node[modbox, fill=green!15] (eegnet) at (1.5,-3.6) {\textbf{EEGNet}\\2.5K params};
\node[modbox, fill=teal!15, right=1.5cm of eegnet] (cnn2d) {\textbf{2D-CNN}\\148K params};
\node[modbox, fill=cyan!15, right=1.5cm of cnn2d] (transf) {\textbf{Transformer}\\502K params};
\node[modbox, fill=lime!15, right=1.5cm of transf] (ensem) {\textbf{Ensemble}\\Soft Voting};
\node[outbox, right=1.0cm of ensem] (pred) {\textbf{Prediction}\\Class + Conf.};

\draw[arrow] (feat127.south) -- ++(0,-0.3) -| (eegnet.north);
\draw[arrow] (tfr.south) -- ++(0,-0.5) -| (cnn2d.north);
\draw[arrow] (stft.south) -- ++(0,-0.3) -| (transf.north);
\draw[arrow] (eegnet.east) -- ++(0.3,0) |- (ensem.west);
\draw[arrow] (cnn2d.east) -- ++(0.2,0) |- (ensem.170);
\draw[arrow] (transf) -- (ensem);
\draw[arrow] (ensem) -- (pred);

%% Row 4: RAG Module
\node[grplabel] at (-1.2,-5.2) {RAG Module};
\node[ragbox, fill=violet!15] (embed) at (2.2,-5.2) {\textbf{Embed KB}\\50K abstracts};
\node[ragbox, fill=purple!15, right=1.5cm of embed] (faiss) {\textbf{FAISS Retrieve}\\Top-$K$=5};
\node[ragbox, fill=magenta!15, right=1.5cm of faiss] (llm) {\textbf{LLM Synthesise}\\llama3.2:3b};
\node[outbox, fill=pink!20, right=1.0cm of llm] (expl) {\textbf{Explanation}\\Biomarkers};

\draw[arrow] (pred.south) -- ++(0,-0.4) -| (embed.north);
\draw[arrow] (embed) -- (faiss);
\draw[arrow] (faiss) -- (llm);
\draw[arrow] (llm) -- (expl);

%% Group boxes (dashed)
\draw[dashed, gray, rounded corners=5pt] (-0.2,0.7) rectangle (13.5,-0.6);
\draw[dashed, gray, rounded corners=5pt] (-0.2,-1.1) rectangle (10.5,-2.6);
\draw[dashed, gray, rounded corners=5pt] (-0.2,-2.9) rectangle (14.5,-4.4);
\draw[dashed, gray, rounded corners=5pt] (-0.2,-4.5) rectangle (13.5,-5.9);

\end{tikzpicture}%
}
\caption[Agentic AI Disease Finder System Architecture]{End-to-End System Architecture of the Agentic AI Disease Finder: Raw EEG Undergoes Five-Stage Preprocessing, Dual-Path Feature Extraction (127 Engineered Features + CWT/STFT), Ensemble Classification (EEGNet, 2D-CNN, Transformer), and RAG-Based Explanation Generation}
\label{app_fig:eeg_disease_finder_arch}
\vspace{0.5em}\noindent\textit{Note.} CAR = common average reference; ICA = independent component analysis; CWT = continuous wavelet transform; STFT = short-time Fourier transform; FAISS = Facebook AI Similarity Search; KB = knowledge base; Conn.\ = connectivity features. EEGNet operates on raw filtered signals; 2D-CNN on CWT scalograms ($64 \times 128$); Transformer on STFT spectrograms ($129 \times T$). Ensemble uses soft voting with learned weights $w_i$ optimised on validation set. \textit{Source:}~\cite{asthana2025eegmaster}.
\end{figure}

%% ============================================================================
%% AGENTIC AI DISEASE FINDER — ROC AND PRECISION-RECALL CURVES
%% ============================================================================

Figure~\ref{fig:eeg_roc_pr} presents the receiver operating characteristic (ROC) and precision-recall curves for three representative datasets from the Agentic AI Disease Finder evaluation~\cite{asthana2025eegmaster}. The ROC analysis demonstrates that the ensemble achieves AUC values of 0.974 (CHB-MIT, epilepsy), 0.958 (SAM40, stress), and 0.872 (BCI-IV, motor imagery), indicating strong discriminative performance across neurological conditions. The precision-recall curves confirm sustained precision even at high recall thresholds, which is clinically essential for screening applications where false negatives carry greater cost than false positives.
\needspace{20\baselineskip}
\begin{figure}[!htbp]
\centering
\resizebox{0.97\textwidth}{!}{%
\begin{tikzpicture}
%% ----- (a) ROC Curves -----
\begin{scope}[shift={(0,0)}]
  \draw[-{Stealth[length=4mm,width=3mm]}, thick] (0,0) -- (6.2,0) node[right, font=\small] {FPR};
  \draw[-{Stealth[length=4mm,width=3mm]}, thick] (0,0) -- (0,6.2) node[above, font=\small] {TPR};
  %% Axis ticks
  \foreach \x/\lab in {0/0, 1.5/0.2, 3/0.4, 4.5/0.6, 6/0.8} {
    \draw (\x,0) -- (\x,-0.1) node[below, font=\small] {\lab};
  }
  \node[below, font=\small] at (6.8,-0.1) {1};
  \foreach \y/\lab in {0/0, 1.5/0.2, 3/0.4, 4.5/0.6, 6/0.8} {
    \draw (0,\y) -- (-0.1,\y) node[left, font=\small] {\lab};
  }
  \node[left, font=\small] at (-0.1,6.8) {1};
  %% Diagonal reference
  \draw[gray, dashed, thin] (0,0) -- (6.8,6.8);
  %% CHB-MIT (AUC=0.974) — steep rise
  \draw[blue!70!black, thick] (0,0) -- (0.1,3.0) -- (0.3,4.8) -- (0.6,5.6) -- (1.2,6.2) -- (3.0,6.6) -- (6.0,6.8) -- (6.8,6.8);
  %% SAM40 (AUC=0.958) — slightly lower
  \draw[red!70!black, thick] (0,0) -- (0.2,2.4) -- (0.5,4.2) -- (1.0,5.4) -- (1.8,6.0) -- (3.5,6.5) -- (6.0,6.7) -- (6.8,6.8);
  %% BCI-IV (AUC=0.872) — lower curve
  \draw[teal!80!black, thick] (0,0) -- (0.3,1.5) -- (0.8,3.0) -- (1.5,4.2) -- (2.5,5.2) -- (4.0,6.0) -- (5.5,6.5) -- (6.8,6.8);
  %% Legend
  \node[font=\small, blue!70!black] at (4.5,2.0) {CHB-MIT (0.974)};
  \node[font=\small, red!70!black] at (4.5,1.4) {SAM40 (0.958)};
  \node[font=\small, teal!80!black] at (4.5,0.8) {BCI-IV (0.872)};
  %% Legend lines
  \draw[blue!70!black, thick] (3.2,2.0) -- (3.8,2.0);
  \draw[red!70!black, thick] (3.2,1.4) -- (3.8,1.4);
  \draw[teal!80!black, thick] (3.2,0.8) -- (3.8,0.8);
  %% Title
  \node[font=\small\bfseries] at (3.4,7.5) {(a) ROC Curves};
\end{scope}

%% ----- (b) Precision-Recall Curves -----
\begin{scope}[shift={(9,0)}]
  \draw[-{Stealth[length=4mm,width=3mm]}, thick] (0,0) -- (6.2,0) node[right, font=\small] {Recall};
  \draw[-{Stealth[length=4mm,width=3mm]}, thick] (0,0) -- (0,6.2) node[above, font=\small] {Precision};
  %% Axis ticks
  \foreach \x/\lab in {0/0, 1.5/0.2, 3/0.4, 4.5/0.6, 6/0.8} {
    \draw (\x,0) -- (\x,-0.1) node[below, font=\small] {\lab};
  }
  \node[below, font=\small] at (6.8,-0.1) {1};
  \foreach \y/\lab in {0/0, 1.5/0.2, 3/0.4, 4.5/0.6, 6/0.8} {
    \draw (0,\y) -- (-0.1,\y) node[left, font=\small] {\lab};
  }
  \node[left, font=\small] at (-0.1,6.8) {1};
  %% CHB-MIT (high precision maintained)
  \draw[blue!70!black, thick] (0,6.8) -- (1.0,6.7) -- (2.5,6.5) -- (4.0,6.2) -- (5.0,5.8) -- (5.8,5.0) -- (6.5,3.8) -- (6.8,2.5);
  %% SAM40 (slightly lower)
  \draw[red!70!black, thick] (0,6.6) -- (1.0,6.4) -- (2.5,6.0) -- (4.0,5.5) -- (5.0,4.8) -- (5.8,4.0) -- (6.5,2.8) -- (6.8,1.8);
  %% BCI-IV (lower)
  \draw[teal!80!black, thick] (0,6.0) -- (1.0,5.6) -- (2.5,5.0) -- (4.0,4.2) -- (5.0,3.2) -- (5.8,2.4) -- (6.5,1.5) -- (6.8,1.0);
  %% Legend
  \node[font=\small, blue!70!black] at (2.0,2.0) {CHB-MIT};
  \node[font=\small, red!70!black] at (2.0,1.4) {SAM40};
  \node[font=\small, teal!80!black] at (2.0,0.8) {BCI-IV};
  \draw[blue!70!black, thick] (0.8,2.0) -- (1.4,2.0);
  \draw[red!70!black, thick] (0.8,1.4) -- (1.4,1.4);
  \draw[teal!80!black, thick] (0.8,0.8) -- (1.4,0.8);
  %% Title
  \node[font=\small\bfseries] at (3.4,7.5) {(b) Precision-Recall Curves};
\end{scope}
\end{tikzpicture}%
}
\caption[ROC and Precision-Recall Curves for EEG Disease Finder]{(a) ROC Curves with AUC Values: CHB-MIT (0.974) $>$ SAM40 (0.958) $>$ BCI-IV (0.872). (b) Precision-Recall Curves Showing Average Precision Scores for Three Representative Datasets}
\label{app_fig:eeg_roc_pr}
\vspace{0.5em}\noindent\textit{Note.} ROC = receiver operating characteristic; AUC = area under the curve; FPR = false positive rate; TPR = true positive rate. CHB-MIT = Children's Hospital Boston--MIT (epilepsy); SAM40 = SAM40 stress dataset; BCI-IV = BCI Competition IV (motor imagery). Curves are computed from ensemble soft-voting predictions under leave-one-subject-out cross-validation. \textit{Source:}~\cite{asthana2025eegmaster}.
\end{figure}

Table~\ref{tab:eeg_loso_performance} summarises the per-dataset LOSO classification performance of the Agentic AI Disease Finder ensemble across all 10 evaluation datasets.
\needspace{8\baselineskip}
\noindent Table~\ref{app_tab:eeg_loso_performance} loso classification performance across 10 eeg datasets.

{\tablestripes\tablefontsize
\begin{xltabular}{\textwidth}{@{} l c c c c L @{}}
\caption[LOSO Classification Performance Across 10 EEG Datasets]{Classification Performance (Mean $\pm$ Std Across LOSO Folds) for the Agentic AI Disease Finder Ensemble}\label{app_tab:eeg_loso_performance} \\
\toprule
\thead{Dataset} & \thead{Acc (\%)} & \thead{F1} & \thead{AUC} & \thead{$\kappa$} & \thead{MCC} \\
\midrule
\endfirsthead
\multicolumn{6}{@{}l}{\textit{(\,Tab.~\ref{app_tab:eeg_loso_performance} continued from previous page\,)}} \\
\toprule
\thead{Dataset} & \thead{Acc (\%)} & \thead{F1} & \thead{AUC} & \thead{$\kappa$} & \thead{MCC} \\
\midrule
\endhead
\bottomrule
\endlastfoot
SAM40        & 91.5 $\pm$ 3.2 & .903 $\pm$ .04 & .958 $\pm$ .02 & .831 & .832 \\
DASPS        & 86.2 $\pm$ 4.8 & .849 $\pm$ .06 & .931 $\pm$ .03 & .792 & .794 \\
MODA         & 83.1 $\pm$ 5.1 & .821 $\pm$ .06 & .912 $\pm$ .04 & .761 & .763 \\
CHB-MIT      & 92.8 $\pm$ 2.9 & .921 $\pm$ .03 & .974 $\pm$ .02 & .856 & .857 \\
TUH-EEG      & 79.3 $\pm$ 6.2 & .782 $\pm$ .07 & .895 $\pm$ .04 & .712 & .715 \\
BCI-IV       & 76.2 $\pm$ 7.8 & .748 $\pm$ .08 & .872 $\pm$ .05 & .683 & .686 \\
PhysioNet-MI & 82.8 $\pm$ 4.5 & .819 $\pm$ .05 & .918 $\pm$ .03 & .771 & .773 \\
UNM-PD       & 89.2 $\pm$ 3.8 & .885 $\pm$ .04 & .951 $\pm$ .02 & .784 & .786 \\
Sleep-EDF    & 85.4 $\pm$ 4.2 & .842 $\pm$ .05 & .932 $\pm$ .03 & .805 & .807 \\
SHHS         & 86.7 $\pm$ 3.9 & .858 $\pm$ .04 & .941 $\pm$ .02 & .817 & .819 \\
\addlinespace
\textbf{Mean} & \textbf{87.3 $\pm$ 4.8} & \textbf{.861} & \textbf{.934} & \textbf{.781} & \textbf{.783} \\
\end{xltabular}}
\vspace{0.5em}\noindent\textit{Note.} Acc = accuracy; F1 = macro-averaged F1-score; AUC = area under the ROC curve (one-vs-rest); $\kappa$ = Cohen's kappa; MCC = Matthews correlation coefficient. All metrics computed under leave-one-subject-out (LOSO) cross-validation. Standard deviations reflect inter-fold variability. \textit{Source:}~\cite{asthana2025eegmaster}.

\subsection{Training Configuration}

%% ENGINEERING DETAIL SUPPRESSED — hyperparameters, optimiser settings, regularisation
\iffalse
\subsubsection{VotingClassifier Ensemble (Deployed Model)}
VotingClassifier: single-pass fit, soft voting, balanced weights, Gaussian noise augmentation, 5-fold CV seed 42.
\subsubsection{Deep Learning Models (Compared)}
Adam optimiser (lr $10^{-3}$, $\beta_1=0.9$, $\beta_2=0.999$), batch 16--64, early stopping patience 10--20, dropout 0.25--0.5, L2 decay $10^{-4}$.
\fi

Training followed two protocols, summarised in Table~\ref{app_tab:training_protocols}. All hyperparameters were selected via grid search within inner cross-validation folds. Full training specifications are documented in the Supplementary Volume.
\needspace{8\baselineskip}
{\tablestripes\tablefontsize
\begin{xltabular}{\textwidth}{@{} p{2.5cm} L L @{}}
\caption{Training Protocols: Deployed VotingClassifier vs Compared DL Models}\label{app_tab:training_protocols} \\
\toprule
\thead{Aspect} & \thead{VotingClassifier (Deployed)} & \thead{DL Models (Compared)} \\
\midrule
\endfirsthead
\multicolumn{3}{@{}l}{\textit{(\,Tab.~\ref{app_tab:training_protocols} continued from previous page\,)}} \\
\toprule
\thead{Aspect} & \thead{VotingClassifier (Deployed)} & \thead{DL Models (Compared)} \\
\midrule
\endhead
\bottomrule
\endlastfoot
Training method & Single-pass fitting (no gradient training) & Iterative gradient descent (Adam optimiser) \\
\addlinespace
Validation & 5-fold stratified CV (seed 42) & Domain-adaptive CV (5-fold or LOSO) \\
\addlinespace
Class handling & Balanced class weights & Class-weighted loss function \\
\addlinespace
Regularisation & Built-in (tree bagging) & Dropout 0.25--0.5; L2 decay $10^{-4}$ \\
\addlinespace
Convergence & N/A (single pass) & Early stopping (patience 10--20 epochs) \\
\addlinespace
Batch size & N/A & 16--64 (domain-adaptive) \\
\addlinespace
Learning rate & N/A & $10^{-3}$ ($\beta_1$=0.9, $\beta_2$=0.999) \\
\end{xltabular}}
\vspace{0.5em}\noindent\textit{Note.} CV = cross-validation; LOSO = leave-one-subject-out; DL = deep learning. The VotingClassifier requires no GPU and trains in under 30\,s per fold; DL models require GPU and train for 50--200 epochs per domain.

%% ENGINEERING DETAIL SUPPRESSED — seed settings, learning rate scheduling
\iffalse
All experiments use a fixed random seed (42) across PyTorch, TensorFlow, NumPy, and Scikit-learn. Learning rate scheduling: reduce-on-plateau (factor=0.5, patience=5).
\fi

Before training commences, each model architecture undergoes readiness assessment across seven criteria. Table~\ref{tab:model_readiness} documents the pass/fail status, enforcing an all-or-nothing deployment gate: a model that fails any single criterion does not proceed to deployment.
\needspace{8\baselineskip}
\noindent Table~\ref{app_tab:model_readiness} model readiness assessment: seven pre-deployment quality gates.

{\tablestripes\tablefontsize
\begin{xltabular}{\textwidth}{@{} p{2.2cm} L p{2.5cm} p{1.2cm} L @{}}
\caption{Model Readiness Assessment: Seven Pre-Deployment Quality Gates}\label{app_tab:model_readiness} \\
\toprule
\thead{Criterion} & \thead{What Is Verified} & \thead{Pass Threshold} & \thead{Status} & \thead{If Failed} \\
\midrule
\endfirsthead
\multicolumn{5}{@{}l}{\textit{(\,Tab.~\ref{app_tab:model_readiness} continued from previous page\,)}} \\
\toprule
\thead{Criterion} & \thead{What Is Verified} & \thead{Pass Threshold} & \thead{Status} & \thead{If Failed} \\
\midrule
\endhead
\bottomrule
\endlastfoot
Accuracy & Per-disease sensitivity on held-out test set & $\geq$85\% & Pass & Retrain; tune hyperparameters \\
\addlinespace
Stability & Standard deviation across 10-fold CV & SD $\leq$3\% & Pass & Investigate fold variance \\
\addlinespace
Fairness & Accuracy parity across demographic subgroups & DI $\geq$0.80 & Pass & Re-balance training data \\
\addlinespace
Explainability & SHAP--LIME--Grad-CAM agreement on top features & $\geq$0.70 concordance & Pass & Review feature pipeline \\
\addlinespace
Calibration & Predicted probability matches observed frequency & Brier $\leq$0.15 & Partial & Platt scaling \\
\addlinespace
Robustness & Performance on noisy/adversarial inputs & $<$5\% accuracy drop & Pass & Augment training data \\
\addlinespace
Governance & 525-module audit pass rate & $\geq$95\% & Pass & Remediate failed modules \\
\end{xltabular}}
\vspace{0.5em}\noindent\textit{Note.} Model readiness is all-or-nothing: partial readiness is not readiness. The Partial status for calibration (Brier score) reflects the known challenge of probability calibration in multi-class EEG classification, addressed through post-hoc Platt scaling. DI = disparate impact ratio; CV = cross-validation; SHAP = SHapley Additive exPlanations.

%% ============================================================================
%% CROSS-DATASET TRANSFER LEARNING (from EEG-RAG paper)
%% ============================================================================
\subsection{Cross-Dataset Transfer Learning Protocol}
\label{app_sec:cross_dataset_transfer}

Cross-dataset generalisation validates deployment readiness: a framework that fails on data from different recording sites cannot be deployed across hospitals. Table~\ref{tab:cross_dataset_transfer} documents the cross-dataset transfer protocol design; numerical results are reported in Chapter~\ref{chap:analysis}.
\needspace{3\baselineskip}
\noindent Table~\ref{app_tab:cross_dataset_transfer} cross-dataset transfer learning.

{\tablestripes\tablefontsize
\begin{xltabular}{\textwidth}{@{} p{3.5cm} L L @{}}
\caption[Cross-Dataset Transfer Learning]{Cross-Dataset Transfer Learning Protocol Design}\label{app_tab:cross_dataset_transfer} \\
\toprule
\thead{Source $\to$ Target} & \thead{DA Method} & \thead{Evaluation} \\
\midrule
\endfirsthead
\multicolumn{3}{@{}l}{\textit{(\,Tab.~\ref{app_tab:cross_dataset_transfer} continued from previous page\,)}} \\
\toprule
\thead{Source $\to$ Target} & \thead{DA Method} & \thead{Evaluation} \\
\midrule
\endhead
\bottomrule
\endlastfoot
DEAP $\to$ SAM-40 & Feature-level alignment with MMD loss; z-score re-normalisation & Accuracy with and without DA compared \\
\addlinespace
SAM-40 $\to$ DEAP & Feature-level alignment with MMD loss; z-score re-normalisation & Accuracy with and without DA compared \\
\addlinespace
DEAP $\to$ Kaggle & Direct transfer (same modality); no DA applied & Baseline transfer accuracy \\
\end{xltabular}}
\vspace{0.5em}\noindent\textit{Note.} DA = domain adaptation; MMD = maximum mean discrepancy. Numerical results are reported in Chapter~\ref{chap:analysis}.

%% ============================================================================
%% ABLATION STUDY DESIGN (from EEG-RAG + NeuroMCP papers)
%% ============================================================================
\subsection{Ablation Study Design}
\label{app_sec:ablation_design}

Ablation studies systematically remove or disable individual architectural components to quantify each component's contribution to overall classification accuracy. This design follows the principle that a claim about a component's value is only credible when supported by evidence of performance degradation upon its removal. Table~\ref{tab:ablation_design} documents the five ablation experiments and their expected contribution. Numerical results are reported in Chapter~\ref{chap:analysis}.
\needspace{3\baselineskip}
\noindent Table~\ref{app_tab:ablation_design} ablation study design.

{\tablestripes\tablefontsize
\begin{xltabular}{\textwidth}{@{} >{\raggedright\arraybackslash}p{0.16\textwidth} >{\raggedright\arraybackslash}p{0.14\textwidth} >{\raggedright\arraybackslash}X >{\raggedright\arraybackslash}p{0.10\textwidth} @{}}
\caption[Ablation Study Design]{Ablation Study Design: Component Removal Impact on Classification Accuracy}\label{app_tab:ablation_design} \\
\toprule
\thead{Component Removed} & \thead{Expected Role} & \thead{Rationale} & \thead{Results} \\
\midrule
\endfirsthead
\multicolumn{4}{@{}l}{\textit{(\,Tab.~\ref{app_tab:ablation_design} continued from previous page\,)}} \\
\toprule
\thead{Component Removed} & \thead{Expected Role} & \thead{Rationale} & \thead{Results} \\
\midrule
\endhead
\bottomrule
\endlastfoot
Bi-LSTM layer & Temporal sequence modelling & Captures EEG dynamics that static features miss & Ch.~\ref{chap:analysis} \\
\addlinespace
Text encoder & Cross-subject generalisation & Provides demographic and clinical context & Ch.~\ref{chap:analysis} \\
\addlinespace
Attention mechanism & Temporal segment weighting & Weights clinically relevant temporal segments & Ch.~\ref{chap:analysis} \\
\addlinespace
Consistency loss & Feature path agreement & Enforces agreement between parallel feature extraction paths & Ch.~\ref{chap:analysis} \\
\addlinespace
CNN only (no ensemble) & Baseline isolation & Isolates CNN from ensemble to measure multi-model combination value & Ch.~\ref{chap:analysis} \\
\end{xltabular}}
\vspace{0.5em}\noindent\textit{Note.} Bi-LSTM = bidirectional long short-term memory; CNN = convolutional neural network. All ablation experiments use the same test set and cross-validation protocol (5-fold stratified CV, seed 42) to ensure fair comparison. Numerical accuracy impact results are reported in Chapter~\ref{chap:analysis}.


% [SPACE-FIX] clearpage removed to eliminate empty page
\subsection{RAG/GenAI Technical Specifications}
\label{app:sec7_rag_genai}


%% ENGINEERING DETAIL SUPPRESSED — MMR equation, pipeline implementation details
\iffalse
MMR re-ranking: $\text{MMR} = \arg\max_{d_i \in R \setminus S} [\lambda \cdot \text{sim}(d_i,q) - (1-\lambda) \cdot \max_{d_j \in S} \text{sim}(d_i,d_j)]$, $\lambda=0.7$.
\fi

The Retrieval-Augmented Generation (RAG) pipeline addresses a governance requirement specified by the EU AI Act (Art.~14) and FDA SaMD guidance: every AI prediction must be accompanied by an evidence-grounded explanation that clinicians can independently verify~\cite{lewis2020rag, gao2024rag}. Unlike post-hoc attribution methods (SHAP, LIME) that explain \textit{which features} drove a prediction, RAG explains \textit{why} the prediction is clinically meaningful by linking it to peer-reviewed biomedical evidence.

Table~\ref{tab:xai_method_comparison} compares the four explainability methods employed in RGAIG. Each method addresses a distinct clinical question, and their triangulation---measuring cross-method agreement on the same prediction---provides convergent evidence that no single technique alone can offer.
\needspace{8\baselineskip}
\noindent Table~\ref{app_tab:xai_method_comparison} explainability method comparison.

{\tablestripes\tablefontsize
\begin{xltabular}{\textwidth}{@{} p{2cm} L L p{2.5cm} p{2.5cm} @{}}
\caption[Explainability Method Comparison]{Explainability Method Comparison: Scope, Clinical Question, Output, and Disease Applicability}\label{app_tab:xai_method_comparison} \\
\toprule
\thead{XAI Method} & \thead{What It Explains} & \thead{Clinical Question Answered} & \thead{Output Format} & \thead{Diseases Applied} \\
\midrule
\endfirsthead
\multicolumn{5}{@{}l}{\textit{(\,Tab.~\ref{app_tab:xai_method_comparison} continued from previous page\,)}} \\
\toprule
\thead{XAI Method} & \thead{What It Explains} & \thead{Clinical Question Answered} & \thead{Output Format} & \thead{Diseases Applied} \\
\midrule
\endhead
\bottomrule
\endlastfoot
SHAP & Feature-level attribution using Shapley values from cooperative game theory & Which EEG channels and frequency bands drove this specific prediction? & Beeswarm plot; force plot; feature ranking & All the ASD domain (ASD, PD, Epilepsy, Stress, Cancer) \\
\addlinespace
LIME & Local instance-level perturbation of the input space & Why was this specific patient flagged as positive? & Weighted feature list with local fidelity score & ASD, PD, Epilepsy \\
\addlinespace
Grad-CAM & Spatial activation heatmap from gradient-weighted class activations in CNN layers & Which brain regions activated the convolutional network's decision? & Overlay heatmap on EEG spectrogram & Epilepsy, Schizophrenia \\
\addlinespace
RAG Narrative & Natural-language rationale grounded in retrieved peer-reviewed literature & What does this result mean in clinical terms, supported by published evidence? & Text paragraph with PubMed citations and confidence score & All the ASD domain \\
\addlinespace
Triangulation & Cross-method agreement score measuring convergence across SHAP, LIME, and Grad-CAM & Do independent explanation methods agree on the same predictive features? & Agreement score ($\geq$0.70 required) & All the ASD domain \\
\end{xltabular}}
\vspace{0.5em}\noindent\textit{Note.} SHAP = SHapley Additive exPlanations~\cite{lundberg2017shap}; LIME = Local Interpretable Model-Agnostic Explanations~\cite{ribeiro2016lime}; Grad-CAM = Gradient-weighted Class Activation Mapping~\cite{selvaraju2017gradcam}; RAG = Retrieval-Augmented Generation~\cite{lewis2020rag}. Triangulation threshold of $\geq$0.70 follows the convergent validity standard. Methods not applicable to a domain (e.g., Grad-CAM requires CNN spatial layers) are excluded from triangulation for that domain.

\noindent Figure~\ref{fig:rag_stages} describes the five pipeline stages.
\needspace{20\baselineskip}
\begin{figure}[H]
\centering
\resizebox{0.97\textwidth}{!}{%
\begin{tikzpicture}[
  stage/.style={draw=ggunavy, fill=ggunavy!8, rounded corners=3pt, text width=2.7cm, minimum height=2.6cm, font=\small, align=center, inner sep=4pt},
  arr/.style={-{Stealth[length=4mm, width=3mm]}, very thick, ggunavy},
  node distance=0.8cm
]
\node[stage, fill=lightblue] (s1) {\textbf{\small 1. Knowledge}\\[2pt]\textbf{\small Base}\\[4pt]50{,}000+ PubMed\\abstracts + guidelines\\512-token window\\384-dim vectors\\{\small ASD, 256 stride}};
\node[stage, fill=lightorange, right=of s1] (s2) {\textbf{\small 2. Vector}\\[2pt]\textbf{\small Indexing}\\[4pt]FAISS IVF-PQ\\256 Voronoi clusters\\384-dim $\to$ 48 bytes\\recall $>$95\%\\{\small Built once, incremental}};
\node[stage, fill=lightteal, right=of s2] (s3) {\textbf{\small 3. Retrieval}\\[4pt]Query = label +\\top-5 SHAP + domain\\cosine similarity ($K$=5)\\MMR re-ranking\\($\lambda$=0.7)\\{\small Sub-linear search}};
\node[stage, fill=lightpurple, right=of s3] (s4) {\textbf{\small 4. Generation}\\[4pt]Structured prompt\\$\to$ LLM narrative:\\diagnosis, features,\\citations, caveats\\{\small Temp 0.3, max 512 tok}};
\node[stage, fill=lightgreen, right=of s4] (s5) {\textbf{\small 5. Post-}\\[2pt]\textbf{\small Processing}\\[4pt]Hallucination detect,\\citation accuracy\\(cosine $\geq$0.65),\\confidence score\\{\small Automated QA}};
\draw[arr] (s1) -- (s2);
\draw[arr] (s2) -- (s3);
\draw[arr] (s3) -- (s4);
\draw[arr] (s4) -- (s5);
\end{tikzpicture}%
}% end resizebox
\caption{RAG Pipeline: Five Sequential Stages}
\label{app_fig:rag_stages}
\vspace{0.5em}\noindent\textit{Note.} FAISS = Facebook AI Similarity Search; IVF-PQ = inverted file with product quantisation; MMR = maximal marginal relevance; MDS = Movement Disorder Society.
\end{figure}

Explainability quality is assessed via expert-rated coherence, citation relevance, and clinical alignment scores, as described in the expert panel validation protocol (Section~\ref{sec:validity}).

%% SUPPRESSED — redundant with fig:rag_stages (TikZ version preferred over PNG)
\iffalse
While the preceding five-stage diagram (Figure~\ref{fig:rag_stages}) shows the logical stages, Figure~\ref{fig:rag_pipeline_visual} provides an implementation-level view of the data flow through the RAG system.
\noindent The RAG pipeline visual is presented in the main text as Figure~\ref{fig:rag_pipeline_visual} (Chapter~\\ref{chap:methods}); not reproduced here to avoid duplication.

\noindent The visual shows the concrete data flow: biomedical abstracts are embedded into 384-dimensional vectors, stored in FAISS, and retrieved via cosine similarity for each prediction; the retrieved evidence is then assembled into a structured prompt for the LLM, which generates a citation-grounded clinical narrative. This implementation-level detail confirms that every generated explanation is traceable to specific peer-reviewed sources, directly satisfying the EU AI Act Article~14 transparency requirement and operationalising H5 (RAG explainability $\geq$80\% expert agreement). Expert panel evaluation scores for these narratives are reported in Chapter~\ref{chap:analysis}.
\fi

Figure~\ref{fig:phase6_rag_explainability} provides a comprehensive two-column view of the Phase~6 RAG-powered explainability pipeline, detailing the nine-step process from model prediction through SHAP extraction, knowledge base preparation, vector storage, and LLM-based report generation, alongside the rationale for each design decision and its alignment with multi-stakeholder needs.

\needspace{4\baselineskip}
\input{figures/fig_phase6_rag_explainability}

%% ============================================================================
%% A1: COMPLETE AGENTIC RAG ARCHITECTURE (EXPANDED)
%% ============================================================================
\subsection{Agentic RAG: Nine-Stage Pipeline Architecture}
\label{app_sec:agentic_rag_architecture}

%% ENGINEERING DETAIL SUPPRESSED — 9-stage pipeline implementation, JSON messaging, GPU latency
\iffalse
Nine-stage agentic pipeline: pre-retrieval intelligence, output evaluation, guardrail enforcement. Each stage autonomous via JSON messaging. 2.3s median latency, 1,200 queries/day on RTX 4090.
\fi

The production implementation extends the five core RAG stages into a nine-stage \textit{agentic} pipeline (Figure~\ref{fig:rag_9stage}) that adds pre-retrieval intelligence, output evaluation, and guardrail enforcement---capabilities addressing governance requirements for clinical AI systems (EU AI Act, NIST AI RMF) where every generated explanation must pass safety checks before delivery to a clinician. The ``agentic'' design means each stage can autonomously rewrite queries, select retrieval strategies, and escalate uncertain outputs without human intervention, addressing a governance requirement for self-monitoring AI systems.
\needspace{20\baselineskip}
\begin{figure}[!htbp]
\centering
\resizebox{0.97\textwidth}{!}{%
\begin{tikzpicture}[
  node distance=0.8cm and 0.12cm,
  stage/.style={rectangle, draw=ggunavy, fill=ggunavy!8, rounded corners=3pt,
    minimum height=1.1cm, minimum width=2.3cm, align=center, font=\small},
  arrow/.style={-{Stealth[length=5pt]}, thick, ggunavy}
]
% Row 1: Stages 1-5
\node[stage, fill=lightblue] (s1) {{\small\bfseries 1. Query Input}\\{\small Intent classification,}\\{\small domain routing}};
\node[stage, fill=lightorange, right=of s1] (s2) {{\small\bfseries 2. Pre-Retrieval}\\{\small Query rewriting,}\\{\small PII scrubbing}};
\node[stage, fill=lightteal, right=of s2] (s3) {{\small\bfseries 3. Embedding}\\{\small Sentence-BERT,}\\{\small 384-dim vectors}};
\node[stage, fill=lightpurple, right=of s3] (s4) {{\small\bfseries 4. Vector Search}\\{\small FAISS + ChromaDB,}\\{\small top-$K$=5}};
\node[stage, fill=lightgreen, right=of s4] (s5) {{\small\bfseries 5. Post-Retrieval}\\{\small MMR re-ranking,}\\{\small context compression}};
% Row 2: Stages 6-9
\node[stage, fill=lightyellow, below=0.9cm of s1, xshift=1.3cm] (s6) {{\small\bfseries 6. Context Assembly}\\{\small Retrieved docs +}\\{\small SHAP + history}};
\node[stage, fill=lightpink, right=of s6] (s7) {{\small\bfseries 7. LLM Generation}\\{\small Ollama llama3.2:3b,}\\{\small temp=0.3}};
\node[stage, fill=lightcoral, right=of s7] (s8) {{\small\bfseries 8. Output Eval}\\{\small Faithfulness,}\\{\small hallucination check}};
\node[stage, fill=lightsky, right=of s8] (s9) {{\small\bfseries 9. Guardrail Gate}\\{\small 6 safety checks,}\\{\small approve/escalate}};
% Arrows
\draw[arrow] (s1) -- (s2);
\draw[arrow] (s2) -- (s3);
\draw[arrow] (s3) -- (s4);
\draw[arrow] (s4) -- (s5);
\draw[arrow] (s5.south) -- ++(0,-0.2) -| (s6.north);
\draw[arrow] (s6) -- (s7);
\draw[arrow] (s7) -- (s8);
\draw[arrow] (s8) -- (s9);
% Annotations
\node[below=0.08cm of s4, font=\small\itshape, text=gray] {cosine $\geq$0.7};
\node[below=0.08cm of s9, font=\small\itshape, text=gray] {pass or escalate};
\end{tikzpicture}%
}
\caption[Agentic RAG Pipeline (9 Stages)]{Nine-Stage Agentic RAG Pipeline: From Query Input to Guardrail-Approved Output}
\label{app_fig:rag_9stage}
\vspace{0.5em}\noindent\textit{Note.} FAISS = Facebook AI Similarity Search; MMR = maximal marginal relevance; SHAP = SHapley Additive exPlanations; PII = personally identifiable information; LLM = large language model. Stages 1--5 handle retrieval; Stages 6--9 handle generation and safety verification. The pipeline uses Ollama llama3.2:3b (3 billion parameters, open-source, locally deployed) for HIPAA-compliant on-premise inference.
\end{figure}

\noindent\textit{RAG figures summary.} Figures~\ref{fig:rag_stages} and~\ref{fig:rag_9stage} present two progressively detailed views of the RAG explainability system: the five-stage logical pipeline (Figure~\ref{fig:rag_stages}) and the nine-stage agentic extension with guardrail enforcement (Figure~\ref{fig:rag_9stage}). Collectively, these figures demonstrate that RAG explainability is not a single-step process but a multi-stage architecture with embedded safety verification, pre-retrieval intelligence, and hallucination detection---directly supporting H5 (evidence-grounded explainability) and the EU AI Act transparency requirements (O5).

%% MOVED TO SUPPLEMENTARY VOLUME — RAG database architecture table (tab:rag_databases),
%% chunking strategy comparison table (tab:chunking_strategies), embedding model
%% comparison table (tab:embedding_models), and three-phase retrieval table (tab:retrieval_phases)
\paragraph{RAG Retrieval Architecture.}

RAG retrieval uses ChromaDB (1.17~GB) for persistence and FAISS for in-memory vector search over 50K+ biomedical abstracts. Documents are chunked recursively (512 characters, 50-character overlap) and embedded via all-MiniLM-L6-v2 (384 dimensions, 14,200 sentences/second on CPU). Retrieval operates in three phases: pre-retrieval (query rewriting, PII scrubbing, domain routing with $>$98\% accuracy), retrieval (cosine similarity $\geq$0.7, hybrid dense + BM25 search), and post-retrieval (MMR re-ranking with $\lambda$=0.7, context compression, citation verification $\geq$90\%). Architecture specifications---including database architecture, chunking strategy comparison, embedding model selection rationale, and phase-specific quality metrics---are in the Supplementary Volume.

Table~\ref{tab:rag_retrieval_specs} summarises the RAG retrieval pipeline specifications across its three operational phases.
\needspace{8\baselineskip}
\noindent Table~\ref{app_tab:rag_retrieval_specs} rag retrieval pipeline specifications.

{\tablestripes\tablefontsize
\begin{xltabular}{\textwidth}{@{} p{2cm} p{3cm} L p{2.5cm} @{}}
\caption{RAG Retrieval Pipeline Specifications}\label{app_tab:rag_retrieval_specs} \\
\toprule
\thead{Phase} & \thead{Component} & \thead{Parameter} & \thead{Value} \\
\midrule
\endfirsthead
\multicolumn{4}{@{}l}{\textit{(\,Tab.~\ref{app_tab:rag_retrieval_specs} continued from previous page\,)}} \\
\toprule
\thead{Phase} & \thead{Component} & \thead{Parameter} & \thead{Value} \\
\midrule
\endhead
\bottomrule
\endlastfoot
Storage & ChromaDB persistence & Vector database size & 1.17 GB \\
Storage & FAISS in-memory index & Biomedical abstracts indexed & 50,000+ \\
\addlinespace
Chunking & Recursive text splitter & Chunk size / overlap & 512 / 50 characters \\
Embedding & all-MiniLM-L6-v2 & Dimensions; throughput & 384-d; 14,200 sent/sec (CPU) \\
\addlinespace
Pre-retrieval & Query rewriting + PII scrub & Domain routing accuracy & $>$98\% \\
Retrieval & Hybrid dense + BM25 & Cosine similarity threshold & $\geq$0.70 \\
Post-retrieval & MMR re-ranking & Diversity parameter ($\lambda$) & 0.70 \\
Post-retrieval & Citation verification & Verification threshold & $\geq$90\% \\
\end{xltabular}}
\vspace{0.5em}\noindent\textit{Note.} MMR = maximal marginal relevance; PII = personally identifiable information; BM25 = Best Matching 25 (term-frequency retrieval). All thresholds are predefined acceptance criteria evaluated in Chapter~\ref{chap:analysis}.

%% MOVED TO SUPPLEMENTARY VOLUME — EEG-Stress RAG Integration section
%% (Tables tab:stress_feature_mapping, tab:stress_knowledge_base, tab:three_tier_explanation)
\subsection{EEG-Stress RAG Integration}
\label{app_sec:eeg_stress_rag}

The EEG-Stress RAG module demonstrates domain-specific integration, grounding classification explanations in neurophysiological biomarkers (alpha suppression, beta elevation, frontal asymmetry, theta elevation, gamma coherence) and 50K+ PubMed abstracts. Feature-to-explanation mappings, knowledge base specifications, and three-tier explanation architecture (classification statement, SHAP-based reasoning, evidence-grounded citations) are documented in the Supplementary Volume.

%% ============================================================================
%% A3: GENAI OUTPUT EVALUATION FRAMEWORK
%% ============================================================================
\subsection{GenAI Output Evaluation Framework}
\label{app_sec:genai_eval}

%% ENGINEERING DETAIL SUPPRESSED — framework comparison narrative
\iffalse
Three frameworks: DeepEval (14 metrics), RAGAS (8 metrics), G-Eval (LLM-based scoring). Detailed comparison in Supplementary Volume.
\fi

Evaluating RAG-generated clinical explanations requires multi-framework quality assurance covering faithfulness, retrieval quality, and narrative coherence. RGAIG combines metrics from three evaluation frameworks (the corresponding table): DeepEval for faithfulness and hallucination detection, RAGAS for retrieval quality, and G-Eval for clinical narrative coherence. This multi-framework triangulation ensures that no single quality dimension is left unmonitored---a regulatory requirement for clinical AI systems where explanation quality directly affects clinician trust and patient safety.
\needspace{8\baselineskip}
\noindent Table~\ref{app_tab:genai_eval_frameworks} genai evaluation frameworks: deepeval, ragas, and g-eval.

{\tablestripes\tablefontsize
\begin{xltabular}{\textwidth}{@{} p{2.0cm} L L p{2.0cm} L @{}}
\caption{GenAI Evaluation Frameworks: DeepEval, RAGAS, and G-Eval}\label{app_tab:genai_eval_frameworks} \\
\toprule
\thead{Framework} & \thead{Metrics} & \thead{Evaluation Method} & \thead{Open Source} & \thead{RGAIG Usage} \\
\midrule
\endfirsthead
\multicolumn{5}{@{}l}{\textit{(\,Tab.~\ref{app_tab:genai_eval_frameworks} continued from previous page\,)}} \\
\toprule
\thead{Framework} & \thead{Metrics} & \thead{Evaluation Method} & \thead{Open Source} & \thead{RGAIG Usage} \\
\midrule
\endhead
\bottomrule
\endlastfoot
DeepEval & 14 metrics: faithfulness, hallucination, bias, toxicity, answer relevancy & Reference-free and reference-based; custom metric support & Yes (Python) & Primary framework for faithfulness and hallucination scoring \\
\addlinespace
RAGAS & 8 metrics: context precision, context recall, faithfulness, answer relevancy, answer correctness & Component-wise evaluation of retrieval and generation stages independently & Yes (Python) & Retrieval quality: context precision and recall measurement \\
\addlinespace
G-Eval & 4 criteria: coherence, consistency, fluency, relevance & LLM-based evaluation with chain-of-thought scoring (GPT-4 as judge) & Partial & Coherence and fluency scoring for clinical narrative quality \\
\end{xltabular}}
\vspace{0.5em}\noindent\textit{Note.} DeepEval serves as the primary evaluation harness; RAGAS provides retrieval-specific metrics; G-Eval provides LLM-judged narrative quality assessment. Using multiple frameworks provides triangulation of output quality.

\paragraph{RGAIG Output Quality Metrics.}

Table~\ref{app_tab:genai_thresholds} documents the predefined acceptance thresholds applied during RGAIG validation testing ($n$=500 generated explanations across five clinical domains). Numerical results are reported in Chapter~\ref{chap:analysis}.
\needspace{8\baselineskip}
{\tablestripes\tablefontsize
\begin{xltabular}{\textwidth}{@{} p{3cm} c p{2cm} L @{}}
\caption{GenAI Output Quality: Predefined Acceptance Thresholds}\label{app_tab:genai_thresholds} \\
\toprule
\thead{Metric} & \thead{Threshold} & \thead{Framework} & \thead{What It Measures} \\
\midrule
\endfirsthead
\multicolumn{4}{@{}l}{\textit{(\,Tab.~\ref{app_tab:genai_thresholds} continued from previous page\,)}} \\
\toprule
\thead{Metric} & \thead{Threshold} & \thead{Framework} & \thead{What It Measures} \\
\midrule
\endhead
\bottomrule
\endlastfoot
Faithfulness & $\geq$0.85 & DeepEval & Consistency between generated text and retrieved source evidence \\
\addlinespace
Answer relevancy & $\geq$0.80 & RAGAS & Relevance of generated answer to the clinical query \\
\addlinespace
Context recall & $\geq$0.80 & RAGAS & Proportion of relevant documents retrieved from knowledge base \\
\addlinespace
Hallucination rate & $\leq$5\% & DeepEval & Unsupported claims in generated clinical narrative \\
\end{xltabular}}
\vspace{0.5em}\noindent\textit{Note.} $n$ = 500 generated explanations across five clinical domains. DeepEval and RAGAS are open-source evaluation frameworks; thresholds are predefined acceptance criteria.


%% MOVED TO SUPPLEMENTARY VOLUME — RAG Governance 15 Assurance Layers figure
%% (fig:rag_governance_layers)
\paragraph{RAG Governance: Fifteen Assurance Layers.}
\label{app_fig:rag_governance_layers}

The RAG pipeline implements 15 automated assurance layers (A1--A15) organised into seven governance categories. Nine layers execute synchronously per query; six operate asynchronously on scheduled monitoring. Table~\ref{app_tab:rag_assurance_summary} summarises these categories with their execution mode and acceptance thresholds. Detailed layer specifications are in the Supplementary Volume.
\needspace{8\baselineskip}
{\tablestripes\tablefontsize
\begin{xltabular}{\textwidth}{@{} p{2.8cm} p{1.5cm} c L @{}}
\caption{RAG Governance: Fifteen Assurance Layers by Category}\label{app_tab:rag_assurance_summary} \\
\toprule
\thead{Category} & \thead{Layers} & \thead{Execution} & \thead{Key Threshold / Tool} \\
\midrule
\endfirsthead
\multicolumn{4}{@{}l}{\textit{(\,Tab.~\ref{app_tab:rag_assurance_summary} continued from previous page\,)}} \\
\toprule
\thead{Category} & \thead{Layers} & \thead{Execution} & \thead{Key Threshold / Tool} \\
\midrule
\endhead
\bottomrule
\endlastfoot
Faithfulness validation & A1--A2 & Sync & DeepEval $\geq$0.85; RAGAS context recall $\geq$0.80 \\
\addlinespace
Input sanitisation & A3--A4 & Sync & Prompt injection blocking; SQL injection detection \\
\addlinespace
Access control & A5--A6 & Sync & RBAC policy engine; API key verification \\
\addlinespace
Contradiction detection & A7--A8 & Sync & NLI-based; entailment score $\geq$0.70 \\
\addlinespace
Hallucination monitoring & A9--A11 & Sync & Citation accuracy $\geq$90\%; atomic claim verification \\
\addlinespace
Lineage versioning & A12--A13 & Async & SHA-256 hash-based immutable registry \\
\addlinespace
Regulatory compliance & A14--A15 & Async & HIPAA \S164.312; GDPR Art.~30; scheduled audit \\
\end{xltabular}}
\vspace{0.5em}\noindent\textit{Note.} Sync = synchronous (evaluated per query); Async = asynchronous (scheduled monitoring). RBAC = role-based access control; NLI = natural language inference; RAGAS = Retrieval-Augmented Generation Assessment. Layer identifiers A1--A15 correspond to the full specifications in the Supplementary Volume.

%% ---------------------------------------------------------------------------
%% THREE ARCHITECTURE DIAGRAMS from Agentic AI Disease Finder Master Analysis
%% Source: eeg_complete_master.pdf (Asthana, 2025)
%% ---------------------------------------------------------------------------

the corresponding figure presents the four-layer system architecture of the Agentic AI Disease Finder following the C4 model~\cite{asthana2025eegmaster}. The architecture separates concerns across Presentation, Application, Domain (ML/AI), and Infrastructure layers, enabling independent scaling of compute-intensive model inference without affecting the user-facing API gateway.
\needspace{20\baselineskip}
\begin{figure}[!htbp]
\centering
\resizebox{0.97\textwidth}{!}{%
\begin{tikzpicture}[
  layer/.style={draw=ggunavy, thick, dashed, rounded corners=4pt, inner sep=10pt},
  component/.style={draw=ggunavy, fill=#1, rounded corners=2pt,
    font=\small\sffamily, minimum height=1.0cm, minimum width=2.2cm,
    text width=2cm, align=center},
  layerlabel/.style={font=\small\bfseries\sffamily\color{ggunavy}, anchor=west},
]

% --- Infrastructure Layer (bottom) ---
\node[layer, fill=lightsalmon!30, minimum width=17.5cm, minimum height=1.6cm]
  (infra) at (0,0) {};
\node[layerlabel] at ([xshift=5pt]infra.north west) {Infrastructure Layer};

\node[component=lightsalmon!60] (pg) at (-6.5,0) {PostgreSQL /\\TimescaleDB};
\node[component=lightsalmon!60] (redis) at (-3.5,0) {Redis\\Cache};
\node[component=lightsalmon!60] (mlflow) at (-0.5,0) {MLflow\\Registry};
\node[component=lightsalmon!60] (s3) at (2.5,0) {S3\\Storage};
\node[component=lightsalmon!60] (k8s) at (5.0,0) {Kubernetes\\Cluster};
\node[component=lightsalmon!60] (prom) at (7.5,0) {Prometheus /\\Grafana};

% --- Domain Layer (ML/AI) ---
\node[layer, fill=lightgreen!30, minimum width=17.5cm, minimum height=1.6cm]
  (domain) at (0,2.6) {};
\node[layerlabel] at ([xshift=5pt]domain.north west) {Domain Layer (ML/AI)};

\node[component=lightgreen!60] (sig) at (-6.5,2.6) {Signal\\Processor (MNE)};
\node[component=lightgreen!60] (feat) at (-4.0,2.6) {Feature\\Extractor (NumPy)};
\node[component=lightgreen!60] (cnn) at (-1.5,2.6) {CNN Model\\(PyTorch)};
\node[component=lightgreen!60] (trans) at (1.0,2.6) {Transformer\\Model (PyTorch)};
\node[component=lightgreen!60] (riem) at (3.5,2.6) {Riemannian\\Classifier};
\node[component=lightgreen!60] (ens) at (5.8,2.6) {Ensemble\\Voter};
\node[component=lightgreen!60] (shap) at (8.0,2.6) {Explainer\\(SHAP)};

% --- Application Layer ---
\node[layer, fill=lightblue!30, minimum width=17.5cm, minimum height=1.6cm]
  (app) at (0,5.2) {};
\node[layerlabel] at ([xshift=5pt]app.north west) {Application Layer};

\node[component=lightblue!60] (auth) at (-5.5,5.2) {Auth\\Service};
\node[component=lightblue!60] (sched) at (-2.8,5.2) {Job\\Scheduler};
\node[component=lightblue!60] (preproc) at (0,5.2) {Preprocessing\\Orchestrator};
\node[component=lightblue!60] (model) at (3.0,5.2) {Model\\Orchestrator};
\node[component=lightblue!60] (agg) at (6.0,5.2) {Result\\Aggregator};

% --- Presentation Layer (top) ---
\node[layer, fill=lightyellow!30, minimum width=17.5cm, minimum height=1.6cm]
  (pres) at (0,7.8) {};
\node[layerlabel] at ([xshift=5pt]pres.north west) {Presentation Layer};

\node[component=lightyellow!60] (web) at (-5.5,7.8) {Web Dashboard\\(React)};
\node[component=lightyellow!60] (mob) at (-2.8,7.8) {Mobile App\\(Flutter)};
\node[component=lightyellow!60] (cli) at (0,7.8) {CLI Tool\\(Python)};
\node[component=lightyellow!60] (rest) at (3.0,7.8) {REST API\\(FastAPI)};
\node[component=lightyellow!60] (ws) at (6.0,7.8) {WebSocket\\Gateway};

% --- Inter-layer arrows ---
\draw[-stealth, ggunavy!60, thick]
  ([yshift=-3pt]pres.south) -- ([yshift=3pt]app.north);
\draw[-stealth, ggunavy!60, thick]
  ([yshift=-3pt]app.south) -- ([yshift=3pt]domain.north);
\draw[-stealth, ggunavy!60, thick]
  ([yshift=-3pt]domain.south) -- ([yshift=3pt]infra.north);

\end{tikzpicture}%
}% end resizebox
\caption[C4 Four-Layer System Architecture (Appendix)]{C4 four-layer system architecture of the Agentic AI Disease Finder. Each layer encapsulates a distinct concern---presentation, orchestration, ML/AI domain logic, and infrastructure---enabling independent deployment and horizontal scaling.}
\label{app_fig:c4_system_architecture}

\vspace{0.5em}\noindent\textit{Note.} C4 = Context, Container, Component, Code architectural model. MNE = MNE-Python EEG processing library. SHAP = SHapley Additive exPlanations. MLflow = open-source ML lifecycle platform. TimescaleDB = time-series PostgreSQL extension.
\end{figure}

the corresponding figure illustrates the end-to-end ML pipeline architecture that processes raw EEG signals through six sequential stages---from signal acquisition to clinical explanation. Each stage has defined input/output contracts and quality gates, ensuring that failures are caught early and propagated to the monitoring stack.
\needspace{20\baselineskip}
\begin{figure}[!htbp]
\centering
\resizebox{0.97\textwidth}{!}{%
\begin{tikzpicture}[
  stage/.style={draw=ggunavy, thick, rounded corners=4pt, fill=#1,
    minimum width=2.4cm, minimum height=2.8cm, inner sep=4pt},
  stagetitle/.style={font=\small\bfseries\sffamily\color{ggunavy}},
  stageitem/.style={font=\small\sffamily, text width=2cm, align=center},
  arrow/.style={-stealth, ggunavy, very thick, >=latex, shorten >=2pt, shorten <=2pt},
]

% Stage 1: Input
\node[stage=lightyellow!50] (s1) at (0,0) {};
\node[stagetitle] at ([yshift=-4pt]s1.north) {1. Input};
\node[stageitem] at (s1.center) {Raw EEG\\256\,Hz, 22\,Ch};
\node[stageitem, font=\small\sffamily, text=gray] at ([yshift=4pt]s1.south) {MNE-Python\\EDF/BDF\\Montage};

% Stage 2: Preprocessing
\node[stage=lightblue!50] (s2) at (3.4,0) {};
\node[stagetitle] at ([yshift=-4pt]s2.north) {2. Preprocess};
\node[stageitem] at (s2.center) {Filtered\\Cleaned\\Re-referenced};
\node[stageitem, font=\small\sffamily, text=gray] at ([yshift=4pt]s2.south) {Notch/BPF\\Bandpass\\ICA};

% Stage 3: Feature Engineering
\node[stage=lightgreen!50] (s3) at (6.8,0) {};
\node[stagetitle] at ([yshift=-4pt]s3.north) {3. Feature Eng.};
\node[stageitem] at (s3.center) {Spectral\\Connectivity\\Bio-markers};
\node[stageitem, font=\small\sffamily, text=gray] at ([yshift=4pt]s3.south) {PSD\\CWT\\Connectivity};

% Stage 4: Modeling
\node[stage=lightpurple!50] (s4) at (10.2,0) {};
\node[stagetitle] at ([yshift=-4pt]s4.north) {4. Modeling};
\node[stageitem] at (s4.center) {CNN/Trans\\Riemannian\\Ensemble};
\node[stageitem, font=\small\sffamily, text=gray] at ([yshift=4pt]s4.south) {EEGNet\\TDN\\ViT};

% Stage 5: Inference
\node[stage=lightorange!50] (s5) at (13.6,0) {};
\node[stagetitle] at ([yshift=-4pt]s5.north) {5. Inference};
\node[stageitem] at (s5.center) {Prediction\\Probability\\Class};
\node[stageitem, font=\small\sffamily, text=gray] at ([yshift=4pt]s5.south) {Softmax\\Calibration\\Threshold};

% Stage 6: RAG Explainability
\node[stage=lightcoral!30] (s6) at (17.0,0) {};
\node[stagetitle] at ([yshift=-4pt]s6.north) {6. RAG};
\node[stageitem] at (s6.center) {FAISS Retrieve\\LLM Synthesise\\Evidence-based};
\node[stageitem, font=\small\sffamily, text=gray] at ([yshift=4pt]s6.south) {ChromaDB\\llama3.2:3b\\SHAP+Citations};

% Stage 7: Output
\node[stage=lightpink!50] (s7) at (20.4,0) {};
\node[stagetitle] at ([yshift=-4pt]s7.north) {7. Output};
\node[stageitem] at (s7.center) {Diagnosis\\Explanation\\Report};
\node[stageitem, font=\small\sffamily, text=gray] at ([yshift=4pt]s7.south) {B2B/B2C\\PDF Report\\Guardrail};

% Arrows between stages
\draw[arrow] (s1.east) -- (s2.west);
\draw[arrow] (s2.east) -- (s3.west);
\draw[arrow] (s3.east) -- (s4.west);
\draw[arrow] (s4.east) -- (s5.west);
\draw[arrow] (s5.east) -- (s6.west);
\draw[arrow] (s6.east) -- (s7.west);

\end{tikzpicture}%
}% end resizebox
\caption[ML Pipeline Architecture (7 Stages)]{End-to-End ML Pipeline Architecture with Seven Sequential Stages Including RAG Explainability. Each stage enforces input/output contracts and quality gates before passing data to the next stage.}
\label{app_fig:ml_pipeline_arch}

\vspace{0.5em}\noindent\textit{Note.} EEG = electroencephalography; Hz = hertz; Ch = channels; BPF = bandpass filter; ICA = independent component analysis; PSD = power spectral density; CWT = continuous wavelet transform; CNN = convolutional neural network; Trans = Transformer; TDN = temporal difference network; ViT = Vision Transformer; SHAP = SHapley Additive exPlanations; RAG = retrieval-augmented generation; FAISS = Facebook AI Similarity Search; LLM = large language model. Stage~6 (RAG) retrieves evidence from 50{,}000+ PubMed abstracts via ChromaDB/FAISS and synthesises clinical explanations using llama3.2:3b. B2B = business-to-business; B2C = business-to-consumer.
\end{figure}

the corresponding figure depicts the data flow through the processing pipeline, identifying six processing stages, three data stores, and two external entities. The flow ensures complete data lineage from raw EEG signal acquisition through preprocessing, feature extraction, model prediction, and RAG-based explanation to persistent storage and clinician delivery.
\needspace{20\baselineskip}
\begin{figure}[!htbp]
\centering
\resizebox{0.97\textwidth}{!}{%
\begin{tikzpicture}[
  process/.style={draw=ggunavy, thick, circle, fill=lightblue!40,
    minimum size=1.6cm, font=\small\sffamily\bfseries, align=center},
  entity/.style={draw=ggunavy, thick, fill=lightpink!50, rounded corners=2pt,
    minimum width=2cm, minimum height=0.9cm, font=\small\sffamily\bfseries, align=center},
  datastore/.style={draw=ggunavy, thick, fill=lightyellow!50,
    minimum width=2.2cm, minimum height=1.0cm, font=\small\sffamily\bfseries, align=center},
  flowlabel/.style={font=\small\sffamily, fill=white, inner sep=1pt},
  arrow/.style={-stealth, ggunavy, thick, shorten >=2pt, shorten <=2pt},
]

% External entities
\node[entity] (eeg) at (-1, 5) {EEG Device};
\node[entity] (clin) at (15, 5) {Clinician};

% Processes (circles)
\node[process] (p1) at (2, 3.5) {1.0\\Receive};
\node[process] (p2) at (5, 3.5) {2.0\\Preprocess};
\node[process] (p3) at (8, 3.5) {3.0\\Extract};
\node[process] (p4) at (11, 3.5) {4.0\\Predict};
\node[process, fill=lightpurple!40] (p5) at (11, 0.8) {5.0\\RAG\\Explain};
\node[process] (p6) at (5, 0.8) {6.0\\Store};

% Data stores
\node[datastore] (d1) at (2, 0.8) {D1: Raw EEG};
\node[datastore] (d2) at (15, 0.8) {D2: Results};
\node[datastore] (db) at (8, -0.8) {Patient DB};
\node[datastore, fill=lightpurple!20] (kb) at (14, -0.8) {D3: Knowledge Base\\(50K abstracts)};

% Arrows: External to Process
\draw[arrow] (eeg) -- node[flowlabel, above, sloped] {EEG Signal} (p1);
\draw[arrow] (p4) -- node[flowlabel, above, sloped] {Diagnosis} (clin);

% Arrows: Process to Process
\draw[arrow] (p1) -- node[flowlabel, above] {Raw Data} (p2);
\draw[arrow] (p2) -- node[flowlabel, above] {Clean Signal} (p3);
\draw[arrow] (p3) -- node[flowlabel, above] {Features} (p4);
\draw[arrow] (p4) -- node[flowlabel, right] {Prediction} (p5);
\draw[arrow] (p5) -- node[flowlabel, above] {Explanation} (p6);

% Arrows: Process to Data Store
\draw[arrow] (p1) -- node[flowlabel, left] {Store Raw} (d1);
\draw[arrow] (p5) -- node[flowlabel, above, sloped] {Results} (d2);
\draw[arrow] (p6) -- node[flowlabel, left] {Archive} (db);

% Arrows: Data Store to Process (read-back)
\draw[arrow, dashed] (d1) -- node[flowlabel, right] {Retrieve} (p2.south west);
\draw[arrow, dashed] (db) -- node[flowlabel, right] {Patient Info} (p5.south);
\draw[arrow, dashed, purple!70!black] (kb) -- node[flowlabel, above, sloped] {FAISS Top-5} (p5.south east);

\end{tikzpicture}%
}% end resizebox
\caption[Data Flow Diagram: Six Processing Stages (Appendix)]{Data flow diagram showing six processing stages (circles), two external entities (rounded rectangles), and three data stores (rectangles). Solid arrows indicate primary data flow; dashed arrows indicate secondary look-up queries.}
\label{app_fig:data_flow_diagram}

\vspace{0.5em}\noindent\textit{Note.} EEG = electroencephalography; DB = database; D1, D2 = data stores for raw signals and computed results respectively. The six processes map to the pipeline stages in the corresponding figure. RAG-based explanation (Process 5.0) retrieves context from the Patient DB before generating clinician-facing narratives.
\end{figure}

%% ============================================================================
\subsection{System Architecture Flows}
\label{app_sec:system_flows}

%% ENGINEERING DETAIL SUPPRESSED — individual flow diagrams moved to Appendix C
\iffalse
%% Eight flow diagrams: Data Flow, Model Training, Hybrid DL, RAG Explainability,
%% GPU/vLLM Deployment, PII Protection, Guardrail AI, Statistical Analysis
%% Full TikZ diagrams preserved in suppressed_content/chapter3_suppressed_content.tex
\fi

%% Phantom labels for cross-references (fig:data_flow now in figures/fig_data_flow.tex)
\label{app_fig:model_flow}
\label{app_fig:hybrid_flow}
\label{app_fig:rag_flow}
\label{app_fig:gpu_deployment_flow}
\label{app_fig:pii_flow}
\label{app_fig:guardrail_flow}
\label{app_fig:statistical_flow}

The RGAIG pipeline comprises eight processing flows, each governed by quality gates that enforce fail-fast semantics: a stage that fails its gate triggers retry or human escalation, never silent pass-through. This governance-by-design ensures that every data transformation, model prediction, and explanation generation passes through a defined checkpoint before reaching the next stage. Table~\ref{tab:system_flows_summary} consolidates all eight flows with their quality gates and output artefacts. Individual flow diagrams are documented in Appendix~\ref{chap:appendix_ch3}.
\needspace{8\baselineskip}
\noindent Table~\ref{app_tab:system_flows_summary} end-to-end system architecture: eight processing flows.

{\tablestripes\tablefontsize
\begin{xltabular}{\textwidth}{@{} p{0.6cm} p{1.8cm} L p{2cm} L @{}}
\caption{End-to-End System Architecture: Eight Processing Flows}\label{app_tab:system_flows_summary} \\
\toprule
\thead{\#} & \thead{Flow} & \thead{Key Stages} & \thead{Quality Gate} & \thead{Output Artefact} \\
\midrule
\endfirsthead
\multicolumn{5}{@{}l}{\textit{(\,Tab.~\ref{app_tab:system_flows_summary} continued from previous page\,)}} \\
\toprule
\thead{\#} & \thead{Flow} & \thead{Key Stages} & \thead{Quality Gate} & \thead{Output Artefact} \\
\midrule
\endhead
\bottomrule
\endlastfoot
1 & Data Flow (Fig.~\ref{fig:data_flow}) & Ingestion $\to$ Validation $\to$ Preprocessing $\to$ Feature extraction & SNR $>$10~dB; completeness $>$95\% & 127-feature vector per sample \\
\addlinespace
2 & Model Training Flow (Fig.~\ref{fig:model_flow}) & Stratified split $\to$ 5-fold CV $\to$ DL training $\to$ Adaptive model selection & All folds converge; no data leakage & Best model per domain (VotingClassifier for EEG; GAN+ResNet-18 for cancer) \\
\addlinespace
3 & Ensemble Flow (Fig.~\ref{fig:hybrid_flow}) & VotingClassifier (ExtraTrees + RF, soft voting) for EEG; DL ensemble compared & Ensemble $>$ individual models & Final prediction with confidence \\
\addlinespace
4 & RAG Explainability Flow (Fig.~\ref{fig:rag_flow}) & Prediction $\to$ Query $\to$ FAISS retrieval $\to$ LLM generation $\to$ Hallucination check & Citation accuracy $\geq$90\%; coherence $\geq$4/5 & Evidence-grounded clinical narrative \\
\addlinespace
5 & GPU/vLLM Deployment (Fig.~\ref{fig:gpu_deployment_flow}) & Model export $\to$ Quantisation $\to$ vLLM serving $\to$ API gateway & Latency $<$200~ms; accuracy drop $<$1\% & Production REST API endpoint \\
\addlinespace
6 & PII Protection Flow (Fig.~\ref{fig:pii_flow}) & PII detection $\to$ De-identification $\to$ Encryption $\to$ Access control $\to$ Audit & $k$-anonymity $\geq$5; zero PII in model input & HIPAA/GDPR-compliant data pipeline \\
\addlinespace
7 & Guardrail AI Flow (Fig.~\ref{fig:guardrail_flow}) & Input validation $\to$ Content filter $\to$ Hallucination detection $\to$ Confidence threshold & All 6 gates pass; else human escalation & Clinically safe, verified output \\
\addlinespace
8 & Statistical Flow (Fig.~\ref{fig:statistical_flow}) & Hypothesis $\to$ Test selection $\to$ CV execution $\to$ Bootstrap CI $\to$ Effect size $\to$ Bonferroni & $p < .05$; CI excludes null; $d > 0.5$ & Accept/reject verdict per hypothesis \\
\end{xltabular}}
\vspace{0.5em}\noindent\textit{Note.} Each flow operates autonomously but shares data through structured JSON messages via the multi-agent orchestration layer (Figure~\ref{fig:agent_responsibilities}). Quality gates enforce fail-fast semantics: a stage that fails its gate triggers retry or human escalation, never silent pass-through. vLLM = virtual large language model serving engine; FAISS = Facebook AI Similarity Search; PII = personally identifiable information.

%% ============================================================================
%% A5: MODEL CONTEXT PROTOCOL (MCP)
%% ============================================================================
\subsection{Model Context Protocol (MCP)}
\label{app_sec:mcp}

The Model Context Protocol (MCP) provides a standardised governance layer that controls how AI agents access tools, data sources, and models within the RGAIG pipeline. MCP governs how autonomous agents access tools, data, and models: agents must operate within defined safety boundaries while retaining sufficient autonomy to execute pipeline tasks. Without a formal context protocol, each agent would require bespoke access controls, creating governance gaps and audit inconsistencies.

The six MCP capabilities implemented in RGAIG---model registry, version control, audit logging, role-based access control, deployment gates, and automated rollback---are documented in Table~\ref{tab:mcp_capabilities}. Figure~\ref{fig:mcp_flow} illustrates the governance flow through which every agent request is validated.
\needspace{8\baselineskip}
\noindent Table~\ref{app_tab:mcp_capabilities} model context protocol: six governance capabilities.

{\tablestripes\tablefontsize
\begin{xltabular}{\textwidth}{@{} p{2cm} L L p{3.2cm} @{}}
\caption{Model Context Protocol: Six Governance Capabilities}\label{app_tab:mcp_capabilities} \\
\toprule
\thead{Capability} & \thead{Function} & \thead{Implementation} & \thead{Enforcement} \\
\midrule
\endfirsthead
\multicolumn{4}{@{}l}{\textit{(\,Tab.~\ref{app_tab:mcp_capabilities} continued from previous page\,)}} \\
\toprule
\thead{Capability} & \thead{Function} & \thead{Implementation} & \thead{Enforcement} \\
\midrule
\endhead
\bottomrule
\endlastfoot
Model registry & Central catalogue of all deployed models with version, hash, and metadata & MLflow model registry; SHA-256 integrity hash per model artefact & No model loads without registry entry \\
\addlinespace
Version control & Track model lineage: data version $\to$ feature version $\to$ model version $\to$ deployment & Git-based versioning; DVC for data and model artefacts & Deployment blocked if version chain incomplete \\
\addlinespace
Audit logs & Immutable record of all model predictions, inputs, outputs, and configuration & Append-only SQLite log; daily rotation; 90-day retention & Every inference call logged; no silent execution \\
\addlinespace
Access control & Role-based permissions: who can train, deploy, query, or rollback models & RBAC policy engine; three roles (admin, clinician, patient) & API key + role verification per request \\
\addlinespace
Deployment gates & Pre-deployment quality checks: accuracy $\geq$85\%, fairness variance $\leq$2\%, latency $<$200\,ms & Automated gate in CI/CD pipeline; all gates must pass & Failed gate blocks deployment; alerts to admin \\
\addlinespace
Rollback & Instant reversion to previous model version if post-deployment metrics degrade & Blue-green deployment; previous version kept hot & Auto-rollback if accuracy drops $>$3\% in 24h \\
\end{xltabular}}
\vspace{0.5em}\noindent\textit{Note.} MCP = Model Context Protocol; RBAC = role-based access control; DVC = Data Version Control; CI/CD = continuous integration/continuous deployment. MCP operates at the orchestration layer, governing all seven pipeline agents (Figure~\ref{fig:agent_responsibilities}).
\needspace{20\baselineskip}
\begin{figure}[!htbp]
\centering
\resizebox{0.97\textwidth}{!}{%
\begin{tikzpicture}[
  node distance=0.8cm and 0.15cm,
  block/.style={rectangle, draw=ggunavy, fill=ggunavy!8, rounded corners=3pt,
    minimum height=0.9cm, minimum width=2.0cm, align=center, font=\small},
  policyblock/.style={block, fill=ggunavy!20},
  arrow/.style={-{Stealth[length=5pt]}, thick, ggunavy},
  redarrow/.style={-{Stealth[length=5pt]}, thick, red!60!black, dashed}
]
\node[block, fill=lightblue] (req) {Agent\\Request};
\node[policyblock, fill=lightorange, right=of req] (policy) {MCP Policy\\Engine};
\node[policyblock, fill=lightteal, right=of policy] (auth) {Access\\Control};
\node[block, fill=lightpurple, right=of auth] (tool) {Tool / Model\\Execution};
\node[policyblock, fill=lightgreen, right=of tool] (audit) {Audit\\Logger};
\node[block, fill=lightyellow, right=of audit] (resp) {Response\\to Agent};
% Rejection path
\node[block, below=0.6cm of auth, fill=red!10, draw=red!60!black] (deny) {Deny +\\Alert Admin};
\draw[arrow] (req) -- (policy);
\draw[arrow] (policy) -- (auth);
\draw[arrow] (auth) -- (tool);
\draw[arrow] (tool) -- (audit);
\draw[arrow] (audit) -- (resp);
\draw[redarrow] (auth) -- (deny);
\node[below=0.08cm of policy, font=\small\itshape, text=gray] {6 capabilities};
\node[below=0.08cm of tool, font=\small\itshape, text=gray] {sandboxed};
\node[below=0.08cm of audit, font=\small\itshape, text=gray] {immutable log};
\end{tikzpicture}%
}%
\caption{Model Context Protocol (MCP) Governance Flow}
\label{app_fig:mcp_flow}
\vspace{0.5em}\noindent\textit{Note.} MCP = Model Context Protocol. Every agent request passes through the policy engine, access control, and audit logger before reaching the target tool or model. Failed access control checks trigger denial with admin notification. The audit logger creates an immutable record of all operations for regulatory compliance.
\end{figure}

%% ============================================================================
%% MULTI-AGENT MCP ARCHITECTURE (from NeuroMCP paper)
%% ============================================================================

The NeuroMCP system~\cite{asthana2025eegmaster} extends the governance-level MCP (Table~\ref{tab:mcp_capabilities}) into a four-layer multi-agent architecture where disease-specific agents operate autonomously under centralised orchestration. Figure~\ref{fig:neuromcp_layers} illustrates the four architectural layers. This architecture avoids the retraining burden of monolithic AI systems: adding a new disease domain requires training only a new agent, not rebuilding the entire system.
\needspace{20\baselineskip}
\begin{figure}[!htbp]
\centering
\resizebox{0.97\textwidth}{!}{%
\begin{tikzpicture}[
  layer/.style={rectangle, draw=ggunavy, rounded corners=4pt,
    minimum width=13.5cm, minimum height=1.6cm, align=center, font=\small, inner sep=6pt},
  l1/.style={layer, fill=ggunavy!30},
  l2/.style={layer, fill=ggunavy!20},
  l3/.style={layer, fill=ggunavy!12},
  l4/.style={layer, fill=ggunavy!5},
  agentbox/.style={rectangle, draw=ggunavy, fill=white, rounded corners=2pt,
    minimum height=1.0cm, minimum width=2.8cm, align=center, font=\small, inner sep=3pt},
  arr/.style={-{Stealth[length=4mm, width=3mm]}, very thick, ggunavy},
  lbl/.style={font=\small\bfseries, text=ggunavy, anchor=east},
  node distance=0.8cm
]
% Layer 1: Model Context Portal
\node[l1, fill=lightblue] (layer1) {\textbf{\color{ggunavy}Layer 1: Model Context Portal} --- REST API gateway, authentication, rate limiting, request routing};

% Layer 2: MCP Orchestrator
\node[l2, fill=lightorange, below=of layer1] (layer2) {\textbf{\color{ggunavy}Layer 2: MCP Orchestrator} --- JSON-RPC 2.0, tool registry, agent lifecycle management, audit logging};

% Layer 3: Disease-Specific Agents
\node[l3, fill=lightteal, below=of layer2] (layer3) {};
\node[font=\small\bfseries, text=ggunavy] at ([yshift=0.3cm]layer3.center) {Layer 3: Disease-Specific Agents};
\node[agentbox, fill=lightgreen] (asd) at ([xshift=-4cm, yshift=-0.2cm]layer3.center) {ASD Agent};
\node[agentbox, fill=lightyellow] (pd) at ([yshift=-0.2cm]layer3.center) {PD Agent};
\node[agentbox, fill=lightpink] (stress) at ([xshift=4cm, yshift=-0.2cm]layer3.center) {Stress Agent};

% Layer 4: Data Processing Pipeline
\node[l4, fill=lightpurple, below=of layer3] (layer4) {\textbf{\color{ggunavy}Layer 4: Data Processing Pipeline} --- EEG preprocessing, feature extraction, model inference, RAG generation};

% Arrows between layers
\draw[arr] (layer1) -- (layer2);
\draw[arr] (layer2) -- (layer3);
\draw[arr] (layer3) -- (layer4);

% Layer labels
\node[lbl] at ([xshift=-0.3cm]layer1.west) {};
\end{tikzpicture}%
}
\caption{NeuroMCP Four-Layer Multi-Agent Architecture}
\label{app_fig:neuromcp_layers}
\vspace{0.5em}\noindent\textit{Note.} MCP = Model Context Protocol; REST = representational state transfer; JSON-RPC = JavaScript Object Notation Remote Procedure Call. Each disease-specific agent encapsulates its own preprocessing, model, and explanation pipeline. The orchestrator layer manages agent lifecycle, tool access, and audit logging via JSON-RPC 2.0. \textit{Source:}~\cite{asthana2025eegmaster}.
\end{figure}

%% ============================================================================
%% A6: GUARDRAIL AI (EXPANDED INPUT/OUTPUT RULES)
%% ============================================================================
\label{app_sec:guardrail_rules}

The Guardrail AI flow (Figure~\ref{fig:guardrail_flow}) implements six safety gates. Tables~\ref{tab:ch3_input_guardrails} and~\ref{tab:ch3_output_guardrails} document the specific input and output rules that govern clinical content generation. The guardrail design reflects the principle that patient-facing AI systems require \textit{both} input and output guardrails: input guardrails prevent dangerous queries from reaching the model, while output guardrails prevent unsafe content from reaching the patient.
\needspace{8\baselineskip}
\noindent Table~\ref{app_tab:ch3_input_guardrails} input guardrail rules: five pre-processing safety checks.

{\tablestripes\tablefontsize
\begin{xltabular}{\textwidth}{@{} p{2cm} L L p{3.5cm} @{}}
\caption{Input Guardrail Rules: Five Pre-Processing Safety Checks}\label{app_tab:ch3_input_guardrails} \\
\toprule
\thead{Rule} & \thead{Trigger Condition} & \thead{Action} & \thead{Example} \\
\midrule
\endfirsthead
\multicolumn{4}{@{}l}{\textit{(\,Tab.~\ref{app_tab:ch3_input_guardrails} continued from previous page\,)}} \\
\toprule
\thead{Rule} & \thead{Trigger Condition} & \thead{Action} & \thead{Example} \\
\midrule
\endhead
\bottomrule
\endlastfoot
Block diagnosis intent & Query requests definitive diagnosis (``do I have cancer?'') & Block; redirect to clinician & ``Please consult your healthcare provider for diagnosis'' \\
\addlinespace
Block prescription & Query requests medication or dosage (``what drug should I take?'') & Block; display disclaimer & ``RGAIG does not provide prescription advice'' \\
\addlinespace
Redirect high-risk & Query mentions self-harm, suicide, or emergency symptoms & Redirect to crisis hotline + human clinician & Display crisis resources (988, emergency contacts) \\
\addlinespace
Refuse out-of-scope & Query outside the five clinical domains (ASD, PD, Cancer, Stress, BCI) & Refuse with explanation & ``This system is designed for neurodegenerative conditions'' \\
\addlinespace
Redact PII & Query contains patient name, DOB, SSN, or medical record number & Redact before processing; log detection event & ``John Smith'' $\to$ ``[PATIENT]'' \\
\end{xltabular}}
\vspace{0.5em}\noindent\textit{Note.} PII = personally identifiable information; DOB = date of birth; SSN = Social Security number. Rules are evaluated in priority order; the first matching rule terminates processing.

\noindent Complementing the input rules, the output guardrails intercept every generated response before delivery to the patient or clinician. Table~\ref{app_tab:ch3_output_guardrails} documents the five post-generation safety checks, their detection methods, corrective actions, and activation thresholds.
\needspace{8\baselineskip}
{\tablestripes\tablefontsize
\begin{xltabular}{\textwidth}{@{} p{2.2cm} L L p{2.8cm} @{}}
\caption{Output Guardrail Rules: Five Post-Generation Safety Checks}\label{app_tab:ch3_output_guardrails} \\
\toprule
\thead{Rule} & \thead{Detection Method} & \thead{Action} & \thead{Threshold} \\
\midrule
\endfirsthead
\multicolumn{4}{@{}l}{\textit{(\,Tab.~\ref{app_tab:ch3_output_guardrails} continued from previous page\,)}} \\
\toprule
\thead{Rule} & \thead{Detection Method} & \thead{Action} & \thead{Threshold} \\
\midrule
\endhead
\bottomrule
\endlastfoot
Downgrade medical authority & Tone classifier detects authoritative medical claims & Inject hedging language (``this may indicate,'' ``consider consulting'') & Authority score $>$0.7 \\
\addlinespace
Inject uncertainty & Confidence below threshold for classification & Add explicit uncertainty statement to output & Model confidence $<$85\% \\
\addlinespace
Block unsupported claims & Claim-to-evidence mapping finds ungrounded assertions & Remove claim; replace with ``insufficient evidence'' & Zero matching passages \\
\addlinespace
Hallucination gate & Faithfulness score below threshold & Replace output with ``I don't have enough information'' & Faithfulness $<$0.70 \\
\addlinespace
Human review escalation & Sensitive topic detected (end-of-life, paediatric, mental health) & Flag for clinician review before patient display & Topic classifier confidence $>$0.8 \\
\end{xltabular}}
\vspace{0.5em}\noindent\textit{Note.} Output rules are evaluated sequentially; multiple rules may apply to the same output. Human review escalation is the final safety net before content reaches the patient.

%% ============================================================================
%% A7: PII PROTECTION AND SECURITY AI
%% ============================================================================

Table~\ref{tab:security_alignment} maps RGAIG security controls to five regulatory frameworks, demonstrating compliance readiness across encryption, access control, audit, pseudonymisation, and human oversight requirements.
\needspace{8\baselineskip}
\noindent Table~\ref{app_tab:security_alignment} security framework alignment: five regulatory standards.

{\tablestripes\tablefontsize
\begin{xltabular}{\textwidth}{@{} p{2.5cm} L L @{}}
\caption{Security Framework Alignment: Five Regulatory Standards}\label{app_tab:security_alignment} \\
\toprule
\thead{Regulation} & \thead{Key Requirement} & \thead{RGAIG Control} \\
\midrule
\endfirsthead
\multicolumn{3}{@{}l}{\textit{(\,Tab.~\ref{app_tab:security_alignment} continued from previous page\,)}} \\
\toprule
\thead{Regulation} & \thead{Key Requirement} & \thead{RGAIG Control} \\
\midrule
\endhead
\bottomrule
\endlastfoot
HIPAA \S164.312 & Technical safeguards: access control, audit, encryption, integrity & RBAC, AES-256, audit trail, hash verification \\
\addlinespace
GDPR Art.~25/32 & Data protection by design; encryption; pseudonymisation & Privacy-by-design architecture; $k$-anonymity $\geq$5 \\
\addlinespace
NIST AI RMF & Risk management throughout AI lifecycle; trustworthy AI & 15 governance layers (A1--A15); continuous monitoring \\
\addlinespace
ISO 27001 & Information security management system; risk assessment & Documented risk register; annual review cycle \\
\addlinespace
EU AI Act Art.~14 & Human oversight for high-risk AI systems & Human-in-the-loop escalation; confidence thresholds \\
\addlinespace
ISO 42001 & AI management system: governance, risk, lifecycle controls for AI & MCP governance layer; model registry; deployment gates; audit logging \\
\addlinespace
ISO 13485 & Medical device quality management system (QMS); design controls & Requirements traceability; design verification/validation; CAPA process \\
\addlinespace
ISO 14971 & Risk management for medical devices; hazard analysis throughout lifecycle & Clinical risk register; FMEA per pipeline stage; residual risk acceptance \\
\end{xltabular}}
\vspace{0.5em}\noindent\textit{Note.} HIPAA = Health Insurance Portability and Accountability Act; GDPR = General Data Protection Regulation; NIST AI RMF = National Institute of Standards and Technology AI Risk Management Framework; EU AI Act = European Union Artificial Intelligence Act; ISO = International Organization for Standardization; CAPA = corrective and preventive action; FMEA = failure mode and effects analysis.

While Table~\ref{tab:security_alignment} maps security controls to five regulatory frameworks, Table~\ref{tab:compliance_matrix} provides a component-level compliance mapping, documenting the specific regulatory mechanism that each RGAIG pipeline component implements for each of five regulatory regimes.
\needspace{8\baselineskip}
\noindent Table~\ref{app_tab:compliance_matrix} regulatory compliance alignment: rgaig components.

{\tablestripes\tablefontsize
\begin{xltabular}{\textwidth}{@{} p{2.2cm} L L L L L @{}}
\caption[Regulatory Compliance Alignment: RGAIG Components]{Regulatory Compliance Alignment: RGAIG Components $\times$ Five Regulatory Frameworks}\label{app_tab:compliance_matrix} \\
\toprule
\thead{Component} & \thead{EU AI Act} & \thead{FDA SaMD} & \thead{GDPR} & \thead{HIPAA} & \thead{NIST AI RMF} \\
\midrule
\endfirsthead
\multicolumn{6}{@{}l}{\textit{(\,Tab.~\ref{app_tab:compliance_matrix} continued from previous page\,)}} \\
\toprule
\thead{Component} & \thead{EU AI Act} & \thead{FDA SaMD} & \thead{GDPR} & \thead{HIPAA} & \thead{NIST AI RMF} \\
\midrule
\endhead
\bottomrule
\endlastfoot
EEG pipeline & Risk assessment & Clinical validation & Data minimisation & De-identification & Map function \\
\addlinespace
Classifier & Transparency log & Performance data & Art.~22 explanation & Access control & Measure function \\
\addlinespace
XAI engine & Human oversight & Intended use doc & Right to explanation & Audit trail & Govern function \\
\addlinespace
RAG system & Accuracy obligation & Labelling requirement & Consent for profiling & Minimum necessary & Manage function \\
\addlinespace
Multi-agent & Robustness testing & Software lifecycle & Data processing record & Business associate & Map + Measure \\
\addlinespace
Governance (525) & Conformity assessment & Quality system & DPIA required & Risk analysis & All four functions \\
\addlinespace
Drift monitor & Post-market monitoring & Vigilance reporting & Storage limitation & Breach notification & Measure function \\
\addlinespace
Wearable (Emotiv) & CE marking & 510(k) pathway & Device consent & Device security & Map function \\
\end{xltabular}}
\vspace{0.5em}\noindent\textit{Note.} Each cell identifies the specific compliance mechanism that RGAIG implements for that component--regulation pair. No component is exempt from regulatory oversight. DPIA = Data Protection Impact Assessment; SaMD = Software as a Medical Device; NIST AI RMF = National Institute of Standards and Technology AI Risk Management Framework.


% [SPACE-FIX] clearpage removed to eliminate empty page
\subsection{System Architecture Implementation Details}
\label{app:sec8_sys_arch}

\label{app_sec:system_architecture_diagrams}

The following diagrams present the RGAIG system architecture from multiple perspectives, each emphasising governance and decision-making rather than implementation detail. These views are complementary: the C4 model (Figure~\ref{fig:c4_model}) provides structural decomposition; the wearable-to-cloud-to-mobile architecture (Figure~\ref{fig:wearable_cloud_mobile}) shows physical deployment; the network flow (Figure~\ref{fig:network_flow_dba}) maps communication pathways; the data flow (Figure~\ref{fig:data_flow}) traces information through governance gates; the decision sequence (Figure~\ref{fig:sequence_flow}) captures stakeholder interactions; and the model lifecycle (Figure~\ref{fig:model_flow_dba}) shows continuous governance over the AI lifecycle.

\input{figures/fig_c4_model}

\input{figures/fig_wearable_cloud_mobile}

\input{figures/fig_network_flow_dba}

\input{figures/fig_data_flow}

\input{figures/fig_sequence_flow}

\input{figures/fig_model_flow_dba}

Figure~\ref{fig:e2e_architecture} presents the RGAIG end-to-end system architecture spanning all 14 layers from EEG wearable capture through signal processing, AI inference, and clinical delivery, with a cross-cutting governance and security layer enforcing compliance throughout.

\needspace{4\baselineskip}
\input{figures/fig_e2e_architecture}

The agentic model control protocol (Figure~\ref{fig:agentic_control_protocol}) details how six autonomous agents communicate through a centralised message bus with four control gates enforcing schema validation, quality assurance, model validity, and safety checks at every transition. The RAG pipeline flowchart (Figure~\ref{fig:rag_pipeline_flowchart}) visualises the ten-stage retrieval-augmented generation process from medical literature segmentation through to patient-readable report generation, with governance gates at critical transitions. The EEG signal processing pipeline (Figure~\ref{fig:signal_processing_pipeline}) traces raw 14-channel EEG through bandpass filtering (0.5--45\,Hz), artefact removal, and time-frequency analysis---including Wigner-Ville Distribution (WVD), Smoothed Pseudo-Wigner-Ville Distribution (SPWVD), and Smoothed Temporal Wigner-Ville Distribution (STWVD)---before conversion to 2D spectrograms for deep learning classification.

\input{figures/fig_agentic_control_protocol}

\input{figures/fig_multi_agent_orchestration}

The multi-agent orchestration architecture (Figure~\ref{fig:multi_agent_orchestration}) presents the disease-specific agent layer: an Agentic AI Orchestrator delegates to three management agents---Coordinator (task scheduling and load balancing), Validator (result verification and multi-agent consensus), and Governor (RAI compliance and audit logging)---which coordinate seven disease-specific agents via an A2A message bus using JSON-RPC~2.0 with mTLS security and OpenTelemetry distributed tracing. Each disease agent specialises in one neurological domain (Epilepsy, Parkinson, Autism, Schizophrenia, Stress, Alzheimer, Depression), enabling parallel multi-disease screening of a single patient's data.

\input{figures/fig_rag_pipeline_flowchart}

Figure~\ref{fig:rag_llm_pipeline} provides a detailed 17-stage view of the RAG + LLM explainability pipeline, organised into five functional groups---AI prediction input, knowledge preparation, storage and retrieval, generation, and delivery---with cross-stage data flows and RAG quality metrics evaluated on 500 clinical queries.

\needspace{4\baselineskip}
\input{figures/fig_rag_llm_pipeline}

\input{figures/fig_hallucination_detection}

The hallucination detection pipeline (Figure~\ref{fig:hallucination_detection}) addresses a critical challenge in clinical AI: ensuring that LLM-generated diagnostic reports are factually grounded in retrieved evidence. The pipeline employs four complementary detection methods---NLI contradiction checking (94.2\% accuracy), entity verification against knowledge bases (91.8\%), atomic claim decomposition (89.5\%), and self-consistency across multiple generations (87.3\%)---to classify each response as grounded, hallucinated, or uncertain. Hallucinated responses trigger regeneration with stricter prompting constraints; uncertain responses are escalated to clinician review, maintaining the human-in-the-loop principle.

Table~\ref{tab:hallucination_detection_methods} presents the accuracy of each hallucination detection method in the multi-layer verification pipeline.
\needspace{8\baselineskip}
\noindent Table~\ref{app_tab:hallucination_detection_methods} hallucination detection method accuracy.

{\tablestripes\tablefontsize
\begin{xltabular}{\textwidth}{@{} p{3.5cm} C L @{}}
\caption{Hallucination Detection Method Accuracy}\label{app_tab:hallucination_detection_methods} \\
\toprule
\thead{Detection Method} & \thead{Accuracy (\%)} & \thead{Description} \\
\midrule
\endfirsthead
\multicolumn{3}{@{}l}{\textit{(\,Tab.~\ref{app_tab:hallucination_detection_methods} continued from previous page\,)}} \\
\toprule
\thead{Detection Method} & \thead{Accuracy (\%)} & \thead{Description} \\
\midrule
\endhead
\bottomrule
\endlastfoot
NLI contradiction checking & 94.2 & Natural language inference flags contradictions between generated text and source evidence \\
\addlinespace
Entity verification & 91.8 & Cross-references named entities (drugs, conditions, biomarkers) against knowledge bases \\
\addlinespace
Atomic claim decomposition & 89.5 & Decomposes generated text into atomic claims and verifies each independently \\
\addlinespace
Self-consistency & 87.3 & Generates multiple responses and flags inconsistencies across generations \\
\end{xltabular}}
\vspace{0.5em}\noindent\textit{Note.} NLI = natural language inference. Each response is classified as grounded, hallucinated, or uncertain based on the consensus of all four methods. Results are reported in Chapter~\ref{chap:analysis}.

\input{figures/fig_signal_processing_pipeline}

\subsection{Agentic Framework Operations}

Table~\ref{tab:agentic_operations} details the operations performed by each of the six autonomous agents in the RGAIG pipeline. Each agent has defined inputs, outputs, and failure modes; all actions are logged to an immutable audit trail. Documenting these operations is essential because agent autonomy without explicit operational boundaries creates governance gaps---each row below defines the contract that constrains agent behaviour during pipeline execution.
\needspace{8\baselineskip}
\noindent Table~\ref{app_tab:agentic_operations} agentic framework operations: each agent's.

{\tablestripes\tablefontsize
\begin{xltabular}{\textwidth}{@{}>{
aggedrightrraybackslash}p{2.2cm} X p{2.5cm}@{}}
\caption[Agentic framework operations: each agent's]{Agentic framework operations: each agent's responsibilities, outputs, and governance controls.}\label{app_tab:agentic_operations} \\
\toprule
\thead{Agent} & \thead{Operations Performed} & \thead{Output} \\
\midrule
\endfirsthead
\multicolumn{3}{@{}l}{\textit{(\,Tab.~\ref{app_tab:agentic_operations} continued from previous page\,)}} \\
\toprule
\thead{Agent} & \thead{Operations Performed} & \thead{Output} \\
\midrule
\endhead
\bottomrule
\endlastfoot
Data Agent & Dataset discovery, format detection (EDF/CSV/NPZ), schema validation, missing value imputation, cross-dataset alignment, metadata extraction & Validated dataset \\[3pt]
Preprocess Agent & Signal quality assessment, bandpass filtering (0.5--45\,Hz), artefact rejection (ICA), epoch segmentation, channel standardisation, SNR estimation & Clean EEG epochs \\[3pt]
Feature Agent & Band power extraction ($\delta$/$\theta$/$\alpha$/$\beta$/$\gamma$), connectivity metrics, SHAP importance ranking, dimensionality reduction, feature selection (47$\to$25) & Feature matrix \\[3pt]
Model Agent & Domain-adaptive model selection, cross-validation (10-fold/LOSO), ensemble construction, risk probability scoring, confidence estimation & Risk scores \\[3pt]
RAG Agent & Query routing, ChromaDB retrieval (1.17\,GB), context assembly, Llama-3.2 generation, citation verification, faithfulness scoring, hallucination check & Clinical narrative \\[3pt]
Governance Agent & 46-module RAI audit, 6-gate guardrail, bias detection (DI$\geq$0.80), drift monitoring (PSI), escalation routing, compliance logging, override management & Audit report \\
\end{xltabular}}

\subsection{RAG Patient Report Types}

The RAG engine generates twelve distinct report types, each tailored to a specific stakeholder (Table~\ref{tab:rag_reports}). Reports range from patient-facing screening summaries to regulatory audit documentation. This differentiation is necessary because the EU AI Act (Art.~13) and FDA SaMD guidance require that AI outputs be communicated at an appropriate level of detail for each recipient---a patient needs plain-language summaries while a regulator needs full audit trails.
\needspace{8\baselineskip}
\noindent Table~\ref{app_tab:rag_reports} twelve patient report types generated by the rag.

{\tablestripes\tablefontsize
\begin{xltabular}{\textwidth}{@{}l X l@{}}
\caption[Twelve patient report types generated by the RAG]{Twelve patient report types generated by the RAG pipeline, with recipient and purpose.}\label{app_tab:rag_reports} \\
\toprule
\thead{\#} & \thead{Report Type} & \thead{Recipient} \\
\midrule
\endfirsthead
\multicolumn{3}{@{}l}{\textit{(\,Tab.~\ref{app_tab:rag_reports} continued from previous page\,)}} \\
\toprule
\thead{\#} & \thead{Report Type} & \thead{Recipient} \\
\midrule
\endhead
\bottomrule
\endlastfoot
R1 & Screening Summary Report: risk score per domain, confidence level, recommended action & Patient \\[2pt]
R2 & Biomarker Explanation Report: SHAP-ranked features with plain-language explanations & Patient \\[2pt]
R3 & Trend Monitoring Report: longitudinal biomarker changes over weeks/months, drift alerts & Patient \\[2pt]
R4 & Literature-Backed Evidence Report: citations from peer-reviewed studies supporting findings & Patient + Doctor \\[2pt]
R5 & Clinical Alert Notification: HL7 FHIR-formatted urgent alert when thresholds exceeded & Clinician \\[2pt]
R6 & Comparative Domain Report: ASD-focused risk comparison across all the ASD domain & Patient + Doctor \\[2pt]
R7 & Treatment Context Report: relevant treatment options from literature (informational, not prescriptive) & Doctor \\[2pt]
R8 & Model Confidence Report: prediction confidence intervals, uncertainty quantification, limitations & Doctor \\[2pt]
R9 & Caretaker Summary: simplified dashboard for family members/caretakers with key metrics only & Caretaker \\[2pt]
R10 & Regulatory Audit Report: full audit trail, model version, data provenance, governance compliance & Regulator \\[2pt]
R11 & Bias and Fairness Report: demographic parity analysis, disparate impact metrics (DI$\geq$0.80) & Governance \\[2pt]
R12 & Incident and Drift Report: model performance degradation, retraining triggers, root-cause analysis & Operations \\
\end{xltabular}}

\subsection{RAG Analysis Capabilities}

Table~\ref{tab:rag_analyses} enumerates the fifteen analysis types supported by the RAG engine, categorised by execution mode (real-time, batch, or on-demand). The breadth of these capabilities demonstrates that RAG extends beyond simple question-answering to encompass longitudinal tracking, differential screening, and guideline mapping---functions that operationalise the clinical decision support requirements of H5.
\needspace{8\baselineskip}
\noindent Table~\ref{app_tab:rag_analyses} fifteen rag analysis capabilities with execution type.

{\tablestripes\tablefontsize
\begin{xltabular}{\textwidth}{@{}l X l@{}}
\caption[Fifteen RAG analysis capabilities with execution type]{Fifteen RAG analysis capabilities with execution type and domain scope.}\label{app_tab:rag_analyses} \\
\toprule
\thead{\#} & \thead{Analysis Type} & \thead{Type} \\
\midrule
\endfirsthead
\multicolumn{3}{@{}l}{\textit{(\,Tab.~\ref{app_tab:rag_analyses} continued from previous page\,)}} \\
\toprule
\thead{\#} & \thead{Analysis Type} & \thead{Type} \\
\midrule
\endhead
\bottomrule
\endlastfoot
A1 & Biomarker Interpretation: EEG features explained in clinical context & Real-time \\[1pt]
A2 & Risk Stratification: patient risk categorised (low/moderate/high) with evidence thresholds & Real-time \\[1pt]
A3 & ASD-Focused Correlation: shared biomarker patterns across ASD, PD, Stress & Batch \\[1pt]
A4 & Literature Synthesis: relevant research papers retrieved and synthesised & Real-time \\[1pt]
A5 & Temporal Pattern Analysis: week-over-week trajectory and change detection & Batch \\[1pt]
A6 & Differential Screening: feature comparison between conditions & ASD-focused \\[1pt]
A7 & Medication Context: treatment information from literature (not prescriptions) & On-demand \\[1pt]
A8 & Comorbidity Assessment: co-occurring conditions flagged from ASD-focused screening & Batch \\[1pt]
A9 & Quality of Life Impact: biomarker patterns mapped to functional impact & Patient-facing \\[1pt]
A10 & Explainability Narrative: Flesch-Kincaid grade-level summaries for patients & Real-time \\[1pt]
A11 & Feature Importance Ranking: SHAP-based ranking with natural language justification & Real-time \\[1pt]
A12 & Confidence Calibration: uncertainty quantification and known limitations & Real-time \\[1pt]
A13 & False Positive Contextualisation: explains when a result may be a false alarm & On-demand \\[1pt]
A14 & Guideline Mapping: findings mapped to NICE, APA, WHO clinical guidelines & On-demand \\[1pt]
A15 & Research Gap Identification: areas where evidence is limited and caution warranted & On-demand \\
\end{xltabular}}

\subsection{EEG Filter Types and Parameters}

EEG preprocessing requires a cascade of frequency-selective filters to isolate neural activity from noise while preserving clinically relevant frequency bands. Table~\ref{tab:eeg_filters} documents the ten filter types used in the RGAIG EEG preprocessing pipeline, including bandpass, sub-band, notch, and adaptive filters. Each filter's frequency range and implementation method are specified to support independent replication.
\needspace{8\baselineskip}
\noindent Table~\ref{app_tab:eeg_filters} eeg filter types, frequency bands, and implementation details.

{\tablestripes\tablefontsize
\begin{xltabular}{\textwidth}{@{} p{2cm} p{1.5cm} L p{3.5cm} @{}}
\caption{EEG filter types, frequency bands, and implementation details.}\label{app_tab:eeg_filters} \\
\toprule
\thead{Filter Type} & \thead{Band (Hz)} & \thead{Purpose} & \thead{Implementation} \\
\midrule
\endfirsthead
\multicolumn{4}{@{}l}{\textit{(\,Tab.~\ref{app_tab:eeg_filters} continued from previous page\,)}} \\
\toprule
\thead{Filter Type} & \thead{Band (Hz)} & \thead{Purpose} & \thead{Implementation} \\
\midrule
\endhead
\bottomrule
\endlastfoot
Bandpass & 0.5--45 & Remove DC offset and high-frequency noise; retain neural signals & Butterworth 4th order \\[2pt]
Delta ($\delta$) & 0.5--4 & Deep sleep, encephalopathy, PD tremor correlation & Sub-band extraction \\[2pt]
Theta ($\theta$) & 4--8 & Drowsiness, memory encoding, ASD connectivity patterns & Sub-band extraction \\[2pt]
Alpha ($\alpha$) & 8--13 & Relaxation, eyes-closed baseline, attention state & Sub-band extraction \\[2pt]
Beta ($\beta$) & 13--30 & Active cognition, motor planning, stress arousal markers & Sub-band extraction \\[2pt]
Gamma ($\gamma$) & 30--45 & Higher cognition, cross-frequency coupling, BCI features & Sub-band extraction \\[2pt]
Notch & 50/60 & Power line interference removal (country-specific) & IIR notch filter \\[2pt]
High-pass & $>$0.5 & Remove slow electrode drift and DC offset & Butterworth 2nd order \\[2pt]
Low-pass & $<$45 & Anti-aliasing before down-sampling; remove muscle artefact & Butterworth 4th order \\[2pt]
Adaptive & Variable & Real-time artefact tracking for wearable (motion, sweat) & LMS/RLS algorithm \\
\end{xltabular}}

\subsection{Time-Frequency Representations}

Selecting the appropriate time-frequency representation directly affects classification performance because each method trades resolution against artefact susceptibility. Table~\ref{tab:tf_methods} compares the five time-frequency analysis methods evaluated in this research. WVD provides the highest resolution but suffers from cross-term interference; SPWVD and STWVD offer improved stability; STFT and CWT are used for the EEG-to-2D-image conversion that enables deep learning classification. The comparison justifies the framework's selection of CWT as the primary representation for CNN-based domains.
\needspace{8\baselineskip}
\noindent Table~\ref{app_tab:tf_methods} time-frequency representation methods: wvd, spwvd,.

{\tablestripes\tablefontsize
\begin{xltabular}{\textwidth}{@{} p{1.5cm} L p{3.5cm} p{3.5cm} @{}}
\caption[Time-frequency representation methods: WVD, SPWVD,]{Time-frequency representation methods: WVD, SPWVD, STWVD, STFT, and CWT.}\label{app_tab:tf_methods} \\
\toprule
\thead{Method} & \thead{Description} & \thead{Strength} & \thead{Limitation} \\
\midrule
\endfirsthead
\multicolumn{4}{@{}l}{\textit{(\,Tab.~\ref{app_tab:tf_methods} continued from previous page\,)}} \\
\toprule
\thead{Method} & \thead{Description} & \thead{Strength} & \thead{Limitation} \\
\midrule
\endhead
\bottomrule
\endlastfoot
WVD & Wigner-Ville Distribution: bilinear time-frequency representation with maximal resolution & Highest resolution & Cross-term interference \\[3pt]
SPWVD & Smoothed Pseudo-Wigner-Ville: independent smoothing in time and frequency; suppresses cross-terms & Reduced cross-terms & Resolution trade-off \\[3pt]
STWVD & Smoothed Temporal Wigner-Ville: temporal smoothing with stability constraints for non-stationary biosignals & Temporally stable & Computationally intensive \\[3pt]
STFT & Short-Time Fourier Transform: fixed-width windowed FFT; standard spectrogram for 2D conversion & Simple, well-understood & Fixed resolution \\[3pt]
CWT & Continuous Wavelet Transform: variable-resolution analysis using Morlet wavelet; adapts to signal frequency & Adaptive resolution & Redundant, slower \\
\end{xltabular}}

\subsection{EEG to 2D Image Conversion Pipeline}

The conversion of one-dimensional EEG time-series into two-dimensional spectrogram images is the critical transformation that enables convolutional neural network architectures---designed for image data---to classify brain signals. Table~\ref{tab:eeg_2d_pipeline} details the ten sequential steps, from epoch selection through augmentation to final classification input. This pipeline is shared across all four EEG domains, with only the channel selection (Step~2) varying per condition.
\needspace{8\baselineskip}
\noindent Table~\ref{app_tab:eeg_2d_pipeline} ten-step pipeline for converting eeg signals to 2d.

{\tablestripes\tablefontsize
\begin{xltabular}{\textwidth}{@{}l X@{}}
\caption[Ten-step pipeline for converting EEG signals to 2D]{Ten-step pipeline for converting EEG signals to 2D spectrograms for deep learning.}\label{app_tab:eeg_2d_pipeline} \\
\toprule
\thead{Step} & \thead{Process} \\
\midrule
\endfirsthead
\multicolumn{2}{@{}l}{\textit{(\,Tab.~\ref{app_tab:eeg_2d_pipeline} continued from previous page\,)}} \\
\toprule
\thead{Step} & \thead{Process} \\
\midrule
\endhead
\bottomrule
\endlastfoot
1. Epoch selection & Extract fixed-length EEG segments (2\,s windows, 50\% overlap) from clean signal \\[2pt]
2. Channel selection & Select relevant channels per domain (frontal for ASD, motor for PD, all for BCI) \\[2pt]
3. STFT computation & Short-Time Fourier Transform (Hann window, 256-point FFT, 50\% overlap) $\to$ frequency$\times$time matrix \\[2pt]
4. CWT computation & Continuous Wavelet Transform (Morlet wavelet, 40 scales) $\to$ scale$\times$time matrix \\[2pt]
5. Magnitude to dB & Convert complex coefficients to power ($|X|^2$) then to decibels ($10\log_{10}$) \\[2pt]
6. Normalisation & Min-max normalise to [0, 1]; apply colour mapping (viridis/jet) for RGB channels \\[2pt]
7. Image resize & Resize to standard dimensions (224$\times$224 or 128$\times$128) for model input \\[2pt]
8. Augmentation & Time-shift, frequency masking, mixup for training data diversity \\[2pt]
9. Multi-channel stack & Stack multiple EEG channels or frequency bands as image channels \\[2pt]
10. Classification & Feed 2D spectrograms into deep learning model for domain-specific risk scoring \\
\end{xltabular}}

\subsection{Monitoring Framework}

Each RGAIG layer has dedicated monitoring, alerting, and escalation mechanisms (Table~\ref{tab:monitoring_framework}). No layer operates without observability.
\needspace{8\baselineskip}
\noindent Table~\ref{app_tab:monitoring_framework} monitoring framework across five rgaig layers with.

{\tablestripes\tablefontsize
\begin{xltabular}{\textwidth}{@{}l X@{}}
\caption[Monitoring framework across five RGAIG layers with]{Monitoring framework across five RGAIG layers with metrics and alerting.}\label{app_tab:monitoring_framework} \\
\toprule
\thead{Layer} & \thead{Monitoring Capabilities} \\
\midrule
\endfirsthead
\multicolumn{2}{@{}l}{\textit{(\,Tab.~\ref{app_tab:monitoring_framework} continued from previous page\,)}} \\
\toprule
\thead{Layer} & \thead{Monitoring Capabilities} \\
\midrule
\endhead
\bottomrule
\endlastfoot
L1: Data Foundation & Data completeness, schema validation pass rate, PII scrub log, missing value alerts, ingestion latency \\[3pt]
L2: Analytics Engine & ASD ensemble accuracy, AUC per domain, model drift (PSI $<$0.10 threshold), CUSUM change-point detection, inference latency \\[3pt]
L3: Trust \& Accountability & SHAP coverage ($\geq$95\% threshold), RAG faithfulness ($\geq$0.85 threshold), hallucination rate ($\leq$5\% threshold), bias audit (DI$\geq$0.80), fairness metrics \\[3pt]
L4: Decision Orchestration & Agent health check (5\,s), pipeline latency, retry count, guardrail gate pass/fail ratio, override log \\[3pt]
L5: Strategic Governance & 15 KPIs tracked, compliance score, escalation count, audit trail completeness, SLA adherence \\
\end{xltabular}}

\subsection{Observability, Logging, and Tracing}

Table~\ref{tab:observability} documents the eight observability capabilities embedded in the RGAIG platform architecture.
\needspace{8\baselineskip}
\noindent Table~\ref{app_tab:observability} observability stack: monitoring, logging, tracing, and.

{\tablestripes\tablefontsize
\begin{xltabular}{\textwidth}{@{} >{
aggedrightrraybackslash}p{3cm} L p{3.5cm} @{}}
\caption[Observability stack: monitoring, logging, tracing, and]{Observability stack: monitoring, logging, tracing, and alerting capabilities.}\label{app_tab:observability} \\
\toprule
\thead{Capability} & \thead{Implementation in RGAIG} & \thead{Tool/Standard} \\
\midrule
\endfirsthead
\multicolumn{3}{@{}l}{\textit{(\,Tab.~\ref{app_tab:observability} continued from previous page\,)}} \\
\toprule
\thead{Capability} & \thead{Implementation in RGAIG} & \thead{Tool/Standard} \\
\midrule
\endhead
\bottomrule
\endlastfoot
Monitoring & Real-time KPI dashboard (8 metrics $\times$ 3 thresholds); PSI drift detection; CUSUM alerts & Prometheus/Grafana \\[3pt]
Observability & Distributed tracing across 6 agents; pipeline stage timing; resource utilisation per container & OpenTelemetry \\[3pt]
Logging & Structured JSON logs with correlation ID; agent-level audit trail; immutable log storage & ELK Stack \\[3pt]
Tracing & End-to-end request trace: EEG upload $\to$ preprocessing $\to$ classification $\to$ RAG $\to$ report & Jaeger / Zipkin \\[3pt]
Alerting & 4-level escalation: L1 automated $\to$ L2 engineer $\to$ L3 committee $\to$ L4 board & PagerDuty \\[3pt]
Health Checks & Agent heartbeat (5\,s); liveness and readiness probes; circuit breaker pattern & Kubernetes probes \\[3pt]
Audit Trail & Immutable append-only log; 21 CFR Part 11 compliant; data provenance chain & Blockchain hash \\[3pt]
Performance & Inference latency $<$500\,ms; throughput metrics; p50/p95/p99 percentiles; error rate & APM dashboards \\
\end{xltabular}}

\subsection{Scalability Architecture}

The RGAIG platform is designed for horizontal scaling to support up to 10,000 concurrent patients (Table~\ref{tab:scalability}). The architecture uses container orchestration, message queues, and auto-scaling to maintain performance under load.
\needspace{8\baselineskip}
\noindent Table~\ref{app_tab:scalability} scalability metrics: 1,000 vs. 10,000 concurrent users.

{\tablestripes\tablefontsize
\begin{xltabular}{\textwidth}{@{} >{\raggedright\arraybackslash}p{3.5cm} L L @{}}
\caption{Scalability metrics: 1,000 vs.\ 10,000 concurrent users.}\label{app_tab:scalability} \\
\toprule
\thead{Metric} & \thead{1,000 Users} & \thead{10,000 Users} \\
\midrule
\endfirsthead
\multicolumn{3}{@{}l}{\textit{(\,Tab.~\ref{app_tab:scalability} continued from previous page\,)}} \\
\toprule
\thead{Metric} & \thead{1,000 Users} & \thead{10,000 Users} \\
\midrule
\endhead
\bottomrule
\endlastfoot
API pods & 3 replicas & 15--30 replicas (auto-scale) \\
Agent pods & 6 (1 per agent) & 30--60 (5--10 per agent) \\
Inference latency & $<$200\,ms (p95) & $<$500\,ms (p95) \\
EEG throughput & 50 uploads/min & 500 uploads/min \\
Database & Single node & Primary + 3 read replicas \\
Cache hit rate & 60\% & 85\% (warm cache) \\
Storage & 30\,GB/month & 300\,GB/month (tiered) \\
Monthly cost & \$2,100 & \$8,500 (economies of scale) \\
\end{xltabular}}

\subsection{Large Data Handling}

Table~\ref{tab:large_data} documents the strategies used at each pipeline stage to manage the 305\,GB data volume across six datasets (131 unique subjects plus 20{,}000 images, yielding 768+ total data instances including augmented epochs).
\needspace{8\baselineskip}
\noindent Table~\ref{app_tab:large_data} large data handling strategy at each pipeline stage.

{\tablestripes\tablefontsize
\begin{xltabular}{\textwidth}{@{} >{
aggedrightrraybackslash}p{3cm} L p{3.5cm} @{}}
\caption{Large data handling strategy at each pipeline stage.}\label{app_tab:large_data} \\
\toprule
\thead{Stage} & \thead{Large Data Strategy} & \thead{Volume} \\
\midrule
\endfirsthead
\multicolumn{3}{@{}l}{\textit{(\,Tab.~\ref{app_tab:large_data} continued from previous page\,)}} \\
\toprule
\thead{Stage} & \thead{Large Data Strategy} & \thead{Volume} \\
\midrule
\endhead
\bottomrule
\endlastfoot
Raw Ingestion & Streaming upload via chunked transfer (10\,MB chunks); parallel multi-dataset loader; checksum validation & 58\,GB raw \\[3pt]
Preprocessing & Lazy loading with memory-mapped files (NumPy memmap); batch processing (1,000 epochs/batch); garbage collection & 247\,GB processed \\[3pt]
Feature Store & Columnar storage (Parquet format); partitioned by domain and subject; indexed for fast retrieval & 12\,GB features \\[3pt]
Model Training & Mini-batch gradient descent (batch=64); data generators; distributed training across GPUs & 198 models \\[3pt]
RAG Knowledge & ChromaDB vector database with HNSW indexing; incremental updates; compressed embeddings & 1.17\,GB vectors \\[3pt]
Report Storage & Object storage (S3) with lifecycle policies; hot/warm/cold tiers; 7-year retention (HIPAA) & Growing \\[3pt]
Archival & Glacier/Archive tier for old EEG; auto-transition after 90 days; retrieval within 12 hours & 58+ GB/year \\
\end{xltabular}}

\subsubsection{Knowledge Graph Layer}
\label{app_sec:knowledge_graph}

Beyond tabular and vector storage, the RGAIG architecture specifies a knowledge graph layer for modelling entity relationships that are poorly represented in relational or vector stores. Clinical AI diagnostics involve complex, interconnected entities---patients linked to multiple recording sessions, sessions linked to biomarker profiles, biomarkers linked to disease conditions, conditions linked to intervention pathways, and evidence linked to clinical recommendations. Table~\ref{tab:knowledge_graph} summarises the six primary relationship types.
\needspace{3\baselineskip}
\noindent Table~\ref{app_tab:knowledge_graph} knowledge graph entity relationships.

{\tablestripes\tablefontsize
\begin{xltabular}{\textwidth}{@{} >{\raggedright\arraybackslash}p{0.22\textwidth} >{\raggedright\arraybackslash}p{0.22\textwidth} >{\raggedright\arraybackslash}X @{}}
\caption{Knowledge Graph Entity Relationships}\label{app_tab:knowledge_graph} \\
\toprule
\thead{Source Entity} & \thead{Target Entity} & \thead{Relationship and Purpose} \\
\midrule
\endfirsthead
\multicolumn{3}{@{}l}{\textit{(\,Tab.~\ref{app_tab:knowledge_graph} continued from previous page\,)}} \\
\toprule
\thead{Source Entity} & \thead{Target Entity} & \thead{Relationship and Purpose} \\
\midrule
\endhead
\bottomrule
\endlastfoot
Patient & Recording Session & \texttt{HAS\_SESSION}: Links longitudinal EEG recordings to enable trend analysis and drift detection across visits \\
\addlinespace
Session & Biomarker Profile & \texttt{PRODUCES}: Maps each session to extracted features (theta power, alpha asymmetry, spectral entropy) for per-session explainability \\
\addlinespace
Biomarker & Disease Condition & \texttt{INDICATES}: Encodes domain knowledge from clinical literature---e.g., elevated theta/beta ratio indicates ADHD-subtype ASD \\
\addlinespace
Condition & Intervention Path & \texttt{TRIGGERS}: Links diagnostic outcomes to evidence-based care protocols retrieved via the RAG pipeline \\
\addlinespace
Evidence Chunk & Recommendation & \texttt{SUPPORTS}: Traces every RAG-generated recommendation to its source PubMed abstract or clinical guideline \\
\addlinespace
Clinician & Case & \texttt{REVIEWS}: Captures human-in-the-loop audit trail---which clinician reviewed which AI-generated output and when \\
\end{xltabular}}
\vspace{0.3em}\noindent\textit{Note.} The knowledge graph complements ChromaDB (vector similarity) and PostgreSQL (structured records) by representing traversable relationships that support multi-hop clinical reasoning and provenance auditing.

In the current implementation, entity relationships are encoded as foreign-key references in PostgreSQL and as metadata filters in ChromaDB. Full graph database deployment (e.g., Neo4j) is specified for Phase~2 clinical deployment, where multi-hop queries---such as ``retrieve all sessions for patients whose biomarker trajectory matches early-stage PD''---require efficient graph traversal beyond relational join capacity.

\subsection{Non-Functional Requirements}

Table~\ref{tab:nfrs} specifies the fifteen non-functional requirements across six categories that the RGAIG platform must satisfy.
\needspace{8\baselineskip}
\noindent Table~\ref{app_tab:nfrs} non-functional requirements: 15 requirements across 6 categories.

{\tablestripes\tablefontsize
\begin{xltabular}{\textwidth}{@{} >{
aggedrightrraybackslash}p{2.8cm} p{1.2cm} L p{3.5cm} @{}}
\caption{Non-functional requirements: 15 requirements across 6 categories.}\label{app_tab:nfrs} \\
\toprule
\thead{Category} & \thead{ID} & \thead{Requirement} & \thead{Target} \\
\midrule
\endfirsthead
\multicolumn{4}{@{}l}{\textit{(\,Tab.~\ref{app_tab:nfrs} continued from previous page\,)}} \\
\toprule
\thead{Category} & \thead{ID} & \thead{Requirement} & \thead{Target} \\
\midrule
\endhead
\bottomrule
\endlastfoot
Performance & NFR-1 & End-to-end inference latency & $\leq$500\,ms (p95) \\
Performance & NFR-2 & EEG upload throughput & $\geq$500/min \\
Performance & NFR-3 & RAG report generation time & $\leq$30\,s \\
Scalability & NFR-4 & Concurrent users supported & 10,000 \\
Scalability & NFR-5 & Horizontal auto-scaling & 3--50 pods \\
Scalability & NFR-6 & Database read throughput & 10,000 qps \\
Reliability & NFR-7 & System uptime SLA & 99.5\% \\
Reliability & NFR-8 & Data durability & 99.999\% (S3) \\
Reliability & NFR-9 & Mean time to recovery & $\leq$15\,min \\
Security & NFR-10 & Encryption at rest & AES-256 \\
Security & NFR-11 & Encryption in transit & TLS 1.3 \\
Security & NFR-12 & Authentication & OAuth 2.0 + MFA \\
Compliance & NFR-13 & Regulatory alignment & HIPAA, GDPR, FDA \\
Compliance & NFR-14 & Audit trail retention & 7 years \\
Usability & NFR-15 & System Usability Scale & $\geq$68 (proj.\ 78.8) \\
\end{xltabular}}

\subsection{Notification and Alert System}

Table~\ref{tab:notifications} documents the eight notification types, their trigger conditions, recipients, and delivery channels.
\needspace{8\baselineskip}
\noindent Table~\ref{app_tab:notifications} notification system: 8 notification types with.

{\tablestripes\tablefontsize
\begin{xltabular}{\textwidth}{@{}l X l l@{}}
\caption[Notification system: 8 notification types with]{Notification system: 8 notification types with recipients and delivery channels.}\label{app_tab:notifications} \\
\toprule
\thead{Type} & \thead{Description} & \thead{Recipient} & \thead{Channel} \\
\midrule
\endfirsthead
\multicolumn{4}{@{}l}{\textit{(\,Tab.~\ref{app_tab:notifications} continued from previous page\,)}} \\
\toprule
\thead{Type} & \thead{Description} & \thead{Recipient} & \thead{Channel} \\
\midrule
\endhead
\bottomrule
\endlastfoot
N1 & Screening Complete: report ready to view & Patient & Push + Email \\[2pt]
N2 & Risk Alert: threshold exceeded ($\geq$0.85) & Patient + Doctor & Push + SMS \\[2pt]
N3 & Clinical Alert: HL7 FHIR notification & Clinician & FHIR + Dashboard \\[2pt]
N4 & Trend Change: PSI $\geq$0.10 or CUSUM trigger & Patient + Doctor & Push + Email \\[2pt]
N5 & Model Drift: accuracy dropped $\geq$2\% & Operations & Slack + PagerDuty \\[2pt]
N6 & Governance Escalation: bias or compliance flag & Governance Board & Email + Dashboard \\[2pt]
N7 & Subscription Renewal: payment or plan update & Patient & Push + Email \\[2pt]
N8 & System Health: agent failure or latency spike & Engineering & PagerDuty + Slack \\
\end{xltabular}}

\subsection{Enterprise Architecture Views}

Table~\ref{tab:enterprise_architecture} presents twelve architect-role perspectives on the RGAIG platform, mapping each role's scope, deliverables, and governance layer.
\needspace{8\baselineskip}
\noindent Table~\ref{app_tab:enterprise_architecture} enterprise architecture: 12 architect views of the rgaig platform.

{\tablestripes\tablefontsize
\begin{xltabular}{\textwidth}{@{}>{\raggedright\arraybackslash}p{2.2cm} X l@{}}
\caption{Enterprise architecture: 12 architect views of the RGAIG platform.}\label{app_tab:enterprise_architecture} \\
\toprule
\thead{Architect Role} & \thead{RGAIG Scope and Deliverables} & \thead{Layer} \\
\midrule
\endfirsthead
\multicolumn{3}{@{}l}{\textit{(\,Tab.~\ref{app_tab:enterprise_architecture} continued from previous page\,)}} \\
\toprule
\thead{Architect Role} & \thead{RGAIG Scope and Deliverables} & \thead{Layer} \\
\midrule
\endhead
\bottomrule
\endlastfoot
Business & B2C subscription strategy, ROI 135--185\%, 15-KPI dashboard, Blue Ocean positioning & L5 \\[2pt]
Data & 305\,GB pipeline, Parquet feature store, ChromaDB 1.17\,GB, data lineage, 7-year HIPAA retention & L1 \\[2pt]
Platform & AWS/Azure cloud, Kubernetes auto-scale (3--50 pods), FPGA edge, PostgreSQL + Redis & L2, L4 \\[2pt]
DevOps & CI/CD pipeline, Docker, Helm, Terraform IaC, blue-green deployment, health probes & L4 \\[2pt]
LLMOps & Llama-3.2 lifecycle, prompt versioning, RAG evaluation, model registry, guardrail tuning & L3, L4 \\[2pt]
Software & Python 3.11 (agenticfinder), 6-agent microservices, A2A MessageBus, REST/WebSocket & L2 \\[2pt]
Integration & HL7 FHIR R4, Emotiv BLE SDK, ChromaDB API, SHAP, Flutter bridge, OAuth 2.0 & L2, L4 \\[2pt]
Test & 525-module taxonomy, cross-validation, bootstrap CI, expert panel, UAT & L2, L3 \\[2pt]
Security & AES-256, TLS 1.3, OAuth 2.0 + MFA, RBAC, API rate limiting, WAF, DDoS protection & All \\[2pt]
Compliance & FDA SaMD, EU AI Act, HIPAA/GDPR, NIST AI RMF, ISO 13485, consent management & L5 \\[2pt]
Encryption & AES-256 at rest, TLS 1.3 in transit, HSM-backed keys, FIPS 140-2, certificate pinning & All \\[2pt]
PII Protection & 7-stage pipeline: pseudonymisation, tokenisation, access control, audit, retention, deletion, breach & L1, L5 \\
\end{xltabular}}

\subsection{Patient Pain Points and Solutions}

Table~\ref{tab:patient_pain} maps ten patient pain points to the specific RGAIG processes, analyses, tracking mechanisms, and monitoring capabilities that address each one.
\needspace{8\baselineskip}
\noindent Table~\ref{app_tab:patient_pain} patient pain points addressed through rgaig process,.

{\tablestripes\tablefontsize
\begin{xltabular}{\textwidth}{@{}l p{2.5cm} X l@{}}
\caption[Patient pain points addressed through RGAIG process,]{Patient pain points addressed through RGAIG process, analysis, tracking, and monitoring.}\label{app_tab:patient_pain} \\
\toprule
\thead{\#} & \thead{Patient Pain} & \thead{RGAIG Solution} & \thead{Evidence} \\
\midrule
\endfirsthead
\multicolumn{4}{@{}l}{\textit{(\,Tab.~\ref{app_tab:patient_pain} continued from previous page\,)}} \\
\toprule
\thead{\#} & \thead{Patient Pain} & \thead{RGAIG Solution} & \thead{Evidence} \\
\midrule
\endhead
\bottomrule
\endlastfoot
P1 & Long wait for diagnosis & Home EEG screening within minutes; early pattern detection & Ch.~\ref{chap:analysis} \\[2pt]
P2 & Late-stage detection & Continuous biomarker tracking detects subtle changes before symptoms & Ch.~\ref{chap:analysis} \\[2pt]
P3 & Black-box AI rejection & RAG plain-language reports with literature citations & Ch.~\ref{chap:analysis} \\[2pt]
P4 & Fragmented tools & One EEG scan, ASD screening & Ch.~\ref{chap:analysis} \\[2pt]
P5 & No inter-visit monitoring & Longitudinal tracking with PSI drift detection and CUSUM alerts & PSI $<$0.10 threshold \\[2pt]
P6 & High specialist cost & Consumer EEG + subscription model & Ch.~\ref{chap:discussion} \\[2pt]
P7 & Rural access barrier & FPGA edge processing, offline mode & $<$100\,ms latency \\[2pt]
P8 & Family anxiety & Caretaker dashboard with simplified health score & 5-screen app \\[2pt]
P9 & Doctor not notified & Auto HL7 FHIR alert to linked clinician & N3 notification \\[2pt]
P10 & Privacy concerns & AES-256, 7-stage PII pipeline, right to erasure & HIPAA/GDPR \\
\end{xltabular}}

\subsubsection{Mobile--Cloud Synchronization Protocol}
\label{app_sec:mobile_sync}

Continuous EEG monitoring in field settings (home, rural clinics, schools) requires robust offline-first synchronisation between the mobile gateway and cloud analytics. Table~\ref{app_tab:mobile_sync_stages} documents the five-stage RGAIG mobile--cloud synchronisation protocol, which ensures zero data loss over intermittent connectivity while maintaining HIPAA~\S164.312 integrity compliance and the NFR-7 availability target (99.5\% SLA). Sessions that fail all retry attempts are flagged for manual upload at the next clinical visit.
\needspace{8\baselineskip}
\noindent Table~\ref{app_tab:mobile_sync_stages} mobile--cloud synchronisation protocol: five sequential stages.

{\tablestripes\tablefontsize
\begin{xltabular}{\textwidth}{@{} p{0.5cm} p{2.2cm} L p{3.5cm} @{}}
\caption{Mobile--Cloud Synchronisation Protocol: Five Sequential Stages}\label{app_tab:mobile_sync_stages} \\
\toprule
\thead{\#} & \thead{Stage} & \thead{Operation} & \thead{Key Parameters} \\
\midrule
\endfirsthead
\multicolumn{4}{@{}l}{\textit{(\,Tab.~\ref{app_tab:mobile_sync_stages} continued from previous page\,)}} \\
\toprule
\thead{\#} & \thead{Stage} & \thead{Operation} & \thead{Key Parameters} \\
\midrule
\endhead
\bottomrule
\endlastfoot
1 & Local capture & Mobile app buffers EEG sessions locally with unique session identifiers, timestamps, and device attestation tokens & Session ID (UUID); device attestation token \\
\addlinespace
2 & Quality gate & Edge-level signal quality checks reject unusable recordings before queuing & Impedance $< 10$\,k$\Omega$; SNR $> 10$\,dB \\
\addlinespace
3 & Store-and-forward queue & Validated sessions enter a persistent FIFO queue with exponential backoff retry for upload over intermittent connectivity & Base delay 5\,s; max 3 attempts; exponential backoff \\
\addlinespace
4 & Cloud receipt & Ingestion layer returns a signed acknowledgement; mobile app marks session as synchronised only upon verified receipt & SHA-256 integrity hash; signed ACK \\
\addlinespace
5 & Conflict resolution & Duplicate uploads from multiple devices resolved via last-writer-wins policy with full version history retained for audit & Last-writer-wins; full version history \\
\end{xltabular}}
\vspace{0.5em}\noindent\textit{Note.} SNR = signal-to-noise ratio; FIFO = first in, first out; ACK = acknowledgement; UUID = universally unique identifier. This protocol supports the NFR-7 availability target (99.5\% SLA) and the HIPAA~\S164.312 integrity requirement.

\subsection{Model Selection and Output Evaluation}

Table~\ref{tab:model_selection} documents the domain-adaptive model selection strategy.
\needspace{8\baselineskip}
\noindent Table~\ref{app_tab:model_selection} domain-adaptive model selection criteria.

{\tablestripes\tablefontsize
\begin{xltabular}{\textwidth}{@{}l X l l@{}}
\caption{Domain-adaptive model selection criteria.}\label{app_tab:model_selection} \\
\toprule
\thead{Domain} & \thead{Selection Criteria} & \thead{Best Model} & \thead{Results} \\
\midrule
\endfirsthead
\multicolumn{4}{@{}l}{\textit{(\,Tab.~\ref{app_tab:model_selection} continued from previous page\,)}} \\
\toprule
\thead{Domain} & \thead{Selection Criteria} & \thead{Best Model} & \thead{Results} \\
\midrule
\endhead
\bottomrule
\endlastfoot
ASD & High-dimensional EEG, small subject count ($n$=19), connectivity features & Ensemble & Ch.~\ref{chap:analysis} \\[2pt]
PD & Motor cortex signals, tremor-related features, high separability & Ensemble & Ch.~\ref{chap:analysis} \\[2pt]
Stress & Multi-band EEG + physiological, arousal-related features & Ensemble & Ch.~\ref{chap:analysis} \\[2pt]
BCI & Motor imagery, spatial patterns, small sample ($n$=9) & Deep learning & Ch.~\ref{chap:analysis} \\[2pt]
Cancer & PBS microscopy images, visual pattern recognition, 2D spatial & Deep learning & Ch.~\ref{chap:analysis} \\
\end{xltabular}}

%% SUPPRESSED — redundant with tab:evaluation_metrics_summary (same metrics, fewer details)
\iffalse
\needspace{8\baselineskip}
\noindent Table~\ref{app_tab:output_evaluation} output evaluation: eight quality metrics for every ai prediction.

{\tablestripes\tablefontsize
\begin{xltabular}{\textwidth}{@{}l X l l@{}}
\caption{Output evaluation: eight quality metrics for every AI prediction.}\label{app_tab:output_evaluation} \\
\toprule
\thead{Metric} & \thead{What It Measures} & \thead{Threshold} & \thead{Results} \\
\midrule
\endfirsthead
\multicolumn{4}{@{}l}{\textit{(\,Tab.~\ref{app_tab:output_evaluation} continued from previous page\,)}} \\
\toprule
\thead{Metric} & \thead{What It Measures} & \thead{Threshold} & \thead{Results} \\
\midrule
\endhead
\bottomrule
\endlastfoot
Accuracy & Overall correct classification rate & $\geq$85\% & Ch.~\ref{chap:analysis} \\[2pt]
AUC-ROC & Discrimination across all thresholds & $\geq$0.90 & Ch.~\ref{chap:analysis} \\[2pt]
Effect Size & Practical significance (Cohen's $d$) & $\geq$0.80 & Ch.~\ref{chap:analysis} \\[2pt]
Faithfulness & RAG consistency with retrieved evidence & $\geq$0.85 & Ch.~\ref{chap:analysis} \\[2pt]
Hallucination & Unsupported claims in RAG output & $\leq$5\% & Ch.~\ref{chap:analysis} \\[2pt]
Citation Accuracy & Correct literature references & $\geq$90\% & Ch.~\ref{chap:analysis} \\[2pt]
Expert Rating & Clinician panel acceptance (Likert) & $\geq$4.0 & Ch.~\ref{chap:analysis} \\[2pt]
Inter-rater & Agreement (Krippendorff's $\alpha$) & $\geq$0.80 & Ch.~\ref{chap:analysis} \\
\end{xltabular}}
\fi

\subsection{Data Capture Types}

Table~\ref{tab:data_capture} documents the eight data sources across three modalities that feed the RGAIG pipeline.
\needspace{8\baselineskip}
\noindent Table~\ref{app_tab:data_capture} data capture types: eight sources across three modalities.

{\tablestripes\tablefontsize
\begin{xltabular}{\textwidth}{@{}l l X l l@{}}
\caption{Data capture types: eight sources across three modalities.}\label{app_tab:data_capture} \\
\toprule
\thead{Source} & \thead{Modality} & \thead{Data Captured} & \thead{Format} & \thead{Volume} \\
\midrule
\endfirsthead
\multicolumn{5}{@{}l}{\textit{(\,Tab.~\ref{app_tab:data_capture} continued from previous page\,)}} \\
\toprule
\thead{Source} & \thead{Modality} & \thead{Data Captured} & \thead{Format} & \thead{Volume} \\
\midrule
\endhead
\bottomrule
\endlastfoot
Emotiv EPOC X & EEG & 14-channel brainwave signals, 128\,Hz & EDF/CSV & 2\,MB/min \\[2pt]
Emotiv Insight & EEG & 5-channel portable EEG, 128\,Hz & EDF & 0.7\,MB/min \\[2pt]
Smartwatch & Physiol. & Heart rate variability, activity, sleep & JSON & 10\,KB/min \\[2pt]
ECG Strap & Cardiac & Cardiac rhythm, R-R intervals & CSV & 50\,KB/min \\[2pt]
Camera & Imaging & Skin lesion photos, PBS microscopy & JPEG/PNG & 5\,MB/image \\[2pt]
Metadata & Structured & Age, gender, medication, demographics & JSON/CSV & 1\,KB/record \\[2pt]
Public Datasets & Mixed & KAU, PhysioNet, DEAP, SAM40, BCI-IV, TCIA & EDF/NPZ & 58\,GB total \\[2pt]
Self-Report & Quest. & Symptom diary, mood, medication adherence & JSON & 2\,KB/entry \\
\end{xltabular}}


% [SPACE-FIX] clearpage removed to eliminate empty page
\subsection{Clinical Assessment Forms and Survey Instruments}
\label{app:sec9_forms}


In the B2B enterprise context, RGAIG integrates with existing hospital Electronic Health Record (EHR) systems via HL7 FHIR R4, enabling clinicians to order AI-assisted screening as part of routine clinical workflows. Table~\ref{app_tab:b2b_pathway} documents the eight-step B2B assessment pathway from patient enrolment to longitudinal trend monitoring.
\needspace{8\baselineskip}
{\tablestripes\tablefontsize
\begin{xltabular}{\textwidth}{@{} p{0.5cm} p{2.5cm} L p{2.5cm} @{}}
\caption{B2B Clinical Assessment Pathway: Eight Sequential Steps}\label{app_tab:b2b_pathway} \\
\toprule
\thead{\#} & \thead{Step} & \thead{Description} & \thead{System Component} \\
\midrule
\endfirsthead
\multicolumn{4}{@{}l}{\textit{(\,Tab.~\ref{app_tab:b2b_pathway} continued from previous page\,)}} \\
\toprule
\thead{\#} & \thead{Step} & \thead{Description} & \thead{System Component} \\
\midrule
\endhead
\bottomrule
\endlastfoot
1 & Patient enrolment & Clinician registers patient in RGAIG via EHR integration (HL7 FHIR R4) & EHR API gateway \\
\addlinespace
2 & EEG capture & Wearable or clinical EEG recorded; signal quality validated at point of capture & Emotiv SDK; quality gate \\
\addlinespace
3 & API submission & EEG data securely uploaded to cloud via encrypted REST endpoint & TLS 1.3; AES-256 \\
\addlinespace
4 & ASD-focused screening & Five-domain parallel classification (ASD) & Model orchestrator \\
\addlinespace
5 & Report generation & RAG-powered clinical narrative with SHAP attribution and literature citations & RAG pipeline \\
\addlinespace
6 & Report delivery & Results delivered to clinician dashboard and patient app simultaneously & FHIR R4; push notification \\
\addlinespace
7 & Clinical decision & Clinician reviews AI output, accepts/modifies recommendation, documents decision & Human-in-the-loop gate \\
\addlinespace
8 & Follow-up tracking & Longitudinal biomarker monitoring with PSI drift detection and CUSUM alerts & Monitoring framework \\
\end{xltabular}}
\vspace{0.5em}\noindent\textit{Note.} EHR = electronic health record; HL7 FHIR = Health Level Seven Fast Healthcare Interoperability Resources; PSI = population stability index; CUSUM = cumulative sum control chart. Steps 4--6 execute in under 30\,s for a single patient screening.

%===============================================================================
\subsection{Patient Assessment Questionnaire Framework}
\label{app_sec:assessment_questionnaires}

The RGAIG framework incorporates structured assessment questionnaires for each target condition to ensure clinical grounding. Tables~\ref{tab:asd_assessment}--\ref{tab:epilepsy_assessment} present condition-specific assessment instruments that inform the AI model's contextual reasoning.

%% ASD Assessment
\noindent The ASD assessment framework (Table~\ref{app_tab:asd_assessment}) integrates behavioural observation, standardised instruments (ADOS-2, ADI-R), and EEG-based biomarkers to provide a multi-modal evaluation that contextualises the AI model's classification output.
\needspace{8\baselineskip}
{\tablestripes\tablefontsize
\begin{xltabular}{\textwidth}{@{} >{\raggedright\arraybackslash}p{3.2cm} L @{}}
\caption{Autism Spectrum Disorder (ASD) Assessment Framework}\label{app_tab:asd_assessment} \\
\toprule
\thead{Assessment Type} & \thead{Questions / Measures} \\
\midrule
\endfirsthead
\multicolumn{2}{@{}l}{\textit{(\,Tab.~\ref{app_tab:asd_assessment} continued from previous page\,)}} \\
\toprule
\thead{Assessment Type} & \thead{Questions / Measures} \\
\midrule
\endhead
\bottomrule
\endlastfoot
Social Interaction & Does the patient avoid eye contact or social engagement? Difficulty understanding social cues? \\
\addlinespace
Communication & Is there delayed speech development? Does the patient repeat words/phrases (echolalia)? \\
\addlinespace
Behaviour Patterns & Repetitive behaviours (hand flapping, rocking)? Insistence on routines and resistance to change? \\
\addlinespace
Sensory Sensitivity & Sensitivity to sound, light, or touch? Over- or under-response to sensory stimuli? \\
\addlinespace
Cognitive & Ability to focus, attention span, learning ability, executive function assessment \\
\addlinespace
Emotional Regulation & Frequent meltdowns or emotional distress? Difficulty managing transitions? \\
\addlinespace
EEG/Neurological & Abnormal brainwave patterns in response to social and auditory stimuli; mu-rhythm suppression analysis \\
\addlinespace
AI-Based (RGAIG) & Pattern recognition in EEG signals for social/cognitive anomalies; SHAP feature attribution \\
\addlinespace
Parent/Guardian & Does the child respond to name? Show interest in peers? Joint attention behaviours? \\
\end{xltabular}}
\vspace{0.5em}\noindent\textit{Note.} ASD assessment combines behavioural observation, standardised instruments (ADOS-2, ADI-R), and EEG-based biomarkers. The RGAIG framework uses EEG pattern analysis with RAG-generated explanations referencing clinical literature. SHAP = SHapley Additive exPlanations.

%% Parkinson's Assessment
\noindent The PD assessment framework (Table~\ref{app_tab:pd_assessment}) combines motor evaluation items aligned with the MDS-UPDRS Part~III, non-motor screening (sleep, cognition, mood, autonomic), and EEG biomarkers (beta-band power, cortical slowing) to capture the full clinical picture used by the RGAIG classifier.
\needspace{8\baselineskip}
\noindent Table~\ref{app_tab:pd_assessment} parkinson's disease assessment framework.

{\tablestripes\tablefontsize
\begin{xltabular}{\textwidth}{@{} >{\raggedright\arraybackslash}p{3.2cm} L @{}}
\caption{Parkinson's Disease  Assessment Framework}\label{app_tab:pd_assessment} \\
\toprule
\thead{Assessment Type} & \thead{Questions / Measures} \\
\midrule
\endfirsthead
\multicolumn{2}{@{}l}{\textit{(\,Tab.~\ref{app_tab:pd_assessment} continued from previous page\,)}} \\
\toprule
\thead{Assessment Type} & \thead{Questions / Measures} \\
\midrule
\endhead
\bottomrule
\endlastfoot
Motor Symptoms & Resting tremor present? Bradykinesia (slowness of movement)? Rigidity in limbs? \\
\addlinespace
Gait \& Balance & Shuffling gait? Postural instability? Freezing episodes during walking? \\
\addlinespace
Non-Motor & Sleep disturbances (REM behaviour disorder)? Cognitive decline? Depression or anxiety? \\
\addlinespace
Speech & Hypophonia (soft speech)? Monotone voice? Difficulty with articulation? \\
\addlinespace
Daily Function & Difficulty with fine motor tasks (writing, buttoning)? Independence in daily activities? \\
\addlinespace
Cognitive & Executive function decline? Memory issues? Visuospatial difficulties? \\
\addlinespace
Autonomic & Orthostatic hypotension? Constipation? Urinary frequency? \\
\addlinespace
EEG/Neurological & Beta-band power changes; cortical slowing patterns; event-related desynchronisation \\
\addlinespace
AI-Based (RGAIG) & Multi-modal pattern recognition (EEG + motor + clinical); severity staging prediction \\
\addlinespace
Medication Response & Levodopa response assessment; on/off period tracking; dosage optimisation markers \\
\end{xltabular}}
\vspace{0.5em}\noindent\textit{Note.} PD assessment uses the MDS-UPDRS (Movement Disorder Society Unified Parkinson's Disease Rating Scale) as the gold standard. EEG biomarkers complement motor assessment for early detection. REM = Rapid Eye Movement; MDS-UPDRS = Movement Disorder Society Unified Parkinson's Disease Rating Scale.

%% Epilepsy Assessment
\noindent The epilepsy assessment framework (Table~\ref{app_tab:epilepsy_assessment}) captures seizure type, frequency, triggers, and EEG findings---including interictal epileptiform discharges---that the RGAIG automated spike-detection module uses to generate classification and alerting outputs.
\needspace{8\baselineskip}
{\tablestripes\tablefontsize
\begin{xltabular}{\textwidth}{@{} >{\raggedright\arraybackslash}p{3.2cm} L @{}}
\caption{Epilepsy Assessment Framework}\label{app_tab:epilepsy_assessment} \\
\toprule
\thead{Assessment Type} & \thead{Questions / Measures} \\
\midrule
\endfirsthead
\multicolumn{2}{@{}l}{\textit{(\,Tab.~\ref{app_tab:epilepsy_assessment} continued from previous page\,)}} \\
\toprule
\thead{Assessment Type} & \thead{Questions / Measures} \\
\midrule
\endhead
\bottomrule
\endlastfoot
Seizure Type & Focal or generalised? Tonic-clonic, absence, myoclonic, or atonic seizures? \\
\addlinespace
Seizure Frequency & How often do seizures occur? Duration of episodes? Clustering patterns? \\
\addlinespace
Triggers & Identified triggers (sleep deprivation, stress, flashing lights, alcohol)? \\
\addlinespace
Aura Symptoms & Warning signs before seizure (visual, auditory, olfactory, emotional)? \\
\addlinespace
Post-Ictal State & Confusion, fatigue, headache, or memory loss after seizure? Duration of recovery? \\
\addlinespace
Medication & Current anti-epileptic drugs? Effectiveness? Side effects? Compliance? \\
\addlinespace
Daily Impact & Impact on work, driving, relationships? Safety concerns? \\
\addlinespace
Cognitive & Memory problems? Concentration difficulties? Word-finding issues? \\
\addlinespace
EEG/Neurological & Interictal epileptiform discharges; spike-wave complexes; focal/generalised patterns \\
\addlinespace
AI-Based (RGAIG) & Automated spike detection in EEG; seizure prediction modelling; pattern classification \\
\addlinespace
Continuous Monitoring & Wearable EEG for ambulatory monitoring; real-time seizure detection and alerting \\
\end{xltabular}}
\vspace{0.5em}\noindent\textit{Note.} Epilepsy assessment relies heavily on EEG interpretation for diagnosis and classification. The RGAIG framework provides automated interictal spike detection with RAG-generated clinical explanations; classification results are reported in Chapter~\ref{chap:analysis}. Continuous wearable monitoring enables real-time seizure alerting outside clinical settings.

%===============================================================================
\subsection{Research Survey Instruments}
\label{app_sec:survey_instruments}

The RGAIG framework evaluation employs three validated survey instruments targeting distinct stakeholder groups: clinicians (B2B), consumers (B2C), and domain experts. Each instrument uses a 5-point Likert scale (1 = Strongly Disagree to 5 = Strongly Agree) unless otherwise specified. the corresponding table--N/A present the complete item sets with theoretical grounding.

Trust in AI diagnostics is measured through a six-item construct adapted from Lee and See~\cite{lee2004trust} and Madsen and Gregor~\cite{madsen2000measuring}, achieving internal consistency of $\alpha = 0.91$ and convergent validity AVE $= 0.69$ in pilot testing. Items assess perceived reliability, predictability, and faith in the system's diagnostic outputs. These six items (C5, C6, C11, C12, C13, C14 in the corresponding table) collectively operationalise the trust construct that mediates between governance antecedents and adoption intention in the structural model (SH1--SH7). High internal consistency confirms that clinicians interpret trust as a coherent latent variable rather than a set of disparate perceptions---a prerequisite for the mediation analysis reported in Chapter~\ref{chap:analysis}.

%% Clinician Survey (B2B)
\noindent Table~\ref{app_tab:clinician_survey} presents the 20-item clinician survey instrument, operationalising six constructs---perceived usefulness, ease of use, cognitive and affective trust, explainability, perceived risk, governance, human control, social influence, and adoption intention---grounded in TAM, UTAUT, and trust theory.
\needspace{8\baselineskip}
\noindent Table~\ref{app_tab:clinician_survey} clinician survey instrument (b2b): ai trust, adoption, and governance.

{\tablestripes\tablefontsize
\begin{xltabular}{\textwidth}{@{} >{\raggedright\arraybackslash}p{0.5cm} >{\raggedright\arraybackslash}p{2.5cm} L p{1.5cm} @{}}
\caption{Clinician Survey Instrument (B2B): AI Trust, Adoption, and Governance}\label{app_tab:clinician_survey} \\
\toprule
\thead{\#} & \thead{Construct} & \thead{Survey Item} & \thead{Theory} \\
\midrule
\endfirsthead
\multicolumn{4}{@{}l}{\textit{(\,Tab.~\ref{app_tab:clinician_survey} continued from previous page\,)}} \\
\toprule
\thead{\#} & \thead{Construct} & \thead{Survey Item} & \thead{Theory} \\
\midrule
\endhead
\bottomrule
\endlastfoot
C1 & Perceived Usefulness & The AI-assisted EEG interpretation system would improve the accuracy of my diagnostic decisions & TAM \\
\addlinespace
C2 & Perceived Usefulness & The system would reduce the time I spend on routine EEG analysis & TAM \\
\addlinespace
C3 & Ease of Use & I would find the AI system easy to integrate into my current clinical workflow & TAM \\
\addlinespace
C4 & Ease of Use & The system interface is intuitive and requires minimal training & TAM \\
\addlinespace
C5 & Cognitive Trust & I trust the AI system's outputs because they are based on validated clinical evidence & Trust Theory \\
\addlinespace
C6 & Affective Trust & I feel confident relying on the AI system for patient screening recommendations & Trust Theory \\
\addlinespace
C7 & Explainability & The AI explanations (SHAP values, RAG citations) help me understand why a recommendation was made & XAI \\
\addlinespace
C8 & Explainability & The literature citations in the AI report are relevant and clinically meaningful & XAI \\
\addlinespace
C9 & Perceived Risk & I am concerned about legal liability if the AI system produces an incorrect recommendation & Risk Theory \\
\addlinespace
C10 & Perceived Risk & I worry that patients might be harmed by false positive or false negative AI results & Risk Theory \\
\addlinespace
C11 & Governance & I believe the AI system has adequate safety controls and governance mechanisms & RGAIG \\
\addlinespace
C12 & Governance & The audit trail and decision logging give me confidence in system accountability & RGAIG \\
\addlinespace
C13 & Human Control & I feel I retain appropriate control over final diagnostic decisions when using the AI system & HITL \\
\addlinespace
C14 & Human Control & The system appropriately escalates uncertain cases to human review & HITL \\
\addlinespace
C15 & Social Influence & My colleagues' positive experiences with AI tools would influence my adoption decision & UTAUT \\
\addlinespace
C16 & Adoption Intent & I would recommend this AI system for use in my department & UTAUT \\
\addlinespace
C17 & Adoption Intent & I intend to use this AI system regularly if it becomes available in my hospital & UTAUT \\
\addlinespace
C18 & Cognitive Load & The AI system reduces my cognitive burden during complex multi-patient assessments & CLT \\
\addlinespace
C19 & Data Quality & I trust the quality of EEG data captured by the wearable device for clinical screening & DQ \\
\addlinespace
C20 & Overall Satisfaction & Overall, the AI-assisted EEG system meets my expectations for clinical decision support & Overall \\
\end{xltabular}}
\vspace{0.5em}\noindent\textit{Note.} 5-point Likert scale: 1 = Strongly Disagree, 2 = Disagree, 3 = Neutral, 4 = Agree, 5 = Strongly Agree. TAM = Technology Acceptance Model; UTAUT = Unified Theory of Acceptance and Use of Technology; XAI = Explainable AI; HITL = Human-in-the-Loop; CLT = Cognitive Load Theory; DQ = Data Quality; RGAIG = Responsible Generative AI Governance.

%% Consumer Survey (B2C)
\noindent The consumer survey (Table~\ref{app_tab:consumer_survey}) targets end-users of the wearable EEG system and measures health literacy, trust, explainability perception, decision support quality, privacy awareness, and behavioural safety across 20 items---constructs that operationalise the B2C adoption pathway tested in Chapter~\ref{chap:analysis}.
\needspace{8\baselineskip}
\noindent Table~\ref{app_tab:consumer_survey} consumer survey instrument (b2c): usability, trust, and.

{\tablestripes\tablefontsize
\begin{xltabular}{\textwidth}{@{} >{\raggedright\arraybackslash}p{0.5cm} >{\raggedright\arraybackslash}p{2.5cm} L p{1.5cm} @{}}
\caption[Consumer Survey Instrument (B2C): Usability, Trust, and]{Consumer Survey Instrument (B2C): Usability, Trust, and Decision Support}\label{app_tab:consumer_survey} \\
\toprule
\thead{\#} & \thead{Construct} & \thead{Survey Item} & \thead{Theory} \\
\midrule
\endfirsthead
\multicolumn{4}{@{}l}{\textit{(\,Tab.~\ref{app_tab:consumer_survey} continued from previous page\,)}} \\
\toprule
\thead{\#} & \thead{Construct} & \thead{Survey Item} & \thead{Theory} \\
\midrule
\endhead
\bottomrule
\endlastfoot
U1 & Understanding & I can understand the health insights the app provides about my EEG data & Literacy \\
\addlinespace
U2 & Understanding & The app explains medical terms in language I can comprehend & Literacy \\
\addlinespace
U3 & Trust & I trust the app to provide accurate information about my brain health & Trust \\
\addlinespace
U4 & Trust & I believe the app uses my EEG data responsibly and securely & Trust \\
\addlinespace
U5 & Explainability & The app clearly explains why it generated a particular health alert & XAI \\
\addlinespace
U6 & Explainability & I understand the confidence level shown with each AI recommendation & XAI \\
\addlinespace
U7 & Decision Support & The app helps me decide when I should consult a healthcare professional & Decision \\
\addlinespace
U8 & Decision Support & The app provides clear next steps after showing me my results & Decision \\
\addlinespace
U9 & Privacy & I am aware of how my EEG data is stored and who can access it & Privacy \\
\addlinespace
U10 & Privacy & The app gives me control over my data sharing preferences & Privacy \\
\addlinespace
U11 & Alert Quality & The alerts I receive are relevant and not excessive (no alert fatigue) & Usability \\
\addlinespace
U12 & Alert Quality & The severity levels of alerts help me prioritise which ones need attention & Usability \\
\addlinespace
U13 & Personalisation & The app provides insights tailored to my individual health patterns & Personal \\
\addlinespace
U14 & Behavioural Safety & The app does not cause me unnecessary anxiety or stress about my health & Safety \\
\addlinespace
U15 & Behavioural Safety & I feel calm and informed (not panicked) after reading the app's results & Safety \\
\addlinespace
U16 & Engagement & I use the app regularly because it provides valuable health insights & Engage \\
\addlinespace
U17 & Device Usability & The wearable EEG headset is comfortable and easy to use at home & Device \\
\addlinespace
U18 & Healthcare Link & If the app recommends seeing a doctor, I know how to connect with one through the app & Escalation \\
\addlinespace
U19 & Adoption Intent & I would recommend this wearable EEG system to family and friends & Adoption \\
\addlinespace
U20 & Overall Satisfaction & Overall, the wearable EEG app helps me better understand and manage my brain health & Overall \\
\end{xltabular}}
\vspace{0.5em}\noindent\textit{Note.} 5-point Likert scale: 1 = Strongly Disagree to 5 = Strongly Agree. Constructs mapped to: Health Literacy Theory, Trust Theory, XAI perception, Decision Support, Privacy Awareness, Usability (ISO 9241), Personalisation, Behavioural Safety Design, Engagement, and Technology Adoption.

%% Expert Validation Survey
\noindent The expert validation survey (Table~\ref{app_tab:expert_survey}) evaluates the RGAIG framework across 18 items spanning architecture, explainability, safety, trust, governance, accuracy, clinical utility, scalability, regulatory readiness, and innovation. A minimum panel of five evaluators from clinical neurology, AI/ML engineering, healthcare governance, and regulatory affairs provides the inter-rater reliability measured by Krippendorff's $\alpha$ in Chapter~\ref{chap:analysis}.
\needspace{8\baselineskip}
\noindent Table~\ref{app_tab:expert_survey} expert validation survey: rgaig framework evaluation.

{\tablestripes\tablefontsize
\begin{xltabular}{\textwidth}{@{} >{\raggedright\arraybackslash}p{0.5cm} >{\raggedright\arraybackslash}p{2.5cm} L p{1.5cm} @{}}
\caption{Expert Validation Survey: RGAIG Framework Evaluation}\label{app_tab:expert_survey} \\
\toprule
\thead{\#} & \thead{Dimension} & \thead{Survey Item} & \thead{RGAIG Layer} \\
\midrule
\endfirsthead
\multicolumn{4}{@{}l}{\textit{(\,Tab.~\ref{app_tab:expert_survey} continued from previous page\,)}} \\
\toprule
\thead{\#} & \thead{Dimension} & \thead{Survey Item} & \thead{RGAIG Layer} \\
\midrule
\endhead
\bottomrule
\endlastfoot
E1 & Architecture & The five-tier RGAIG architecture is well-structured for ASD AI & L1--L5 \\
\addlinespace
E2 & Architecture & The modular design allows effective domain extension to new conditions & L1--L5 \\
\addlinespace
E3 & Explainability & The RAG-based explanation generation provides clinically meaningful outputs & L3 \\
\addlinespace
E4 & Explainability & The SHAP/LIME feature attribution adds value to clinical decision-making & L3 \\
\addlinespace
E5 & Safety & The six-gate guardrail system provides adequate protection against harmful outputs & L4 \\
\addlinespace
E6 & Safety & The confidence threshold and human-fallback mechanisms are appropriately designed & L4 \\
\addlinespace
E7 & Trust & The framework's trust-building mechanisms (audit trail, transparency) are sufficient for clinical adoption & L4 \\
\addlinespace
E8 & Governance & The 525-module quality taxonomy provides comprehensive coverage of AI governance & L4 \\
\addlinespace
E9 & Governance & The responsible AI compliance scoring demonstrates adequate governance maturity & L4 \\
\addlinespace
E10 & Accuracy & The classification accuracy (ASD composite accuracy) is acceptable for clinical screening & L2 \\
\addlinespace
E11 & Accuracy & The ASD-focused validation approach provides adequate evidence of generalisability & L2 \\
\addlinespace
E12 & Clinical Utility & The framework would be useful in real hospital settings for EEG-based screening & Clinical \\
\addlinespace
E13 & Clinical Utility & The RAG-generated reports are suitable for clinician review and patient communication & Clinical \\
\addlinespace
E14 & Scalability & The architecture can scale to handle enterprise-level patient volumes & Technical \\
\addlinespace
E15 & Regulatory & The framework addresses key regulatory requirements (FDA SaMD, EU AI Act, HIPAA) & Compliance \\
\addlinespace
E16 & Innovation & The RGAIG framework represents a meaningful contribution to healthcare AI governance & Novelty \\
\addlinespace
E17 & Completeness & The framework adequately addresses the B2B (hospital) and B2C (consumer) use cases & Scope \\
\addlinespace
E18 & Limitations & The framework honestly acknowledges its limitations and boundary conditions & Rigour \\
\end{xltabular}}
\vspace{0.5em}\noindent\textit{Note.} 5-point Likert scale: 1 = Strongly Disagree to 5 = Strongly Agree. Expert panel: minimum 5 evaluators from clinical neurology, AI/ML engineering, healthcare governance, and regulatory affairs. Inter-rater reliability measured by Krippendorff's $\alpha$ (target $\geq$0.80). RGAIG layers: L1 = Data, L2 = Model, L3 = Explanation, L4 = Governance, L5 = Interface.

%===============================================================================
\subsection{Clinical Assessment Forms}
\label{app_sec:clinical_assessment_forms}

In addition to the research surveys, the RGAIG system uses structured clinical assessment forms for each target condition. These forms capture patient-specific data that contextualises the AI model's EEG analysis. the corresponding table--N/A present the complete assessment instruments.

%% ASD Clinical Assessment Form
\noindent The ASD clinical assessment form (Table~\ref{app_tab:asd_form}) captures 20 structured items across nine domains---demographics, social interaction, communication, behaviour, sensory, cognitive, emotional, EEG findings, and prior diagnosis---that contextualise the AI model's EEG classification by providing the clinical ground truth against which predictions are validated.
\needspace{8\baselineskip}
\noindent Table~\ref{app_tab:asd_form} asd clinical assessment form: structured data capture.

{\tablestripes\tablefontsize
\begin{xltabular}{\textwidth}{@{} >{\raggedright\arraybackslash}p{0.5cm} >{\raggedright\arraybackslash}p{2.0cm} L p{1.4cm} @{}}
\caption{ASD Clinical Assessment Form: Structured Data Capture}\label{app_tab:asd_form} \\
\toprule
\thead{\#} & \thead{Domain} & \thead{Assessment Item} & \thead{Response Type} \\
\midrule
\endfirsthead
\multicolumn{4}{@{}l}{\textit{(\,Tab.~\ref{app_tab:asd_form} continued from previous page\,)}} \\
\toprule
\thead{\#} & \thead{Domain} & \thead{Assessment Item} & \thead{Response Type} \\
\midrule
\endhead
\bottomrule
\endlastfoot
A1 & Demographics & Patient age, sex, handedness, developmental history & Numeric/Select \\
\addlinespace
A2 & Social & Eye contact frequency during interaction (none / occasional / frequent) & 3-point scale \\
\addlinespace
A3 & Social & Response to name when called (no response / delayed / immediate) & 3-point scale \\
\addlinespace
A4 & Social & Interest in peer interaction (avoidant / passive / active) & 3-point scale \\
\addlinespace
A5 & Communication & Speech development stage (non-verbal / single words / phrases / fluent) & 4-point scale \\
\addlinespace
A6 & Communication & Presence of echolalia (none / occasional / frequent) & 3-point scale \\
\addlinespace
A7 & Communication & Non-verbal communication use (gestures, pointing, showing) & Yes/No + notes \\
\addlinespace
A8 & Behaviour & Repetitive behaviours observed (list: hand flapping, rocking, spinning, lining up) & Checklist \\
\addlinespace
A9 & Behaviour & Insistence on sameness / routine disruption reaction & 5-point severity \\
\addlinespace
A10 & Sensory & Sensory sensitivities (auditory, visual, tactile, olfactory, gustatory) & Checklist \\
\addlinespace
A11 & Sensory & Sensory seeking behaviours observed & Yes/No + notes \\
\addlinespace
A12 & Cognitive & Estimated cognitive level (below / at / above age expectations) & 3-point scale \\
\addlinespace
A13 & Emotional & Emotional regulation difficulties (frequency and severity) & 5-point severity \\
\addlinespace
A14 & EEG Quality & EEG signal quality rating (poor / acceptable / good / excellent) & 4-point scale \\
\addlinespace
A15 & EEG Findings & Mu-rhythm suppression observed during social stimuli & Yes/No \\
\addlinespace
A16 & EEG Findings & Abnormal spectral power in frontal/temporal regions & Yes/No + band \\
\addlinespace
A17 & Prior Diagnosis & Previous ASD screening results (ADOS-2, ADI-R, M-CHAT scores) & Numeric \\
\addlinespace
A18 & Medications & Current medications and supplements & Free text \\
\addlinespace
A19 & Comorbidities & Co-occurring conditions (ADHD, anxiety, epilepsy, sleep disorders) & Checklist \\
\addlinespace
A20 & Clinician Notes & Overall clinical impression and recommendation & Free text \\
\end{xltabular}}
\vspace{0.5em}{\hyphenpenalty=50\exhyphenpenalty=50\noindent\textit{Note.} This form is completed by the clinician during or after the EEG session. Items A14--A16 are populated automatically by the RGAIG system. ADOS-2 = Autism Diagnostic Observation Schedule; ADI-R = Autism Diagnostic Interview-Revised; M-CHAT = Modified Checklist for Autism in Toddlers.\par}

%% PD Clinical Assessment Form
\noindent\textbf{Parkinson's Disease Clinical Assessment Form:....} Table~\ref{app_tab:pd_form} below presents the detailed analysis.

\needspace{8\baselineskip}
{\tablestripes\tablefontsize
\begin{xltabular}{\textwidth}{@{} >{\raggedright\arraybackslash}p{0.5cm} >{\raggedright\arraybackslash}p{2.0cm} L p{1.4cm} @{}}
\caption{Parkinson's Disease Clinical Assessment Form: Structured Data Capture}\label{app_tab:pd_form} \\
\toprule
\thead{\#} & \thead{Domain} & \thead{Assessment Item} & \thead{Response Type} \\
\midrule
\endfirsthead
\multicolumn{4}{@{}l}{\textit{(\,Tab.~\ref{app_tab:pd_form} continued from previous page\,)}} \\
\toprule
\thead{\#} & \thead{Domain} & \thead{Assessment Item} & \thead{Response Type} \\
\midrule
\endhead
\bottomrule
\endlastfoot
P1 & Demographics & Patient age, sex, age at symptom onset, family PD history & Numeric/Select \\
\addlinespace
P2 & Motor: Tremor & Resting tremor presence and affected limbs (L/R hand, L/R leg, jaw) & Checklist + side \\
\addlinespace
P3 & Motor: Bradykinesia & Slowness of movement severity (finger tapping, hand movements) & 5-point MDS \\
\addlinespace
P4 & Motor: Rigidity & Limb and axial rigidity on passive movement & 5-point MDS \\
\addlinespace
P5 & Motor: Gait & Gait pattern (normal / mildly impaired / shuffling / freezing / unable) & 5-point scale \\
\addlinespace
P6 & Motor: Balance & Postural stability (pull test result) & 5-point MDS \\
\addlinespace
P7 & Non-Motor: Sleep & Sleep quality, REM behaviour disorder presence, daytime sleepiness & 5-point + Yes/No \\
\addlinespace
P8 & Non-Motor: Cognitive & MoCA score or equivalent cognitive screening result & Numeric (0--30) \\
\addlinespace
P9 & Non-Motor: Mood & Depression (PHQ-9 score), anxiety (GAD-7 score), apathy & Numeric \\
\addlinespace
P10 & Non-Motor: Autonomic & Orthostatic hypotension, constipation, urinary symptoms & Checklist \\
\addlinespace
P11 & Speech & Voice volume, monotone quality, articulation clarity & 5-point scale \\
\addlinespace
P12 & Function & Hoehn and Yahr stage (1--5); Schwab and England ADL score & Stage + \% \\
\addlinespace
P13 & Medication & Current PD medications, dosages, timing, on/off fluctuations & Free text \\
\addlinespace
P14 & Medication & Levodopa equivalent daily dose (LEDD) calculation & Numeric (mg) \\
\addlinespace
P15 & EEG Quality & Signal quality rating; electrode impedance check & 4-point scale \\
\addlinespace
P16 & EEG Findings & Beta-band power changes (frontal/parietal regions) & Auto-populated \\
\addlinespace
P17 & EEG Findings & Cortical slowing pattern detected; event-related desynchronisation & Auto-populated \\
\addlinespace
P18 & Disease Duration & Years since diagnosis; years since first symptom & Numeric \\
\addlinespace
P19 & Comorbidities & Co-occurring conditions (dementia, depression, falls, osteoporosis) & Checklist \\
\addlinespace
P20 & Clinician Notes & MDS-UPDRS Part III total score; overall impression; follow-up plan & Numeric + text \\
\end{xltabular}}
\vspace{0.5em}{\hyphenpenalty=50\exhyphenpenalty=50\noindent\textit{Note.} Motor items P2--P6 align with MDS-UPDRS Part III. Items P15--P17 are auto-populated by RGAIG. MDS = Movement Disorder Society; UPDRS = Unified Parkinson's Disease Rating Scale; MoCA = Montreal Cognitive Assessment; PHQ-9 = Patient Health Questionnaire; GAD-7 = Generalised Anxiety Disorder scale; LEDD = Levodopa Equivalent Daily Dose; ADL = Activities of Daily Living.\par}

%% Epilepsy Clinical Assessment Form
\noindent\textbf{Epilepsy Clinical Assessment Form: Structured Data....} Table~\ref{app_tab:epilepsy_form} below presents the detailed analysis.

\needspace{8\baselineskip}
{\tablestripes\tablefontsize
\begin{xltabular}{\textwidth}{@{} >{\raggedright\arraybackslash}p{0.5cm} >{\raggedright\arraybackslash}p{2.0cm} L p{1.4cm} @{}}
\caption{Epilepsy Clinical Assessment Form: Structured Data Capture}\label{app_tab:epilepsy_form} \\
\toprule
\thead{\#} & \thead{Domain} & \thead{Assessment Item} & \thead{Response Type} \\
\midrule
\endfirsthead
\multicolumn{4}{@{}l}{\textit{(\,Tab.~\ref{app_tab:epilepsy_form} continued from previous page\,)}} \\
\toprule
\thead{\#} & \thead{Domain} & \thead{Assessment Item} & \thead{Response Type} \\
\midrule
\endhead
\bottomrule
\endlastfoot
EP1 & Demographics & Patient age, sex, age at first seizure, family epilepsy history & Numeric/Select \\
\addlinespace
EP2 & Seizure Type & Classification: focal aware / focal impaired / tonic-clonic / absence / myoclonic / atonic & Select \\
\addlinespace
EP3 & Frequency & Seizure frequency (daily / weekly / monthly / yearly / seizure-free) & Select + count \\
\addlinespace
EP4 & Duration & Typical seizure duration (seconds / minutes) & Numeric \\
\addlinespace
EP5 & Triggers & Known triggers (sleep deprivation, stress, photosensitivity, alcohol, illness, missed medication) & Checklist \\
\addlinespace
EP6 & Aura & Aura presence and type (visual, auditory, olfactory, gustatory, emotional, d\'ej\`a vu) & Checklist \\
\addlinespace
EP7 & Ictal & Seizure description: motor activity, consciousness level, automatisms, lateralisation & Free text \\
\addlinespace
EP8 & Post-Ictal & Post-ictal symptoms: confusion (duration), headache, fatigue, Todd's paralysis & Checklist + time \\
\addlinespace
EP9 & Status & History of status epilepticus (prolonged seizure $>$5 minutes) & Yes/No + count \\
\addlinespace
EP10 & Medication & Current anti-seizure medications, doses, levels, adherence & Free text \\
\addlinespace
EP11 & Medication & Drug-resistant epilepsy: failed $\geq$2 appropriate medications at adequate doses & Yes/No \\
\addlinespace
EP12 & Cognitive & Cognitive function: memory, attention, processing speed, language & 5-point scale \\
\addlinespace
EP13 & Psychosocial & Impact on driving, employment, education, relationships & 5-point severity \\
\addlinespace
EP14 & Safety & Seizure-related injuries: falls, burns, drowning risk, nocturnal seizure precautions & Checklist \\
\addlinespace
EP15 & EEG Quality & Signal quality; recording duration; sleep/wake state during recording & 4-point + time \\
\addlinespace
EP16 & EEG Findings & Interictal epileptiform discharges: spike/sharp waves, location, frequency & Auto-populated \\
\addlinespace
EP17 & EEG Findings & Background activity: normal / mildly slow / moderately slow / severely abnormal & Auto-populated \\
\addlinespace
EP18 & EEG Findings & Seizure captured during recording (ictal pattern documented) & Yes/No \\
\addlinespace
EP19 & Imaging & MRI findings: normal / lesional (specify) / pending & Select + text \\
\addlinespace
EP20 & Clinician Notes & Epilepsy syndrome classification; treatment plan; surgery candidacy; follow-up & Free text \\
\end{xltabular}}
\vspace{0.5em}{\hyphenpenalty=50\exhyphenpenalty=50\noindent\textit{Note.} Items EP16--EP18 are auto-populated by the RGAIG system's automated spike detection; classification accuracy is reported in Chapter~\ref{chap:analysis}. Seizure classification follows ILAE 2017 criteria. Drug-resistant epilepsy defined per Kwan et al.\ (2010). ILAE = International League Against Epilepsy.\par}

%% ============================================================================
%% VARIABLE CLASSIFICATION AND MEASUREMENT MODEL
%% ============================================================================

\subsection{Variable Classification and Operationalisation}
\label{app_sec:variable_classification}

A rigorous quantitative study requires that every variable be explicitly classified by its role in the hypothesised causal model. Table~\ref{tab:variable_types} defines the five variable types used in this study, after which Table~\ref{tab:variable_master} provides the complete variable inventory with codes and measurement sources.
\needspace{8\baselineskip}
\noindent Table~\ref{app_tab:variable_types} variable types used in this study.

{\tablestripes\tablefontsize
\begin{xltabular}{\textwidth}{@{} >{\raggedright\arraybackslash}p{3.8cm} p{4.5cm} L @{}}
\caption{Variable Types Used in This Study}\label{app_tab:variable_types} \\
\toprule
\thead{Type} & \thead{Definition} & \thead{Role in Study} \\
\midrule
\endfirsthead
\multicolumn{3}{@{}l}{\textit{(\,Tab.~\ref{app_tab:variable_types} continued from previous page\,)}} \\
\toprule
\thead{Type} & \thead{Definition} & \thead{Role in Study} \\
\midrule
\endhead
\bottomrule
\endlastfoot
Independent Variable (IV) & What is changed or tested & Drives outcomes (trust, adoption, decision confidence) \\
\addlinespace
Dependent Variable (DV) & What is measured as an outcome & Outcome of framework effectiveness \\
\addlinespace
Mediator Variable & Explains how IV affects DV & Intermediate mechanism through which IVs operate \\
\addlinespace
Control Variable (CV) & Kept constant to avoid bias & Eliminates confounding influences on DV \\
\addlinespace
Moderator Variable & Alters the strength of IV$\to$DV & Conditional effect that amplifies or attenuates the relationship \\
\end{xltabular}}
\vspace{0.5em}\noindent\textit{Note.} Variable classification follows Baron and Kenny's (1986) mediation framework and Aiken and West's (1991) moderation framework. Each variable is operationalised with a unique code for traceability across chapters.

